% !TeX program = xelatex

% QSG-Note-002-QSSC-Theorem
% Lang: RU
% Identifier: QSG-002
% Version: v6.0
% Date: 2026-01-18

\documentclass{article}
\usepackage[fontsize=12pt]{fontsize}
\usepackage{fontspec}
\usepackage{amsthm}
\usepackage{mathtools}
\usepackage{unicode-math}
\setmainfont{CMU Serif}
\setmathfont{Latin Modern Math}
\setmonofont{CMU Typewriter Text}[Scale=MatchLowercase]

\usepackage{polyglossia}
\setdefaultlanguage{russian}
\setotherlanguage{english}

\DeclarePairedDelimiter{\abs}{\lvert}{\rvert}
\DeclarePairedDelimiter{\norm}{\lVert}{\rVert}
\numberwithin{equation}{section}

\usepackage{microtype}
\usepackage[a4paper,margin=1in]{geometry}
\usepackage{booktabs}
\usepackage{float}
\usepackage{array}
\usepackage{enumitem}
\setlist{nosep,topsep=3pt,parsep=2pt}

\usepackage{graphicx}
\usepackage{tikz-cd}
\usetikzlibrary{arrows.meta}

\usepackage[square,numbers,sort&compress]{natbib}
\usepackage[hidelinks]{hyperref}
\usepackage{bookmark}
\usepackage[nameinlink,noabbrev]{cleveref}

\renewcommand{\proofname}{Доказательство}
\theoremstyle{plain}
\newtheorem{theorem}{Теорема}
\newtheorem{lemma}{Лемма}
\newtheorem{proposition}{Утверждение}
\newtheorem{corollary}{Следствие}
\newtheorem{assumption}{Предположение}[section]

\theoremstyle{definition}
\newtheorem{definition}{Определение}
\newtheorem{example}{Пример}[section]

\theoremstyle{remark}
\newtheorem*{remark}{Замечание}
\newtheorem{goal}{Цель}[section]

\newcommand{\imp}{\texorpdfstring{$\Rightarrow$}{=>}}

\providecommand{\dd}{\mathrm{d}}
\providecommand{\grad}{\nabla}
\providecommand{\laplace}{\Delta}
\providecommand{\Ric}{\mathrm{Ric}}
\providecommand{\Riem}{\mathrm{Riem}}
\providecommand{\Tr}{\mathrm{Tr}}
\providecommand{\diag}{\mathrm{diag}}

\renewcommand{\qedsymbol}{\ensuremath{□}}
\providecommand{\dalembert}{\Box}
\providecommand{\Boxm}[1]{\dalembert_{#1}}  % 

\providecommand{\Geff}{G_{\mathrm{eff}}}
\providecommand{\Leff}{\Lambda_{\mathrm{eff}}}
\providecommand{\gm}{g^{\mathrm{macro}}}
\providecommand{\Gmn}{G_{\mu\nu}(\gm)}
\providecommand{\Tmn}{T^{\mathrm{m}}_{\mu\nu}}
\providecommand{\Tphi}{T^{(\Phi)}_{\mu\nu}}

\providecommand{\Jphi}{\mathcal{J}_{\Phi}}

\providecommand{\pdv}[2]{\frac{\partial #1}{\partial #2}}

\providecommand{\Ph}{\Phi}
\providecommand{\sig}{\sigma}
\usepackage[most]{tcolorbox}
\usepackage{etex}

\hypersetup{
    pdfencoding=auto,
    colorlinks=true,
    linkcolor=blue,
    citecolor=blue,
    urlcolor=blue,
    pdftitle={Quantum Spectral Self-Consistency Theorem (QSSC-Theorem)},
    pdfauthor={Evgeny Monakhov},
    pdfsubject={Quantum Spectral Geometry, Operator Zeta Functions, Spectral Self-Consistency},
    pdfkeywords={
        spectral zeta function,
        operator zeta,
        self-adjoint operators,
        spectral measures,
        functional equation,
        reflection positivity,
        chiral balance,
        quantum spectral geometry,
        Hilbert--Pólya framework,
        Riemann hypothesis,
        QSSC,
        QSG-002
    }
}

\pdfstringdefDisableCommands{%
    \def\Rightarrow{⇒}%
    \def\Leftarrow{⇐}%
    \def\Leftrightarrow{⇔}%
    \def\rightarrow{→}%
    \def\leftarrow{←}%
    \def\leftrightarrow{↔}%
    \def\phi{φ}\def\Phi{Φ}\def\psi{ψ}\def\Psi{Ψ}%
    \def\theta{θ}\def\Theta{Θ}\def\pi{π}\def\Pi{Π}%
    \def\alpha{α}\def\beta{β}\def\gamma{γ}\def\Gamma{Γ}%
    \def\delta{δ}\def\Delta{Δ}\def\epsilon{ε}\def\zeta{ζ}%
    \def\eta{η}\def\mu{μ}\def\nu{ν}\def\lambda{λ}\def\Lambda{Λ}%
    \def\rho{ρ}\def\sigma{σ}\def\Sigma{Σ}\def\tau{τ}%
    \def\omega{ω}\def\Omega{Ω}\def\hbar{ℏ}\def\nabla{∇}%
    \def\infty{∞}\def\partial{∂}\def\sum{Σ}\def\int{∫}%
    \def\cdot{·}\def\times{×}\def\pm{±}\def\mp{∓}\def\approx{≈}%
    \def\propto{∝}\def\neq{≠}\def\leq{≤}\def\geq{≥}\def\sim{∼}\def\equiv{≡}%
}

\usepackage[ruled,vlined,linesnumbered]{algorithm2e}

\SetKwInput{KwIn}{Вход}
\SetKwInput{KwOut}{Выход}
\SetKwInput{KwData}{Дано}
\SetKwInput{KwResult}{Результат}
\SetKw{KwTo}{до}
\SetKw{KwRet}{вернуть}
\SetKw{KwAnd}{и}
\SetKw{KwOr}{или}
\SetKw{KwBreak}{прервать}
\SetKw{KwContinue}{продолжить}

\SetAlgoCaptionSeparator{.}
\SetAlCapNameFnt{\bfseries}
\SetAlCapFnt{\normalfont}
\SetAlgoNlRelativeSize{-1}
\SetAlgoInsideSkip{smallskip}
\DontPrintSemicolon
\SetNlSty{tiny}{}{:} 
\SetAlgoCaptionLayout{centerline}
\RestyleAlgo{ruled}
\SetKwComment{Comment}{$\triangleright$\ }{}

%------------ 1


\begin{document}
    \sloppy

\begin{titlepage}
    \centering

    {\Large\textbf{Quantum Spectral Geometry Notes}}\\[4pt]
    {\large\textbf{Volume 002}}\\[16pt]
    
    \vspace*{\fill}

    {\Large\textbf{
    Теорема о квантово-спектральной самосогласованности (QSSC-Theorem)
    }}\\[12pt]

    {\large\textit{
    Quantum Spectral Self-Consistency Theorem (QSSC-Theorem)
    }}\\[20pt]

    {\normalsize\textit{
    Операторная ζ-функция, функциональное уравнение и осевая локализация нулей\\
    в языке самосопряжённых операторов и спектральных мер
    }}\\[8pt]

    {\footnotesize\textit{
    Operator ζ-function, functional equation, and axis localization of zeros\\
    in the language of self-adjoint operators and spectral measures
    }}\\[40pt]

    {\large Евгений Монахов}\\[6pt]
    {\itshape Независимый исследователь}\\[2pt]
    \texttt{evgeny.monakhov@voscom.online}\\[2pt]
    \href{https://orcid.org/0009-0003-1773-5476}{ORCID: \texttt{0009-0003-1773-5476}}\\[30pt]

    {\footnotesize Identifier: QSG-002 v7.0}\\[6pt]

    \vspace*{\fill}

    {\large \today}

\end{titlepage}

\clearpage
\tableofcontents


\clearpage
%------------ 2

\begin{abstract}
\noindent
В работе формулируется и строго доказывается \emph{теорема о квантово-спектральной самосогласованности} (QSSC-Theorem), устанавливающая связь между самосопряжёнными операторами, их спектральными мерами и функциональным уравнением ζ-типа. 
В рамках аксиоматического каркаса (A1–A6) определяется операторная ζ-функция
\begin{equation}
\mathcal Z(u,s)=\langle\psi,\,g(H)\,(u^2+H^2)^{-s/2}\psi\rangle,\qquad u\ge0,
\label{eq:QSSC0001}
\end{equation}
где $H=H^\ast$ — самосопряжённый оператор в сепарабельном гильбертовом пространстве, $g$ — допустимая весовая функция, а $\psi$ — фиксированное состояние. 

Доказываются следующие результаты для абстрактного оператора:
\begin{enumerate}
  \item мероморфность $\mathcal Z(u,s)$ по $s\in\mathbb C$ при каждом фиксированном $u\ge0$;
  \item функциональная симметрия
\begin{equation}
  \widetilde{\mathcal Z}(u,s)
  =\pi^{-s/2}\Gamma\!\big(\tfrac{s}{2}\big)u^{s-1}\mathcal Z^{\mathrm{reg}}(u,s)
  =\widetilde{\mathcal Z}(u,1-s);
\label{eq:QSSC0002}
\end{equation}
  \item внутренняя эксплицитная формула (EF$^{\mathrm{qs}}$) как равенство спектральных мер;
  \item эквивалентность хирального баланса и осевой локализации нулей функции
  $\xi_H(s)=\lim_{u\downarrow0}\mathcal Z^{\mathrm{reg}}(u,s)$;
  \item достаточные условия (отражательная положительность ядра и антиунитарная симметрия) $\Rightarrow$ хиральный баланс $\Rightarrow$ осевая локализация.
\end{enumerate}

Для приложений теории к конкретным $L$-функциям вводится дополнительное условие \textbf{фазовой компенсации (аксиома A7)}. Подчеркивается, что данная поправка не является универсальной константой, а \textbf{зависит от самой функции}: для каждого типа $L$-функции оператор фазового сдвига определяется уникальной структурой её архимедовых множителей (гамма-факторов). В частности, для дзета-функции Римана аксиома A7 требует согласования с фазой Римана-Зигеля.

Тем самым QSSC-Theorem формализует статический закон спектральной самосогласованности: в абстрактном операторном каркасе осевая симметрия спектра возникает как внутреннее следствие унитарности, самодвойности и положительности.

Важно подчеркнуть, что в работе \emph{не утверждается} прямое доказательство гипотезы Римана. Однако показано, что \textbf{существование самосопряжённого оператора, удовлетворяющего общим аксиомам (A1–A6) и специфическому условию фазовой совместимости (A7), является достаточным условием для справедливости гипотезы Римана}. Таким образом, проблема сводится к вопросу существования оператора с заданной фазовой структурой в рамках предложенного формализма~\cite{Connes1999}.

%------------ 3
\textbf{Настоящая статья является второй частью серии работ по формированию
нового направления в математической физике, обозначаемого автором как
\emph{квантовая спектральная геометрия (Quantum Spectral Geometry, QSG)}.}~\cite{Monakhov-QSG-000}
Если первая часть цикла~\cite{Monakhov-QSG-001} была посвящена локальной спектральной плотности
(L-SED) и индуцированной конформной геометрии (ICG), задавая строгое
описание \emph{локальных} спектральных инвариантов и их геометрической
интерпретации, то настоящая работа строит \emph{глобальный ζ-уровень}
этой программы: вводится операторная ζ-функция, формулируется и доказывается
квантово-спектральная теорема самосогласованности (QSSC-Theorem), связывающая
мероморфность, функциональное уравнение и осевую локализацию нулей с
структурой самосопряжённого оператора и его спектральной меры.

В дальнейших работах цикла планируется объединить локальные (L-SED/ICG)
и глобальные (QSSC/ζ-уровень) конструкции в единую динамическую рамку
QSG: будут введены вариационное действие, уравнения спектральной динамики~\cite{ChamseddineConnes1997},
классификация допустимых спектральных режимов и связи с моделями
математической физики высоких энергий. Настоящая статья тем самым играет
роль «одномерного ζ-механического ядра» программы QSG, дополняя локальный
фундамент L-SED строго описанным ζ-типом самосогласованности спектра~\cite{ConnesMarcolli2008}.
\end{abstract}

\newpage

%------------ 4
\section {Расширенная аннотация (обзор результатов)}

\noindent
В статье формулируется и строго доказывается \textit{Теорема о квантово-спектральной самосогласованности} (QSSC-Theorem), 
устанавливающая внутреннюю связь между самосопряжёнными операторами, их спектральными мерами и функциональным уравнением ζ-типа.  
Рассматривается сепарабельное гильбертово пространство $\mathcal H$, самосопряжённый оператор $H=H^{\ast}$, фиксированное состояние $\psi\in\mathcal H$ и допустимая весовая функция $g$ из класса $\mathcal G$ (чётная, $C^\infty$, субэкспоненциального роста).  
Определяется операторная ζ-функция

\begin{equation}
\mathcal Z(u,s)
:=\big\langle \psi,\, g(H)\,(u^2+H^2)^{-s/2}\psi \big\rangle,
\qquad u>0.
\end{equation}

Значение при $u=0$ понимается в смысле регуляризованного предела
\begin{equation}
\mathcal Z^{\mathrm{reg}}(0,s)
:=\lim_{u\downarrow 0}\mathcal Z(u,s),
\end{equation}
при условии существования данного предела.


\noindent
для которой доказываются следующие утверждения:
%------------ 5
\begin{enumerate}
  \item[(i)] \textbf{Операторное определение.} 
  Для каждого фиксированного $u \ge 0$ функция $\mathcal Z(u,s)$ корректно определена как аналитическая функция по $s$ в некоторой правой полуплоскости $\{\Re s > \sigma_0\}$ и обладает интегральным представлением
  \begin{equation}
  \mathcal Z(u,s)
  =\int_{\mathbb R} (u^2+\lambda^2)^{-s/2}\,d\nu_g(\lambda),
  \end{equation}
  где $\nu_g$ — (в общем случае комплексная) спектральная мера, определяемая через спектральную теорему для $H$ и функциональное исчисление $g(H)$.


  \item[(ii)] \textbf{Мероморфность.} 
  Для каждого фиксированного $u\ge 0$ функция $\mathcal Z(u,s)$ допускает представление
  \begin{equation}
  \mathcal Z(u,s)
  =\frac{1}{\Gamma\!\big(\tfrac{s}{2}\big)}
   \int_0^\infty t^{\frac{s}{2}-1} e^{-u^2 t}\,K_g(t)\,dt,
  \end{equation}
  где
  \begin{equation}
  K_g(t):=\langle\psi,\,g(H)e^{-tH^2}\psi\rangle.
  \end{equation}
  Этот интеграл сходится и задаёт $\mathcal Z(u,s)$ для тех значений $s$, для которых обеспечена соответствующая асимптотика $K_g(t)$ при $t\downarrow 0$ и $t\to\infty$ (см. аксиомы (A4)–(A6)). 
  При этих предположениях функция $\mathcal Z(u,s)$ продолжается мероморфно по $s$ на всю комплексную плоскость $\mathbb C$, а полюса описываются асимптотикой $K_g(t)$.

  \item[(iii)] \textbf{Функциональное уравнение.}
  Из условий самодвойности ядра (чётности $g$, существования антиунитарной симметрии $\Theta$ и отражательной положительности в смысле~\cite{OsterwalderSchrader1973}) выводится функциональное уравнение
  \begin{equation}
  \widetilde{\mathcal Z}(u,s)
  =\pi^{-s/2}\Gamma\!\big(\tfrac{s}{2}\big)u^{s-1}\mathcal Z^{\mathrm{reg}}(u,s)
  =\widetilde{\mathcal Z}(u,1-s),
  \end{equation}
  т.\,е.\ симметрия $\widetilde{\mathcal Z}(u,s)=\widetilde{\mathcal Z}(u,1-s)$ является следствием операторно-спектральных предположений (A1–A6), а не отдельной аксиомой.


%------------ 6
\item[(iv)] \textbf{Эксплицитная формула (EF$^{\mathrm{qs}}$).}
Для спектральной меры $\nu_g$ имеет место каноническое разложение
\begin{equation}
\nu_g=\nu_{\mathrm{skel}}+\mu_\infty+\nu_{\mathrm{sing}},
\end{equation}
где $\nu_{\mathrm{skel}}$ — атомарная «скелетная» часть,
\begin{equation}
\mu_\infty(d\lambda)=a_g(\lambda)\,d\lambda,\qquad a_g\in L^1(\mathbb R),
\end{equation}
— абсолютно непрерывная архимедова компонента,
а $\nu_{\mathrm{sing}}$ — возможная сингулярно-непрерывная часть спектральной меры.
Тогда для любого тестового $h\in\mathcal S(\mathbb R)$ выполняется равенство
\begin{equation}
\int_{\mathbb R}\widehat h(\lambda)\,d\nu_g(\lambda)
=
\int_{\mathbb R}\widehat h(\lambda)\,d\nu_{\mathrm{skel}}(\lambda)
+
\int_{\mathbb R}\widehat h(\lambda)\,a_g(\lambda)\,d\lambda
+
\int_{\mathbb R}\widehat h(\lambda)\,d\nu_{\mathrm{sing}}(\lambda).
\end{equation}
В рассматриваемом операторном классе (см.\ аксиомы (A4)–(A6)) сингулярно-непрерывная компонента
$\nu_{\mathrm{sing}}$ либо отсутствует, либо не вносит вклада в спектральные функционалы,
индуцированные преобразованием Фурье на $\mathcal S(\mathbb R)$.
Полученное равенство тем самым представляет собой \emph{внутренний операторно-спектральный
аналог классической эксплицитной формулы}, в котором дискретный и непрерывный вклады
идентифицируются свойствами спектральной меры, а не внешними арифметическими суммами.


%------------ 7
  \item[(v)] \textbf{Хиральный баланс и осевая локализация.}
  Вводится хиральный функционал
  \begin{equation}
  \mathbb W_\tau[h]
  :=W_+[h_{\varepsilon,\tau}]-W_-[h_{\varepsilon,\tau}],
  \qquad 
  h_{\varepsilon,\tau}(t)=h_\varepsilon(t)\sinh(\tau t),
  \end{equation}
  где $W_\pm[\cdot]$ — спектральные функционалы, построенные из меры $\nu_g$. 
  Доказывается эквивалентность
  \begin{equation}
  \mathbb W_\tau[h]\equiv 0
  \quad \Longleftrightarrow \quad
  \text{все нули }\xi_H(s)
  =\lim_{u\downarrow0}\mathcal Z^{\mathrm{reg}}(u,s)
  \text{ лежат на оси }\Re s=\tfrac12.
  \end{equation}


  \item[(vi)] \textbf{Достаточные условия.}
  Если выполнены отражательная положительность, чётность $g$ и антиунитарная симметрия $\Theta$ ~\cite{OsterwalderSchrader1973}, то $\mathbb W_\tau[h]\equiv0$ для всех чётных $h$, а значит все нули $\xi_H$ лежат на критической оси.
\end{enumerate}

\noindent
Эти результаты образуют следующую логическую цепочку в рамках рассматриваемого операторного каркаса:

\begin{equation}
	\begin{aligned}
		&\text{(A1–A6)} 
		\;\Longrightarrow\;
		\text{мероморфность}
		\;\Longrightarrow\;
		\text{функциональная симметрия}
		\\
		&\Longrightarrow\;
		(\text{баланс}\;\Longleftrightarrow\;\text{ось}),
		\quad\text{и при RP: осевая локализация автоматически.}
	\end{aligned}
	\label{eq:QSSC0003}
\end{equation}

\medskip
\noindent
Для приложений теории к конкретным $L$-функциям вводится дополнительное условие \textbf{фазовой компенсации (аксиома A7)}. Подчеркивается, что данная поправка не является универсальной константой, а \textbf{зависит от самой функции}: для каждого типа $L$-функции оператор фазового сдвига определяется уникальной структурой её архимедовых множителей (гамма-факторов). В частности, для дзета-функции Римана аксиома A7 требует согласования с фазой Римана-Зигеля.
%------------ 8
\medskip
\noindent
\textbf{По сути.}  
Настоящая теорема формализует \emph{закон спектральной самосогласованности} в абстрактном операторном каркасе: 
в рамках этой модели осевая симметрия нулей ζ-функций ζ-типа может быть понята как следствие внутренних свойств самосопряжённых операторов и положительности, а не как отдельно постулируемое свойство.  
Таким образом, функциональное уравнение и осевая локализация спектра получают естественную операторно-спектральную интерпретацию; классические $\zeta$-функции (включая римановскую) в этом подходе рассматриваются как возможные частные реализации общего квантово-спектрального механизма.
\textbf{RH-следствие (условное).} 
При наличии самосопряжённого оператора, удовлетворяющего аксиомам (A1–A6) и \textbf{условию фазовой совместимости (A7)}, получаем осевую локализацию нулей соответствующей завершённой функции $\xi_H$, что в частности даёт \emph{условное} заключение о гипотезе Римана для стандартной дзета-функции при реализации её архимедовой фазы в рамках (A7)~\cite{Connes1999,PolyaHilbert,HaaKus2024}. Тем самым, ключевой вопрос переносится на задачу построения/существования оператора с требуемой фазовой структурой.

\clearpage
%------------ 9

\section{Введение}
\label{sec:intro}

\subsection{Почему необходим операторный подход к $\zeta$-функциям}

Классическая аналитическая теория $\zeta$-функции Римана, опирающаяся на ряды, интегральные представления и функциональное уравнение, чрезвычайно развита, однако сама структура нулей по-прежнему не выведена из более общих структурных «первых принципов».  
В частности, так называемая \emph{проблема оси} — доказательство того, что все нетривиальные нули функции $\zeta(s)$ лежат на прямой $\Re s=\tfrac12$ — остаётся в значительной мере неразрешённой в рамках существующих чисто аналитических методов: они позволяют получать оценочные результаты и частичную информацию о распределении нулей, но не дают строго сформулированного спектрального критерия, эквивалентного локализации нулей на оси.

Одна из причин этого состоит в следующем.  
Стандартные аналитические построения (ряд Дирихле, гамма-факторы, аналитическое продолжение) описывают поведение функции как аналитического объекта, но не фиксируют явно её возможный \emph{спектральный источник}~\cite{BerryKeating1999}.  
В таком языке информация о симметрии и осевой локализации нулей задаётся \emph{внешне} — через выбранное функциональное уравнение для «завершённой» $\zeta$-функции, — и не привязана напрямую к самосопряжённому оператору и его спектральной мере.

%------------ 10
Чтобы превратить проблему Римана в строго геометрическую задачу, необходимо ввести объект, который:
\begin{enumerate}
  \item[(a)] допускает самосопряжённый оператор $H=H^\ast$ с вещественным спектром $\sigma(H)\subset\mathbb R$;
  \item[(b)] реализует $\zeta$-функцию как коррелятор или интеграл по спектральной мере;
  \item[(c)] содержит внутреннюю симметрию, способную порождать функциональное уравнение без внешних предположений.
\end{enumerate}

Такой подход естественно возникает в рамках \emph{квантово-спектральной геометрии} — области, объединяющей методы спектральной теории операторов и квантовой механики для исследования аналитических объектов через их спектральные аналоги~\cite{Connes1999,ReedSimonI,ConnesMarcolli2008}.

%------------ 11
\medskip
\noindent
Пусть $H$ — самосопряжённый оператор в сепарабельном гильбертовом пространстве $\mathcal H$, и пусть
\begin{equation}
U(t)=e^{itH},\qquad t\in\mathbb R,
\end{equation}
— порождённая им унитарная группа.  
Если $\psi\in\mathcal H$ — фиксированное состояние, то коррелятор
\begin{equation}
\Xi(t):=\langle\psi,\,U(t)\psi\rangle
\end{equation}
представляет собой положительно-определённую функцию (в силу унитарности $U(t)$), 
а по теореме Бохнера~\cite{Bochner1933} ей соответствует конечная положительная спектральная мера $\nu$
(при $\psi\in\mathcal H$):
\begin{equation}
\Xi(t)=\int_{\mathbb R} e^{it\lambda}\,d\nu(\lambda),
\end{equation}
где $\nu$ — конечная положительная мера, являющаяся спектральной мерой, ассоциированной с парой $(H,\psi)$ (векторное состояние). 
Фурье-представление этой меры, а также связанные с ней интегральные трансформы ζ-типа, могут рассматриваться как \emph{спектральный аналог} классической $\zeta$-функции, в котором, как будет показано ниже, свойства оси и симметрии естественно выражаются как следствия самосопряжённости и отражательной положительности~\cite{LaxPhillips1976}.

%------------ 12
\medskip
\noindent
Таким образом, операторный подход переформулирует задачу «найти нули аналитической функции» в задачу «описать геометрию спектра самосопряжённого оператора, порождающего данный коррелятор».  
При этом понятия \emph{унитарности}, \emph{симметрии}, \emph{самосопряжённости} и \emph{отражательной положительности} выступают не в виде внешне навязанных аксиом, а как естественные условия внутренней согласованности выбранного квантово-спектрального каркаса.

\begin{quote}
\textit{Иными словами, если мы хотим рассматривать осевую симметрию нулей $\zeta$-функции как следствие, а не как постулат, естественно переформулировать задачу в спектральном пространстве, где ось выступает инвариантом антиунитарной симметрии.}
\end{quote}

\medskip
\noindent
Этот переход от аналитической к операторно-спектральной форме и составляет основной методологический мотив настоящей работы.


%------------ 13
\subsection{Наш вклад: единая операторная конструкция, функциональная симметрия и эквивалентность «баланса» и осевой локализации}
\label{subsec:our_contribution}

\noindent
Основной результат настоящей работы заключается в построении \emph{единой квантово-спектральной конструкции ζ-типа},
которая формализует в операторном каркасе ключевые аналитические свойства \emph{завершённых} ζ-функций —
мероморфность и функциональную симметрию — и связывает осевую локализацию нулей с внутренними спектральными
условиями (хиральным балансом и отражательной положительностью), не опираясь на \emph{явные} арифметические постулаты.

\medskip
\noindent
Пусть $(\mathcal H,H,\psi,g)$ — четверка, удовлетворяющая аксиомам (A1)–(A6), определённым в разделе~\ref{sec:axioms}.
Задаётся операторная ζ-функция
\begin{equation}
\mathcal Z(u,s)
:=\langle\psi,\,g(H)\,(u^2+H^2)^{-s/2}\psi\rangle,
\qquad u>0,
\end{equation}
которая первоначально определена как аналитическая функция по $s$ в подходящей правой полуплоскости
(в зависимости от роста спектра $H$ и веса $g$), а затем продолжается мероморфно на всю комплексную плоскость
(см.~раздел~\ref{sec:axioms} и~\cite{Seeley1967}).
Данный объект представляет собой квантово-спектральный аналог классической $\zeta$-функции,
определяемый не через ряды, а посредством функционального исчисления самосопряжённого оператора,
так что его аналитические свойства выводятся из \emph{внутренней структуры} тройки $(H,g,\psi)$.
подхода состоит в том, что свойства $\mathcal Z(u,s)$ выводятся \emph{из внутренней структуры} $(H,g,\psi)$, а не из внешних аналитических аппроксимаций.

%------------ 14
\paragraph{(i) Мероморфность и функциональное уравнение.}
Функция $\mathcal Z(u,s)$ имеет интегральное тепловое представление
\begin{equation}
\mathcal Z(u,s)
=\frac{1}{\Gamma(\tfrac s2)}
\int_0^\infty t^{\frac s2-1}e^{-u^2 t}\,\langle\psi,\,g(H)e^{-tH^2}\psi\rangle\,dt,
\end{equation}
что при выполнении аксиом теплового типа (см. (A4)–(A6)) обеспечивает мероморфное продолжение $\mathcal Z(u,s)$ по $s$ и позволяет строго определить регуляризованное продолжение $\mathcal Z^{\mathrm{reg}}(u,s)$ на всю комплексную плоскость~\cite{Seeley1967}.
При выполнении условий самодвойности (sd1–sd3) — чётности $g$, антиунитарной симметрии $\Theta$ и отражательной положительности — эта функция удовлетворяет \emph{функциональному уравнению}
\begin{equation}
\widetilde{\mathcal Z}(u,s)
=\pi^{-s/2}\Gamma\!\big(\tfrac s2\big)u^{s-1}\mathcal Z^{\mathrm{reg}}(u,s)
=\widetilde{\mathcal Z}(u,1-s),
\end{equation}
которое возникает из внутренней геометрии оператора и не требует внешне постулируемых симметрий.

%------------ 15
\paragraph{(ii) Внутренняя эксплицитная формула.}
Для спектральной меры $\nu_g$ (в общем случае комплексной), ассоциированной с коррелятором
\begin{equation}
\Xi_g(t)=\langle\psi,\,g(H)e^{itH}\psi\rangle
\end{equation}
через спектральную теорему для $H$, доказывается разложение
\begin{equation}
\nu_g=\nu_{\mathrm{skel}}+\mu_\infty+\nu_{\mathrm{sing}},
\end{equation}
где $\nu_{\mathrm{skel}}$ — дискретная (атомарная) часть,
\begin{equation}
\mu_\infty(d\lambda)=a_g(\lambda)\,d\lambda,\qquad a_g\in L^1(\mathbb R),
\end{equation}
— абсолютно непрерывная архимедова компонента,
а $\nu_{\mathrm{sing}}$ — возможная сингулярно-непрерывная часть спектральной меры.
Для любого тестового $h\in\mathcal S(\mathbb R)$ имеет место равенство
\begin{equation}
\int_{\mathbb R}\widehat h(\lambda)\,d\nu_g(\lambda)
=
\int_{\mathbb R}\widehat h(\lambda)\,d\nu_{\mathrm{skel}}(\lambda)
+
\int_{\mathbb R}\widehat h(\lambda)\,a_g(\lambda)\,d\lambda
+
\int_{\mathbb R}\widehat h(\lambda)\,d\nu_{\mathrm{sing}}(\lambda).
\end{equation}
В рассматриваемом операторном классе (см.\ аксиомы (A4)–(A6)) сингулярно-непрерывная компонента
$\nu_{\mathrm{sing}}$ либо отсутствует, либо не вносит вклада в \emph{в рассматриваемом классе спектральных функционалов},,
индуцированных  преобразованиями Фурье на $\mathcal S(\mathbb R)$.
Тем самым полученное равенство представляет собой \emph{операторно-спектральный аналог
классической эксплицитной формулы}, в котором дискретный и непрерывный вклады
идентифицируются внутренними свойствами спектральной меры,
без обращения к числовым рядам или суммам по простым числам~\cite{Schwartz1950}.


\paragraph{(iii) Хиральный баланс и эквивалентность с осевой локализацией.}
Вводится хиральный функционал
\begin{equation}
\mathbb W_\tau[h]
:=W_+[h_{\varepsilon,\tau}]-W_-[h_{\varepsilon,\tau}],
\qquad 
h_{\varepsilon,\tau}(t)=h_\varepsilon(t)\sinh(\tau t),
\end{equation}
описывающий асимметрию между «положительной» и «отрицательной» ветвями спектра.
В рассматриваемом операторном каркасе доказывается строгая эквивалентность
\begin{equation}
\mathbb W_\tau[h]\equiv 0
\ \text{для всех чётных } h\in\mathcal S(\mathbb R)
\text{ и достаточно малых } \tau
\quad \Longleftrightarrow \quad
\text{все нули }
\xi_H(s)
=\lim_{u\downarrow 0}\mathcal Z^{\mathrm{reg}}(u,s)
\text{ лежат на оси }\Re s=\tfrac12.
\end{equation}
Тем самым, положение нулей $\xi_H$ становится \emph{спектральным следствием условия хирального баланса},
а не отдельным аналитическим постулатом.

\paragraph{(iv) Достаточные условия.}
Если выполнены антиунитарная симметрия $\Theta$, чётность весовой функции $g$
и отражательная положительность ядра $\Xi_g$, то $\mathbb W_\tau[h]\equiv 0$
для всех чётных $h\in\mathcal S(\mathbb R)$ и достаточно малых $\tau$.
Из этого следует осевая локализация нулей $\xi_H$, что замыкает логическую цепочку:
\begin{equation}
\text{RP}+\Theta+\text{чётность }g
\;\Rightarrow\;
\text{хиральный баланс}
\;\Leftrightarrow\;
\text{осевая локализация}.
\label{eq:QSSC0004}
\end{equation}

\paragraph{(v) Концептуальный итог.}
Предложенная конструкция объединяет ключевые свойства ζ-функций
в \emph{внутреннюю квантово-спектральную структуру}:
\begin{itemize}
  \item мероморфность — следствие интегрального теплового представления;
  \item функциональное уравнение — следствие антиунитарной самодвойности ядра;
  \item осевая локализация нулей — следствие хирального баланса и отражательной положительности.
\end{itemize}
Таким образом, QSSC-Theorem формулирует строгий операторный критерий спектральной самосогласованности ζ-типа функций,
основанный на общих принципах спектральной геометрии~\cite{Connes1999,HelfferRobert1984,BaezSegal2000},
и позволяет рассматривать мероморфность, функциональное уравнение и осевую локализацию нулей
как взаимосвязанные следствия единой операторной структуры.

\subsection{План работы и структура доказательства}
\label{subsec:structure}

\noindent
Настоящая работа имеет строго логическую структуру, в которой каждая последующая глава
формально опирается на аксиомы и результаты предыдущих.
Основная цель статьи — \textit{построить и доказать квантово-спектральную теорему самосогласованности (QSSC-Theorem)},
формализующую закон \emph{внутренней спектральной самосогласованности}
для операторных ζ-функций.

\paragraph{Аксоматический каркас.}
Формулируются шесть базовых аксиом (A1)–(A6), определяющих общий спектральный каркас,
а также дополнительная аксиома (A7), необходимая для точной настройки фазы
в арифметических приложениях.
\begin{itemize}
  \item[(A1)] сепарабельное гильбертово пространство $\mathcal H$;
  \item[(A2)] самосопряжённый оператор $H=H^\ast$ и унитарная группа $U(t)=e^{itH}$;
  \item[(A3)] фиксированное состояние $\psi\in\mathcal H$;
  \item[(A4)] допустимый класс весовых функций $g\in\mathcal G$
  (чётные, $C^\infty$, субэкспоненциального роста);
  \item[(A5)] удвоение $\widetilde{\mathcal H}=\mathcal H\oplus\mathcal H$,
  $\widetilde H=\mathrm{diag}(H,-H)$ и антиунитарная симметрия $\Theta$;
  \item[(A6)] отражательная положительность, индуцированная коррелятором $\Xi_g(t)$
  в смысле Остервальдера–Шрадера;
  \item[(A7)] (опционально) фазовая компенсация $\Phi$ для согласования
  с архимедовой структурой конкретных $L$-функций.
\end{itemize}
Эти аксиомы задают минимальный аналитический и геометрический контекст,
необходимый для доказательства основной теоремы.

\paragraph{Коррелятор и спектральная мера.}
Определяется коррелятор
\begin{equation}
\Xi_g(t):=\langle\psi,\,g(H)e^{itH}\psi\rangle,
\end{equation}
и через спектральную теорему для самосопряжённого оператора $H$
вводится (в общем случае комплексная) внутренняя спектральная мера $\nu_g$:
\begin{equation}
\Xi_g(t)=\int_{\mathbb R} e^{it\lambda}\,d\nu_g(\lambda).
\end{equation}
Доказываются её основные свойства: чётность (в силу чётности $g$ и самосопряжённости $H$),
контроль усечений и корректность обменов интегралов
для класса весов $\mathcal G$.
При дополнительных предположениях на знак $g$ (например, $g\ge 0$ на спектре $H$)
коррелятор $\Xi_g$ становится положительно определённой функцией,
а мера $\nu_g$ — положительной.

\paragraph{Квантово-спектральная ζ-функция.}
Вводится операторная ζ-функция
\begin{equation}
\mathcal Z(u,s):=\langle\psi,\,g(H)\,(u^2+H^2)^{-s/2}\psi\rangle,
\qquad u>0,
\end{equation}
которая первоначально определена как аналитическая функция по $s$
в подходящей правой полуплоскости
(в зависимости от асимптотики спектра $H$ и веса $g$),
а затем продолжается мероморфно на всю комплексную плоскость.
Показано, что $\mathcal Z(u,s)$ имеет тепловое (меллиново) представление
через ядро~\cite{BerlineGetzlerVergne1992}
\begin{equation}
K_g(t)=\langle\psi,\,g(H)e^{-tH^2}\psi\rangle,
\label{eq:QSSC0005}
\end{equation}
и что при выполнении аксиом теплового типа (см.\ (A4)–(A6))
в пределе $u\downarrow 0$ корректно определяется функция
\begin{equation}
\xi_H(s):=\lim_{u\downarrow0}\mathcal Z^{\mathrm{reg}}(u,s),
\label{eq:QSSC0006}
\end{equation}
играющая роль «завершённой» ζ-функции, ассоциированной с оператором $H$.

\paragraph{Самодвойность и функциональное уравнение.}
Вводятся структурные условия (sd1–sd3): чётность $g$, антиунитарная симметрия $\Theta$
и отражательная положительность.
При выполнении этих условий в рассматриваемом операторном каркасе выводится
функциональное уравнение
\begin{equation}
\widetilde{\mathcal Z}(u,s)
=\pi^{-s/2}\Gamma\!\big(\tfrac s2\big)u^{s-1}\mathcal Z^{\mathrm{reg}}(u,s)
=\widetilde{\mathcal Z}(u,1-s),
\end{equation}
то есть симметрия $s\mapsto 1-s$ возникает как следствие самодвойности ядра
и не требует апелляции к классической аналитической теории ζ-функций
\emph{в качестве исходного постулата}.

\paragraph{Внутренняя эксплицитная формула (EF$^{\mathrm{qs}}$).}
Для спектральной меры $\nu_g$ в рассматриваемом операторном классе
(см.~аксиомы (A4)–(A6)) доказывается разложение
\begin{equation}
\nu_g=\nu_{\mathrm{skel}}+\mu_\infty,
\label{eq:QSSC0007}
\end{equation}
где $\nu_{\mathrm{skel}}$ — дискретная (атомарная) часть,
а $\mu_\infty$ абсолютно непрерывна и имеет плотность $a_g\in L^1(\mathbb R)$.
В терминах функционалов на пространстве Шварца $\mathcal S(\mathbb R)$
это разложение эквивалентно равенству
\begin{equation}
\int_{\mathbb R}\widehat h(\lambda)\,d\nu_g(\lambda)
=
\int_{\mathbb R}\widehat h(\lambda)\,d\nu_{\mathrm{skel}}(\lambda)
+
\int_{\mathbb R}\widehat h(\lambda)\,a_g(\lambda)\,d\lambda,
\label{eq:QSSC0008}
\end{equation}
для всех $h\in\mathcal S(\mathbb R)$.
Тем самым формулируется строгий аналог эксплицитной формулы
в чисто операторной форме.

\paragraph{Хиральный баланс и осевая локализация.}
Вводится хиральный функционал
\begin{equation}
\mathbb W_\tau[h]
:=W_+[h_{\varepsilon,\tau}]-W_-[h_{\varepsilon,\tau}],
\qquad h_{\varepsilon,\tau}(t)=h_\varepsilon(t)\sinh(\tau t),
\label{eq:QSSC0009}
\end{equation}
описывающий асимметрию между «положительной» и «отрицательной» ветвями спектра.
В рамках рассматриваемого операторного каркаса доказывается эквивалентность
\begin{equation}
\mathbb W_\tau[h]\equiv 0
\ \text{для всех чётных } h\in\mathcal S(\mathbb R)
\text{ и достаточно малых } \tau
\quad\Longleftrightarrow\quad
\text{все нули }\xi_H(s)\text{ лежат на оси }\Re s=\tfrac12.
\label{eq:QSSC0010}
\end{equation}
Эта теорема формулирует строгий операторный критерий осевой локализации
нулей $\xi_H$.


\paragraph{Достаточные условия (RP $\Rightarrow$ баланс).}
Показывается, что при выполнении условий отражательной положительности
(аксиомы Остервальдера–Шрадера), чётности весовой функции $g$
и наличии антиунитарной симметрии $\Theta$ выполняется равенство
$\mathbb W_\tau[h]\equiv0$ для всех чётных $h\in\mathcal S(\mathbb R)$
и достаточно малых $\tau$.
Это автоматически приводит к осевой локализации нулей $\xi_H$,
то есть реализует импликацию
\[
\text{RP}+\Theta+\text{чётность }g \;\Rightarrow\; \text{баланс} \;\Rightarrow\; \text{ось}.
\]

\paragraph{Технические оценки, дискретизации и специализации.}
Формулируются оценки сходимости интегралов и правила обмена пределов
(Фубини, Тонелли, теорема о доминированной сходимости)~\cite{Folland1999},
а также корректные схемы дискретизации
(в частности, логарифмическая $\varphi$-сетка)
для численной аппроксимации спектральной меры $\nu_g$.
Отдельный раздел посвящён возможным специализациям общей конструкции,
включая модельные примеры и обсуждение того,
каким образом классическая ζ-функция может быть реализована
в данном операторном каркасе.

\paragraph{Заключение и следствия.}
Подводятся итоги: при выполнении аксиом (A1)–(A6)
из них логически следуют мероморфность $\mathcal Z(u,s)$,
функциональная симметрия
$\widetilde{\mathcal Z}(u,s)=\widetilde{\mathcal Z}(u,1-s)$
и эквивалентность
«хиральный баланс $\Leftrightarrow$ осевая локализация нулей $\xi_H$».
Дополнительно, при наличии отражательной положительности (RP)
осевая локализация нулей $\xi_H$ возникает автоматически.
Отмечается, что все эти результаты выведены
\emph{без обращения к арифметическим постулатам}.

\bigskip
\noindent
\textbf{Интуитивное резюме.}
Данная работа переводит ключевые вопросы о ζ-типа функциях
из области \emph{классической аналитической теории ζ-функций}
в область \emph{спектральной геометрии}.
В рамках предлагаемого операторного каркаса
осевая симметрия нулей $\xi_H(s)$ перестаёт быть внешним постулатом
и становится внутренним следствием симметрии самосопряжённого оператора
и отражательной положительности его спектрального ядра.
Тем самым, квантово-спектральный подход формулирует общий
геометрический закон самосогласованности для ζ-типа функций;
в этом языке классическая гипотеза Римана может интерпретироваться
как утверждение о существовании подходящего самосопряжённого оператора $H$,
для которого ассоциированная функция $\xi_H$ совпадает
(с точностью до стандартных архимедовых нормировок)
с римановской ζ-функцией и удовлетворяет осевой локализации нулей.

\clearpage
%------------ 26

\section{Аксиоматический каркас}
\label{sec:axioms}


\noindent
Для строгого построения квантово-спектральной ζ-функции и доказательства теоремы самосогласованности необходимо задать минимальный, но достаточно богатый набор аксиом, обеспечивающих корректность всех последующих операций — функционального исчисления, интегральных преобразований, обменов пределов и симметрий~\cite{ReedSimon1}.

В данном разделе формулируются аксиомы (A1)–(A6), задающие общий \emph{операторно-спектральный каркас} для всей теории.  
По сути, это стандартные предпосылки функционального анализа и спектральной теории самосопряжённых операторов, адаптированные под структуру будущей ζ-конструкции.

%------------ 27
\subsection{(A1) Пространство: сепарабельное гильбертово пространство}

\noindent
Пусть $(\mathcal H,\langle\cdot,\cdot\rangle)$ — комплексное сепарабельное гильбертово пространство.
Скалярное произведение определяется в конвенции, линейной по \emph{второму} аргументу
и антилинейной по первому; далее в работе используется именно эта конвенция:
\begin{equation}
\langle x,y\rangle = \overline{\langle y,x\rangle}, 
\qquad 
\langle x,ay_1+by_2\rangle = a\,\langle x,y_1\rangle + b\,\langle x,y_2\rangle,
\end{equation}
для всех $x,y,y_1,y_2\in\mathcal H$ и $a,b\in\mathbb C$.
Сепарабельность обеспечивает существование счётного ортонормированного базиса
$\{e_n\}_{n\ge1}$ и позволяет рассматривать стандартные классы
векторнозначных интегралов, возникающих в спектральной теории,
в частности интегралы в смысле Бохнера при выполнении соответствующих
условий измеримости и интегрируемости~\cite{HillePhillips1957}.

\emph{Комментарий.}
Сепарабельность и полнота являются минимальными структурными предпосылками,
необходимыми для применения спектральной теоремы к самосопряжённым операторам
и для интерпретации $\mathcal H$ как пространства состояний в квантовом контексте.

\subsection{(A2) Оператор: самосопряжённость и унитарная группа}

\noindent
На $\mathcal H$ задан (возможно неограниченный) плотно определённый линейный оператор
$H:\mathcal D(H)\to\mathcal H$, который является самосопряжённым:
\begin{equation}
H=H^\ast .
\end{equation}
Эквивалентно, $H$ симметричен и его область определения совпадает с областью определения
сопряжённого оператора, то есть
\begin{equation}
\langle Hx,y\rangle=\langle x,Hy\rangle,
\qquad \forall\,x,y\in\mathcal D(H),
\end{equation}
и $\mathcal D(H)=\mathcal D(H^\ast)$.

По теореме Стоуна самосопряжённость $H$ эквивалентна существованию единственной
\emph{сильно непрерывной унитарной группы} $(U(t))_{t\in\mathbb R}$ на $\mathcal H$,
для которой $H$ является генератором:
\begin{equation}
U(t)U(s)=U(t+s),\qquad U(0)=I,\qquad
\text{и}\quad t\mapsto U(t)x \text{ непрерывна } \forall\,x\in\mathcal H .
\end{equation}
При этом $U(t)$ задаётся спектральным функциональным исчислением и обозначается
\begin{equation}
U(t)=e^{itH},\qquad t\in\mathbb R.
\end{equation}
Оператор $H$ далее называется \emph{гамильтонианом}, а $U(t)$ — унитарной динамикой.

\emph{Комментарий.}
Самосопряжённость гарантирует, что спектр $\sigma(H)\subset\mathbb R$ и к $H$ применима спектральная теорема.
Это позволяет определять через борелевское функциональное исчисление операторы вида $f(H)$.
В частности, для $u>0$ и тех значений $s$, при которых функция
$\lambda\mapsto (u^2+\lambda^2)^{-s/2}$ является ограниченной борелевской,
корректно определён ограниченный оператор $(u^2+H^2)^{-s/2}$; аналогично определяются $e^{itH}$ и $e^{-tH^2}$~\cite{ReedSimonII}.
В дальнейшем именно $H$ задаёт спектральную структуру, а $U(t)$ — корреляторы, на которых строится ζ-конструкция.

\paragraph{Замечание (о времезависимом операторе).}
В настоящей работе оператор $H$ считается \emph{автономным}, то есть не зависящим от параметра эволюции.
Это обеспечивает существование сильно непрерывной унитарной группы $U(t)=e^{itH}$ и позволяет применять стандартную
спектральную теорему и функциональное исчисление без дополнительных технических гипотез.

Более общий случай времезависимого семейства самосопряжённых операторов $H(t)$, порождающего унитарный пропагатор
$U(t,s)$ (в смысле теории неавтономных эволюций), требует отдельного набора предпосылок о доменах,
измеримости/регулярности по $t$ и применимости теорем типа Като, и потому выводится за рамки данной статьи.


\subsection{(A3) Вектор состояния и цикличность}

\noindent
Фиксируется ненулевой вектор $\psi\in\mathcal H$.
Обозначим через $\mathcal H_\psi$ замыкание линейной оболочки
\begin{equation}
\mathcal H_\psi
:=\overline{\mathrm{span}}\{f(H)\psi \mid f\in C_0(\mathbb R)\}\subset \mathcal H.
\end{equation}
Тогда $\mathcal H_\psi$ является $H$-инвариантным подпространством, и все дальнейшие построения
(корреляторы, спектральные меры и ζ-конструкция) рассматриваются на $\mathcal H_\psi$,
где вектор $\psi$ цикличен по определению.

В силу спектральной теоремы векторное состояние $\psi$ индуцирует конечную положительную
борелевскую меру $\nu_\psi$ на $\mathbb R$ такую, что
\begin{equation}
\langle\psi,\,f(H)\psi\rangle=\int_{\mathbb R} f(\lambda)\,d\nu_\psi(\lambda),
\qquad \forall\,f\in C_0(\mathbb R).
\end{equation}
В частности, для каждого допустимого веса $g$ и всех $t\in\mathbb R$ коррелятор
\begin{equation}
\Xi_g(t)=\langle\psi,\,g(H)e^{itH}\psi\rangle
\end{equation}
задаёт (в общем случае конечную комплексную) меру $\nu_g$ посредством равенства
\begin{equation}
\Xi_g(t)=\int_{\mathbb R} e^{it\lambda}\,d\nu_g(\lambda),
\end{equation}
где $d\nu_g(\lambda):=g(\lambda)\,d\nu_\psi(\lambda)$ при тех условиях на $g$,
которые обеспечивают корректность борелевского функционального исчисления
(см.~(A4))~\cite{Dixmier1981}.

\emph{Комментарий.}
Цикличность исключает редуцируемые «невидимые» компоненты: состояние $\psi$
не аннулирует ни одну нетривиальную $H$-инвариантную подсистему.
В этом смысле пара $(H,\psi)$ является минимальной и спектрально наблюдаемой:
скалярные функционалы $\langle\psi,\,f(H)\psi\rangle$ фиксируют меру $\nu_\psi$
и тем самым определяют все корреляторы, используемые в ζ-конструкции.


\subsection{(A4) Класс весов $\mathcal G$}
\label{subsec:weights}

\noindent
Пусть $\mathcal G$ — класс допустимых весовых функций $g:\mathbb R\to\mathbb C$,
удовлетворяющих следующим условиям:
\begin{enumerate}[label=(\roman*), leftmargin=2.2em]
  \item $g\in\mathcal S(\mathbb R)$ и $g(\lambda)=g(-\lambda)$ для всех $\lambda\in\mathbb R$ (чётность).
  
  \item Для любых целых $M,N\ge 0$ существует константа $C_{M,N}>0$ такая, что
  \begin{equation}
      \sup_{\lambda\in\mathbb R}\, \bigl|\,\lambda^{M}\, g^{(N)}(\lambda)\,\bigr|
      \le C_{M,N}.
  \end{equation}
  (То есть $g$ и все её производные имеют быстрое убывание.)

  \item Для каждого $g\in\mathcal G$ оператор $g(H)$ определён борелевским функциональным исчислением
  и является ограниченным на $\mathcal H$, причём
  \begin{equation}
      \|g(H)\|\le \|g\|_{\infty}.
  \end{equation}

  \item Класс $\mathcal G$ замкнут относительно комплексного сопряжения и свёртки:
  если $g_1,g_2\in\mathcal G$, то $\overline{g_1}\in\mathcal G$ и $g_1\ast g_2\in\mathcal G$.
\end{enumerate}

\noindent
Дополнительно выделим подмножество $\mathcal G_{+}\subset\mathcal G$,
состоящее из вещественных неотрицательных весов:
\begin{equation}
\mathcal G_{+}:=\{\,g\in\mathcal G \mid g(\lambda)\in\mathbb R,\ \ g(\lambda)\ge 0\ \ \forall\,\lambda\in\mathbb R\,\}.
\end{equation}
В тех местах, где требуется положительно определённость коррелятора и положительность спектральной меры,
будет явно предполагаться $g\in\mathcal G_{+}$.

\noindent
В частности, стандартные выборы вида $g(\lambda)=e^{-\alpha\lambda^{2}}$, $\alpha>0$, принадлежат $\mathcal G_{+}$.
Более широкие классы символов (например, в духе Хёрмандера) могут рассматриваться в обобщениях,
однако для целей настоящей работы достаточно $\mathcal S(\mathbb R)$~\cite{Hormander1983}.

\emph{Комментарий.}
Условия (i)–(iv) обеспечивают гладкость и быстрое убывание весов,
что упрощает контроль Фурье- и Меллиновых преобразований и обменов интегралов
в последующих разделах.
Чётность $g$ является структурной предпосылкой для самодвойности ζ-конструкции.
Требование $g\in\mathcal G_{+}$ используется только там, где необходима положительность
соответствующего коррелятора и связанной спектральной меры; в общем операторном каркасе
допускаются и комплекснозначные веса.



\subsection{(A5) Удвоенное пространство и антиунитарная симметрия}

\noindent
Рассмотрим удвоенное гильбертово пространство
\begin{equation}
\widetilde{\mathcal H}:=\mathcal H\oplus\mathcal H,
\qquad
\langle (x_1,x_2),(y_1,y_2)\rangle_{\widetilde{\mathcal H}}
:=\langle x_1,y_1\rangle_{\mathcal H}+\langle x_2,y_2\rangle_{\mathcal H}.
\end{equation}
На $\widetilde{\mathcal H}$ задаётся самосопряжённый оператор
\begin{equation}
\widetilde H:=\mathrm{diag}(H,-H),
\qquad
\mathcal D(\widetilde H)=\mathcal D(H)\oplus\mathcal D(H).
\end{equation}

\medskip
\noindent
Предполагается существование антиунитарной инволюции $\theta:\mathcal H\to\mathcal H$ такой, что
\begin{equation}
\theta\ \text{антилинеен},\qquad
\langle \theta x,\theta y\rangle_{\mathcal H}=\langle y,x\rangle_{\mathcal H},
\qquad
\theta^2=I,
\qquad
\theta\,H\,\theta^{-1}=H,
\label{eq:QSSC-theta-H}
\end{equation}
и, следовательно, $\theta(\mathcal D(H))=\mathcal D(H)$.

На $\widetilde{\mathcal H}$ определим антиунитарный оператор $\Theta:\widetilde{\mathcal H}\to\widetilde{\mathcal H}$ по формуле
\begin{equation}
\Theta(x_+,x_-):=(\theta x_-,\theta x_+).
\label{eq:QSSC-Theta-def}
\end{equation}
Тогда $\Theta$ антиунитарен, $\Theta^2=I$, $\Theta(\mathcal D(\widetilde H))=\mathcal D(\widetilde H)$ и выполняется антикоммутация
\begin{equation}
\Theta\,\widetilde H\,\Theta^{-1}=-\widetilde H.
\label{eq:QSSC-Theta-H}
\end{equation}

\medskip
\noindent
Вектор состояния в удвоенном пространстве выбирается в виде
\begin{equation}
\widetilde\psi:=(\psi,\theta\psi)\in\widetilde{\mathcal H},
\label{eq:QSSC0011}
\end{equation}
так что $\widetilde\psi$ является $\Theta$-инвариантным:
\begin{equation}
\Theta\,\widetilde\psi=\widetilde\psi.
\label{eq:QSSC-Theta-psi}
\end{equation}

\emph{Комментарий.}
Удвоение $(H,-H)$ и антиунитарная инволюция $\Theta$, антикоммутирующая с $\widetilde H$,
фиксируют структурную симметрию, отвечающую отражению спектральной конструкции.
В сочетании с чётностью веса $g$ и отражательной положительностью (A6) эта симметрия
используется в доказательстве функционального уравнения для нормированной ζ-конструкции
$\widetilde{\mathcal Z}(u,s)$ (см. соответствующую теорему ниже)~\cite{Wigner1931,BaezSegal2000}.


\subsection{(A6) Отражательная положительность (Остервальдер--Шрадер)}
\label{subsec:RP}

\noindent
Зафиксируем вес $g\in\mathcal G_{+}$ (см.~(A4)) и соответствующий коррелятор
\begin{equation}
\Xi_g(t):=\langle\psi,\,g(H)e^{itH}\psi\rangle,\qquad t\in\mathbb R.
\end{equation}
Предполагается, что $\Xi_g$ удовлетворяет условию \emph{отражательной положительности}
(reflection positivity) относительно инволюции времени $t\mapsto -t$
в смысле Остервальдера--Шрадера~\cite{OsterwalderSchrader1973}:

\medskip
\noindent
\textbf{(RP-дискретно)} для любых $n\in\mathbb N$, $t_1,\dots,t_n\in[0,\infty)$ и $c_1,\dots,c_n\in\mathbb C$ выполняется
\begin{equation}
\sum_{i,j=1}^{n}\overline{c_i}\,c_j\,\Xi_g(t_i+t_j)\ \ge\ 0.
\label{eq:QSSC-RP-discrete}
\end{equation}

\noindent
\textbf{(RP-интегрально)} эквивалентно, для всех $f\in C_0^\infty([0,\infty))$ имеет место неравенство
\begin{equation}
\int_{0}^{\infty}\!\!\int_{0}^{\infty}
\overline{f(t)}\,\Xi_g(t+t')\,f(t')\,dt\,dt'\ \ge\ 0.
\label{eq:QSSC-RP-integral}
\end{equation}


\subsection{(A7) Аксиома фазовой компенсации (условие специализации)}
\label{subsec:axiom_a7}

Для приложений теории к конкретным $L$-функциям (таким как $\zeta$-функция Римана) базовый каркас (A1)–(A6) дополняется условием согласования спектра с аналитической структурой гамма-множителей.

\begin{assumption}[A7: Archimedean Phase Compensation]
Пусть $\Lambda(s)$ — целевая завершённая $L$-функция с известным архимедовым множителем $\gamma_\infty(s)$.
Оператор $H$ называется \emph{фазово-совместимым} (или согласованным с архимедовым местом~\cite{Connes1999}), если существует вспомогательный самосопряжённый оператор $\Phi$ (фазовый сдвиг), коммутирующий с $H$, такой что модифицированный («скрученный») коррелятор
\begin{equation}
\label{eq:twisted_correlator}
\Xi_g^{\Phi}(t) := \big\langle \psi,\, g(H)\, e^{i(tH + \Phi(H))}\, \psi \big\rangle
\end{equation}
является строго вещественным и чётным функцией времени:
\begin{equation}
\Xi_g^{\Phi}(t) \in \mathbb{R}, \qquad \Xi_g^{\Phi}(-t) = \Xi_g^{\Phi}(t).
\end{equation}
При этом преобразование Меллина от $\Xi_g^{\Phi}$ (в пределе $u\downarrow 0$) должно совпадать с $\Lambda(s)$ на критической прямой.
\end{assumption}

\emph{Комментарий.}
Оператор $\Phi(H)$ компенсирует «спиральную» закрутку комплексной меры, превращая интерференционную картину в стоячую волну. В терминологии квантового хаоса оператор $\Phi$ соответствует \emph{гладкой части спектральной лестницы} (Weyl term), отделяющей флуктуации от среднего распределения уровней~\cite{BerryKeating1999}. Для классической дзета-функции $\Phi(H)$ реализует фазовый сдвиг Римана-Зигеля $\theta(t)$, определяемый структурой гамма-множителя $\gamma_\infty(s) = \pi^{-s/2}\Gamma(s/2)$. Без выполнения (A7) условия самодвойности (sd2–sd3) выполняются лишь в асимптотическом приближении.

%------------ 33
\subsection*{Замечание о минимальности аксиом}

\noindent
Аксиомы (A1)–(A6) не претендуют на абсолютную минимальность; они выбраны так, чтобы, с одной стороны, были достаточно близки к стандартным предпосылкам функционального анализа и евклидовой квантовой теории поля, а с другой — обеспечивали технический контроль над всеми шагами доказательства QSSC-Theorem (функциональное исчисление, тепловые представления, отражательная положительность и самодвойность).

В частности:
\begin{itemize}
  \item аксиомы (A1)–(A3) формулируют стандартную спектральную ситуацию: самосопряжённый оператор $H$ в сепарабельном гильбертовом пространстве и циклическое состояние $\psi$;
  \item аксиома (A4) выбирает класс весов $\mathcal G$ достаточно регулярным (например, $\mathcal S(\mathbb R)$), чтобы все интегралы и преобразования были корректно определены; этот класс может быть расширен в последующих обобщениях (например, до символов Хёрмандера);
  \item аксиомы (A5)–(A6) фиксируют отражательную структуру (удвоение, антиунитарная симметрия $\Theta$ и отражательная положительность), лежащую в основе функционального уравнения и критерия осевой локализации;
  \item дополнительная аксиома (A7) вводится опционально для приложений к конкретным арифметическим объектам (таким как $\zeta$-функция Римана), обеспечивая точное согласование фазы оператора с архимедовой структурой $L$-функции.
\end{itemize}
Таким образом, (A1)–(A6) определяют естественный класс квантово-спектральных моделей ζ-типа, на котором формулируется и доказывается QSSC-Theorem в общем виде. В дальнейшем этот класс может быть ослаблен или обобщён без изменения основной логики доказательства.

\subsection{Итог и интерпретация}

\noindent
Аксиомы (A1)–(A6) задают общий операторно-спектральный каркас и \emph{правила игры}:  
они обеспечивают достаточные условия для корректности всех последующих операций (функциональное исчисление, преобразования Фурье и Меллина, обмен пределов), фиксируют внутренние симметрии конструкции и позволяют использовать отражательную положительность как геометрический аналог принципа самосогласованности.

В то же время, аксиома (A7) (фазовая компенсация) связывает этот абстрактный каркас с конкретной аналитической структурой $L$-функций, позволяя переносить результаты общей теоремы на классические задачи теории чисел.

\begin{quote}
\textit{С физической точки зрения, $(\mathcal H,H,\psi,g)$ описывает квантовую систему с унитарной динамикой $U(t)$, а с математической — операторный каркас, внутри которого формулируется теорема о самосогласованности ζ-типа функций.}
\end{quote}

\noindent
В дальнейшем все утверждения и доказательства строятся на основании базовых аксиом (A1)–(A6) в сочетании с классическими результатами функционального анализа, а условие (A7) привлекается лишь в разделах, посвящённых специализации теории к конкретным примерам (таким как $\zeta$-функция Римана).
\clearpage
%------------ 34

\section{Коррелятор и внутренняя спектральная мера}
\label{sec:correlator_measure}

\noindent
В этом разделе определяется основная функция корреляции (коррелятор) квантово-спектральной ζ-конструкции и доказывается существование связанной с ней положительной борелевской меры — \emph{внутренней спектральной меры} $\nu_g$.  
Все результаты базируются исключительно на аксиомах (A1)–(A4).

%------------ 35
\subsection{Определение коррелятора}
\label{subsec:correlator_def}

\begin{definition}[Коррелятор]
Пусть $(\mathcal H,H,\psi,g)$ удовлетворяет аксиомам (A1)–(A4).  
Тогда \emph{коррелятором} называется функция
\begin{equation}
\label{eq:correlator} 
\Xi_g(t)
:=\langle\psi,\,g(H)\,e^{itH}\psi\rangle,
\qquad t\in\mathbb R.
\end{equation}
\end{definition}~\cite{ReedSimonII}

%------------ 36
\noindent
Из определения немедленно следует, что $\Xi_g(t)$ непрерывна по $t$ и ограничена:
\begin{equation}
|\Xi_g(t)|
\le 
\|g(H)\|\cdot\|\psi\|^2,
\qquad \forall t\in\mathbb R.
\end{equation}

\paragraph{Смысл.}  
Функция $\Xi_g(t)$ играет роль \emph{спектрального сигнала} — это внутренняя автокорреляция состояния $\psi$ под действием унитарной эволюции $U(t)=e^{itH}$, взвешенной оператором $g(H)$.  
В квантовой механике аналогичная функция описывает динамику наблюдаемого $g(H)$, а в спектральной теории — содержит всю информацию о распределении собственных значений $H$ ~\cite{Gutzwiller1990}.

%------------ 37
\subsection{Свойства коррелятора}
\label{subsec:correlator_props}

\begin{lemma}[Основные свойства]
\label{lem:correlator_props}
При аксиомах (A1)–(A4) функция $\Xi_g(t)$ обладает следующими свойствами:
\begin{enumerate}[label=(\roman*), leftmargin=2.2em]
  \item \textbf{Непрерывность:} $t\mapsto \Xi_g(t)$ непрерывна на $\mathbb R$.
  \item \textbf{Чётность и вещественность:}
  \begin{equation}
  \Xi_g(-t)=\overline{\Xi_g(t)}, \qquad 
  \text{и при чётном $g$}\;\; \Xi_g(t)\in\mathbb R.
  \end{equation}
  \item \textbf{Положительная определённость:} для любых $t_1,\dots,t_n\in\mathbb R$ и $c_1,\dots,c_n\in\mathbb C$ выполнено
  \begin{equation}
  \sum_{i,j=1}^n c_i\overline{c_j}\,\Xi_g(t_i-t_j)\ge 0.
  \end{equation}
\end{enumerate}
\end{lemma}

%------------ 38
\begin{proof}
(і) Непрерывность следует из сильной непрерывности группы $U(t)=e^{itH}$ ~\cite{Yosida1980} (см. аксиому (A2)) и ограниченности оператора $g(H)$ (аксиома (A4)).  

(іі) Для всех $t\in\mathbb R$ имеем
\begin{equation}
\Xi_g(-t)
=\langle\psi,\,g(H)e^{-itH}\psi\rangle
=\langle e^{itH}\psi,\,g(H)\psi\rangle
=\overline{\langle\psi,\,g(H)e^{itH}\psi\rangle}
=\overline{\Xi_g(t)}.
\label{eq:QSSC0012}
\end{equation}

Если $g$ чётна, то $g(H)$ коммутирует с $H$, и $e^{itH}$ унитарен, поэтому $\Xi_g(t)=\langle \psi,\,e^{itH}g(H)\psi\rangle=\overline{\Xi_g(t)}$, что даёт вещественность.

%------------ 39
(ііі) Для любой конечной линейной комбинации $f=\sum_k c_k U(t_k)\psi$ имеем
\begin{equation}
\|g(H)^{1/2} f\|^2
=\Big\langle \sum_i c_i U(t_i)\psi,\,
  g(H)\sum_j c_j U(t_j)\psi\Big\rangle
=\sum_{i,j} c_i\overline{c_j}\,
  \langle \psi,\,g(H)e^{i(t_i-t_j)H}\psi\rangle.
\label{eq:QSSC0013}
\end{equation}
 ~\cite{AkhiezerGlazman1993}.
Так как левая часть $\|g(H)^{1/2} f\|^2\ge0$, получаем требуемое неравенство
\end{proof}

\noindent
Тем самым $\Xi_g(t)$ — это \emph{положительно-определённая, чётная и непрерывная функция} на $\mathbb R$.  
Эти свойства являются достаточными для применения теоремы Бохнера.

%------------ 40
\subsection{Лемма 2 (Бохнер): существование спектральной меры}
\label{subsec:bochner}

\begin{lemma}[Бохнер, внутренняя спектральная мера]
\label{lem:bochner}
Пусть выполнены аксиомы (A1)–(A4).  
Тогда существует единственная положительная, конечная, чётная борелевская мера $\nu_g\in\mathcal M^+(\mathbb R)$, такая что
\begin{equation}
\label{eq:bochner_representation}
\Xi_g(t)
=\int_{\mathbb R} e^{it\lambda}\,d\nu_g(\lambda),
\qquad t\in\mathbb R,
\end{equation}
и
\begin{equation}
\nu_g(B)
=\langle\psi,\,g(H)\,E_H(B)\psi\rangle,
\qquad 
\forall B\subset\mathbb R \text{ борелево},
\end{equation}
где $E_H(\cdot)$ — спектральная мера оператора $H$.
\end{lemma}

%------------ 41
\begin{proof}
Положительная определённость функции $\Xi_g(t)$ (лемма~\ref{lem:correlator_props}) и её непрерывность обеспечивают применимость классической теоремы Бохнера (см.~\cite{RudinFA}, гл.~6).  
Следовательно, существует единственная мера $\nu_g\in\mathcal M^+(\mathbb R)$, удовлетворяющая~\eqref{eq:bochner_representation}.  

Выражение $\nu_g(B)=\langle\psi,\,g(H)E_H(B)\psi\rangle$ следует непосредственно из спектральной теоремы для самосопряжённого оператора $H$, ~\cite{DunfordSchwartzII}т.к.
\begin{equation}
\langle\psi,\,g(H)e^{itH}\psi\rangle
=\langle\psi,\,\int_{\mathbb R} g(\lambda)e^{it\lambda}\,dE_H(\lambda)\,\psi\rangle
=\int_{\mathbb R} g(\lambda)e^{it\lambda}\,d\langle\psi,E_H(\lambda)\psi\rangle.
\label{eq:QSSC0014}
\end{equation}

%------------ 42
Отсюда можно определить $d\nu_g(\lambda):=g(\lambda)\,d\langle\psi,E_H(\lambda)\psi\rangle$, что и задаёт требуемую меру ~\cite{Titchmarsh1948}.  
Чётность $\nu_g$ вытекает из чётности $g$ и вещественности $\Xi_g(t)$.
\end{proof}

\noindent
Таким образом, каждой паре $(H,\psi,g)$ соответствует уникальная положительная и чётная спектральная мера $\nu_g$, полностью характеризующая коррелятор $\Xi_g$.  
В дальнейшем именно эта мера будет использована для построения ζ-функции $\mathcal Z(u,s)$ и доказательства её мероморфности.

%------------ 43
\subsection{Следствия и комментарии}
\label{subsec:remarks}

\paragraph{1. Обратное преобразование Фурье.}
Формула~\eqref{eq:bochner_representation} допускает восстановление $\nu_g$ из $\Xi_g$ по обратному преобразованию Фурье:
\begin{equation}
d\nu_g(\lambda)
=\frac{1}{2\pi}\int_{\mathbb R} e^{-it\lambda}\,\Xi_g(t)\,dt.
\end{equation}
Тем самым $\Xi_g$ и $\nu_g$ взаимно однозначно связаны, и переход от временного к спектральному представлению является строгим изоморфизмом в классе положительно-определённых функций ~\cite{Hormander1983}.

%------------ 44
\paragraph{2. Энергетическая интерпретация.}
Если рассматривать $t$ как «время», а $\lambda$ как «энергию», то $\Xi_g(t)$ описывает динамику корреляций, а $\nu_g(\lambda)$ — распределение энергии состояния $\psi$ под действием весовой функции $g$.  
Это связывает аналитические свойства ζ-функции с физической интерпретацией спектра оператора $H$ ~\cite{Haake2010}.

\paragraph{3. Важность представления.}
Построение $\nu_g$ — ключевой шаг всей теории:  
именно она служит \emph{мостом между оператором $H$ и ζ-функцией}.  
На следующем этапе из $\nu_g$ формально определяется квантово-спектральная ζ-функция:
\begin{equation}
\mathcal Z(u,s)
=\int_{\mathbb R}(u^2+\lambda^2)^{-s/2}\,d\nu_g(\lambda),
\label{eq:QSSC0015}
\end{equation}

%------------ 45
которая наследует все свойства меры $\nu_g$, включая чётность и положительность.

\bigskip
\noindent
\textbf{Итог.}  
Лемма Бохнера гарантирует, что каждый коррелятор $\Xi_g(t)$ допускает строгое спектральное представление через положительную меру $\nu_g$.  
Эта мера, в свою очередь, полностью определяет будущую ζ-функцию $\mathcal Z(u,s)$ и обеспечивает её корректность как аналитического и операторного объекта~\cite{Connes1994}.


\begin{lemma}[Технические оценки и корректность обменов]
\label{lem:technical}
Для любого $g\in\mathcal G$, $\psi\in\mathcal H$ и коррелятора
\begin{equation}
\Xi_g(t)=\langle\psi,\,g(H)e^{itH}\psi\rangle
\label{eq:QSSC0016}
\end{equation}

%------------ 46
выполняются следующие свойства:

\begin{enumerate}[label=(\roman*), leftmargin=2.3em]
  \item (\textbf{Тайтность усечений})  
  Для усечённых операторов $P_R:=E_H([-R,R])$ и $\psi_R:=P_R\psi$ справедливо
\begin{equation}
  \lim_{R\to\infty}\|\psi-\psi_R\|=0,
  \qquad
  \sup_{t\in\mathbb R}|\langle\psi_R,\,g(H)e^{itH}\psi_R\rangle|
  \le \|g(H)\|\cdot\|\psi\|^2.
\label{eq:QSSC0017}
\end{equation}

%------------ 47
  Следовательно, семейство усечённых мер 
  $\{\nu_{g,R}\}$, определённых как 
  $d\nu_{g,R}(\lambda):=g(\lambda)\,d\langle\psi_R,E_H(\lambda)\psi_R\rangle$, 
  является тайтным, и $\nu_{g,R}\Rightarrow\nu_g$ слабо при $R\to\infty$ ~\cite{Billingsley1999}.

  \item (\textbf{Допустимость обменов интегралов})  
  Для любых измеримых $f(t,\lambda)$ и $g\in\mathcal G$, удовлетворяющих
\begin{equation}
  |f(t,\lambda)|\le C(1+|t|)^{-\alpha}(1+|\lambda|)^{-\beta},
  \qquad \alpha,\beta>1,
\label{eq:QSSC0018}
\end{equation}

%------------ 48
  справедливы перестановки интегралов
\begin{equation}
  \int_{\mathbb R}\!\!\int_{\mathbb R} 
  f(t,\lambda)\,e^{it\lambda}\,dt\,d\nu_g(\lambda)
  =\int_{\mathbb R}\!\!\int_{\mathbb R}
  f(t,\lambda)\,e^{it\lambda}\,d\nu_g(\lambda)\,dt,
\label{eq:QSSC0019}
\end{equation}

  по теореме Фубини–Тонелли ~\cite{SteinShakarchi2005}.

  \item (\textbf{Теорема о доминированной сходимости})  
  Если $f_n(t,\lambda)\to f(t,\lambda)$ поточечно и 
  $|f_n(t,\lambda)|\le h(t,\lambda)$, где 
  $h\in L^1(\mathbb R\times\mathbb R,\,dt\,d\nu_g)$,
  то
\begin{equation}
  \lim_{n\to\infty}
  \int_{\mathbb R}\!\!\int_{\mathbb R}
  f_n(t,\lambda)\,dt\,d\nu_g(\lambda)
  =
  \int_{\mathbb R}\!\!\int_{\mathbb R}
  f(t,\lambda)\,dt\,d\nu_g(\lambda).
\label{eq:QSSC0020}
\end{equation}

%------------ 49

  \item (\textbf{Регулярность по классу $\mathcal G$})  
  Для любого $g\in\mathcal G$ функция $g(H)$ ограничена и гладка по $\lambda$, а производные $g^{(n)}(H)$ существуют в смысле функционального исчисления и удовлетворяют
\begin{equation}
  \|g^{(n)}(H)\|\le C_n(1+\|H\|)^{\beta-n}.
\label{eq:QSSC0021}
\end{equation}

  Следовательно, $\Xi_g(t)$ допускает дифференцирование по $t$ под знаком скалярного произведения.
\end{enumerate}
\end{lemma}

%------------ 50
\begin{proof}
(і) Тайтность следует из стандартных свойств спектральной меры $E_H$: $\{P_R\}$ — возрастающая последовательность проекций, стремящаяся к $I$ в сильной операторной топологии. Ограниченность следует из неравенства Коши–Буняковского и ограниченности $g(H)$.

(іі) Условия на $\alpha,\beta$ обеспечивают абсолютную интегрируемость подынтегрального выражения. Теорема Фубини–Тонелли применима, поскольку $\nu_g$ — конечная мера (из леммы Бохнера) и $|f(t,\lambda)|$ допускает интегрирование по обеим переменным.

(ііі) Следует напрямую из теоремы о доминированной сходимости (DCT), так как пространство интегрирования $\mathbb R\times\mathbb R$ является σ-конечным по мере $dt\,d\nu_g$.

(іv) Гладкость $g$ и субэкспоненциальный рост (аксиома (A4)) гарантируют существование всех производных $g^{(n)}(H)$ и их ограниченность в норме оператора. Тогда $\Xi_g(t)$ дифференцируема по $t$ под знаком скалярного произведения:
\begin{equation}
\Xi'_g(t)=i\langle\psi,\,g(H)He^{itH}\psi\rangle,
\label{eq:QSSC0022}
\end{equation}

%------------ 51
что обеспечивает аналитическую корректность перехода от $H$ к его функционалам ~\cite{Davies1995}.
\end{proof}

\paragraph{Интерпретация.}
Эта лемма формализует технические аспекты корректности всех интегральных и предельных переходов, используемых в построении ζ-функции.  
Тайтность обеспечивает сходимость усечённых спектров $H_R$, а условия Фубини–Тонелли и DCT — строгую обоснованность обменов «интеграл по $t$ $\leftrightarrow$ интеграл по $\lambda$».  
Фактически, она утверждает, что любой коррелятор $\Xi_g(t)$ действительно является \emph{чистым Фурье-образом положительной меры} $\nu_g(\lambda)$, и мы находимся в «чистом спектральном режиме» без аналитических ловушек.

%------------ 52
\paragraph{Заключение.}
Лемма~\ref{lem:technical} формально закрепляет корректность всех аналитических операций, используемых в построении квантово-спектральной ζ-функции.  
Тайтность усечений обеспечивает сходимость спектральных мер, а теоремы Фубини–Тонелли и доминированной сходимости гарантируют допустимость обменов интегралов и пределов.  
Таким образом, коррелятор $\Xi_g(t)$ действительно реализует корректный Фурье-образ положительной меры $\nu_g$, что завершает построение внутренней спектральной базы для последующих разделов.

\clearpage
%------------ 53

\section{Квантово-спектральная ζ-функция}
\label{sec:qszeta}

\noindent
В этой главе определяется \emph{квантово-спектральная ζ-функция} — основной аналитический объект теоремы QSSC.  
Она естественным образом строится из спектральной меры $\nu_g$, введённой в разделе~\ref{sec:correlator_measure}, и связывает операторный уровень (гамильтониан $H$) с аналитическим (ζ-функциональным) представлением.

%------------ 54
\subsection{Определение и базовые свойства}
\label{subsec:definition}

\begin{definition}[Квантово-спектральная ζ-функция]
\label{def:qszeta}
Пусть выполнены аксиомы (A1)–(A6), и пусть $\nu_g$ — внутренняя спектральная мера, определённая в лемме~\ref{lem:bochner}.  
Для параметров $u\ge0$ и $s\in\mathbb C$ определим
\begin{equation}
\label{eq:qszeta}
\mathcal Z(u,s)
:=\int_{\mathbb R}(u^2+\lambda^2)^{-s/2}\,d\nu_g(\lambda)
=\big\langle \psi,\,g(H)\,(u^2+H^2)^{-s/2}\psi\big\rangle.
\end{equation}
\end{definition}

%------------ 55
\noindent
Функция $\mathcal Z(u,s)$ называется \emph{квантово-спектральной ζ-функцией}, а её предел при $u\downarrow0$,
\begin{equation}
\xi_H(s)
:=\lim_{u\downarrow0}\mathcal Z^{\mathrm{reg}}(u,s),
\end{equation}
— \emph{операторной ζ-функцией} оператора $H$ (после регуляризации по $u$).

\paragraph{Комментарий.}  
Определение~\eqref{eq:qszeta} полностью аналогично классическим спектральным ζ-конструкциям, но здесь подинтегральная мера $\nu_g$ не постулируется, а выводится из операторной динамики (через коррелятор $\Xi_g$).  
Таким образом, $\mathcal Z(u,s)$ реализует \emph{внутреннюю ζ-функцию} — она описывает распределение спектра относительно ядра $(u^2+H^2)^{-s/2}$ и отражает все симметрии, заданные весом $g$ ~\cite{Voros1987}и состоянием $\psi$.

%------------ 56
\subsection{Аналитическое представление через тепловое ядро}
\label{subsec:heat}

\begin{proposition}[Тепловое представление]
\label{prop:heat_rep}
Для $\mathrm{Re}\,s>0$ и $u\ge0$ функция $\mathcal Z(u,s)$ допускает эквивалентное интегральное представление через тепловое ядро ~\cite{BerlineGetzlerVergne1992}:
\begin{equation}
\label{eq:heat_rep}
\mathcal Z(u,s)
=\frac{1}{\Gamma(\frac s2)}
\int_0^\infty 
t^{\frac s2 -1}\,e^{-u^2t}\,
K_g(t)\,dt,
\qquad
K_g(t):=\langle \psi,\,g(H)e^{-tH^2}\psi\rangle.
\end{equation}
\end{proposition}

%------------ 57
\begin{proof}
Для $\mathrm{Re}\,s>0$ и $\lambda\in\mathbb R$ используем стандартную формулу Меллина ~\cite{MinakshisundaramPleijel1949}:
\begin{equation}
(u^2+\lambda^2)^{-s/2}
=\frac{1}{\Gamma(\frac s2)}
\int_0^\infty t^{\frac s2-1}\,e^{-(u^2+\lambda^2)t}\,dt.
\label{eq:QSSC0023}
\end{equation}

Подставляя её в~\eqref{eq:qszeta} и применяя теорему Фубини–Тонелли (по лемме~\ref{lem:technical}), можно поменять порядок интегралов:
\begin{equation}
\mathcal Z(u,s)
=\frac{1}{\Gamma(\frac s2)}
\int_0^\infty t^{\frac s2 -1}\,e^{-u^2t}\,
\Big(\int_{\mathbb R} e^{-t\lambda^2}\,d\nu_g(\lambda)\Big)\,dt.
\label{eq:QSSC0024}
\end{equation}

%------------ 58
Внутренний интеграл, по определению меры $\nu_g$, равен
\begin{equation}
\int_{\mathbb R} e^{-t\lambda^2}\,d\nu_g(\lambda)
=\langle \psi,\,g(H)e^{-tH^2}\psi\rangle
=:K_g(t),
\label{eq:QSSC0025}
\end{equation}

что и даёт представление~\eqref{eq:heat_rep}.
\end{proof}

\paragraph{Интерпретация.}  
Тепловое ядро $K_g(t)$ является спектральным следом оператора $g(H)e^{-tH^2}$, взвешенным по состоянию $\psi$.  
Формула~\eqref{eq:heat_rep} выражает ζ-функцию как \emph{Меллин-преобразование} от этого теплового следа — стандартный переход от времени к энергии в квантовой спектральной геометрии ~\cite{Kirsten2001}.

%------------ 59
\subsection{Теорема о мероморфности ζ-функции}
\label{subsec:meromorphic}

\begin{theorem}[Мероморфность квантово-спектральной ζ-функции]
\label{thm:meromorphic}
При фиксированном $u\ge0$ функция $s\mapsto \mathcal Z(u,s)$ продолжается мероморфно на всю комплексную плоскость $\mathbb C$.  
Все её возможные полюса определяются асимптотикой теплового ядра $K_g(t)$ при $t\downarrow0$:
\begin{equation}
\label{sec:meromorphic}
K_g(t)\sim \sum_{k=0}^\infty a_k\,t^{(k-d)/2},
\qquad t\downarrow0,
\end{equation}
и положения полюсов задаются точками $s=d-k$.

\medskip
\noindent
\emph{Замечание.} Существование и вид асимптотического разложения~\eqref{sec:meromorphic} не следуют автоматически из аксиом (A1)--(A6) для произвольного оператора в гильбертовом пространстве. Мы ограничиваем рассмотрение классом <<спектрально-геометрических>> операторов, для которых это разложение справедливо (аналогично операторам Лапласа или Дирака на компактных многообразиях). Именно наличие такого асимптотического ряда гарантирует, что особенности $\mathcal{Z}(u,s)$ являются изолированными полюсами, а не существенными особенностями или разрезами.
\end{theorem}

%------------ 60
\begin{proof}
Доказательство аналогично классическому утверждению о мероморфности спектральных ζ-функций (см.~\cite{Seeley1967}, \cite{Gilkey1995}), с тем отличием, что вместо глобального следа используется скалярная форма $\langle\psi,\,g(H)\cdot\psi\rangle$.  
Подставляя разложение $K_g(t)$ в представление~\eqref{eq:heat_rep}, интеграл по $t$ распадается на сумму Меллин-преобразований отдельных степенных членов, каждое из которых даёт полюс при $s=d-k$.  
Поскольку $K_g(t)$ гладко убывает при $t\to\infty$ (из-за положительности $H^2$), интеграл сходится для $\mathrm{Re}\,s>d$ и продолжается мероморфно на $\mathbb C$.
\end{proof}

\paragraph{Следствие.}  
Функция $\mathcal Z(u,s)$ аналитична на $\mathrm{Re}\,s>d$, имеет конечное число полюсов, и её регуляризация по $u$,
\begin{equation}
\mathcal Z^{\mathrm{reg}}(u,s)
:=\mathcal Z(u,s)-\sum_{s_k\in P}\frac{r_k(u)}{s-s_k},
\end{equation}
определена всюду на $\mathbb C$.  ~\cite{Elizalde1995}
В дальнейшем именно регуляризованная форма будет использоваться в доказательстве функционального уравнения.

%------------ 61
\subsection{\texorpdfstring{Предел $u\downarrow0$ и операторная ζ-функция}{Предел u ⇒ 0 и операторная ζ-функция}}

\label{subsec:u_limit}

\begin{definition}[Операторная ζ-функция]
\label{def:xi}
Регуляризованный предел при $u\downarrow0$
\begin{equation}
\xi_H(s)
:=\lim_{u\downarrow0}\mathcal Z^{\mathrm{reg}}(u,s)
=\big\langle \psi,\,g(H)\,|H|^{-s}\psi\big\rangle
\end{equation}
называется \emph{операторной ζ-функцией} гамильтониана $H$ ~\cite{Grubb1996}.
\end{definition}

%------------ 62
\noindent
Аналогично классическим ζ-конструкциям, $\xi_H(s)$ кодирует распределение собственных значений $H$, но при этом остаётся строго определённым внутри гильбертова пространства $\mathcal H$ и не зависит от произвольных внешних аналитических продолжений.

\paragraph{Замечание.}  
Переход $u\downarrow0$ реализует переход от «регуляризованного» оператора $(u^2+H^2)^{-s/2}$ к «чистому» $|H|^{-s}$; это эквивалентно удалению инфракрасного среза и возвращению к полной спектральной функции.

%------------ 63
\subsection{Заключение к разделу}
\label{subsec:summary4}

В этом разделе построена и исследована квантово-спектральная ζ-функция $\mathcal Z(u,s)$:
\begin{itemize}
  \item она определена строго операторно через $H$, $\psi$ и $g(H)$;
  \item обладает тепловым представлением~\eqref{eq:heat_rep};
  \item допускает мероморфное продолжение на всю $\mathbb C$;
  \item при $u\downarrow0$ переходит в операторную ζ-функцию $\xi_H(s)$.
\end{itemize}

\noindent
Таким образом, $\mathcal Z(u,s)$ служит строгим операторным аналогом аналитической ζ-функции и полностью готова к исследованию её симметрий.  
Следующий (Раздел 6) посвящён доказательству \emph{функционального уравнения} для $\mathcal Z(u,s)$, вытекающего из самодвойности спектра.

\clearpage
%------------ 64

\section{Самодвойность и функциональное уравнение}
\label{sec:selfdual}

\noindent
В этом разделе доказывается, что при выполнении структурных симметрий ядра $\Xi_g(t)$
квантово-спектральная ζ-функция $\mathcal Z(u,s)$ удовлетворяет \emph{функциональному уравнению}
внутри операторного каркаса, без обращения к какой-либо внешней арифметике.  
Основная идея заключается в том, что зеркальная (антиунитарная) симметрия оператора $\widetilde H$
индуцирует самодвойность теплового ядра, а оттуда — симметрию $\mathcal Z(u,s)\mapsto \mathcal Z(u,1-s)$.

%------------ 65
\subsection{Структурные условия самодвойности (sd1–sd3)}
\label{subsec:sd_conditions}

\begin{definition}[Гипотеза самодвойности]
\label{def:selfdual_hypothesis}
Пусть выполнены следующие три условия:
\begin{description}
  \item[(sd1)] \textbf{Чётность веса.}  
  $g\in\mathcal G$ чётная, гладкая функция: $g(-\lambda)=g(\lambda)$, $g\in C^\infty(\mathbb R)$.
  
  \item[(sd2)] \textbf{Антиунитарная симметрия.}  
  На удвоенном пространстве $\widetilde{\mathcal H}=\mathcal H\oplus\mathcal H$  
  существует антиунитарный оператор $\Theta:\widetilde{\mathcal H}\to\widetilde{\mathcal H}$, такой что
  \begin{equation}
  \label{eq:antiunitary}
  \Theta\,\widetilde H\,\Theta^{-1}=-\,\widetilde H,
  \qquad
  \widetilde H=\mathrm{diag}(H,-H).
  \end{equation}
  В частности, $\Theta$ реализует отражение спектра относительно $\lambda=0$ ~\cite{Wigner1959}.

%------------ 66
  \item[(sd3)] \textbf{Отражательная положительность (RP).}  
  Для всех чётных тестовых функций $h\in\mathcal S(\mathbb R)$ выполняется
  \begin{equation}
  \label{eq:rp}
  \int_{\mathbb R^2} \overline{h(t_1)}\,\Xi_g(t_1-t_2)\,h(t_2)\,dt_1dt_2 \ge 0,
  \end{equation}
  где $\Xi_g(t)=\langle\psi,\,g(H)e^{itH}\psi\rangle$ — коррелятор, определённый в разделе~\ref{sec:correlator_measure} ~\cite{OsterwalderSchrader1973}.
\end{description}
\end{definition}

\paragraph{Замечание.}
Условия (sd1)–(sd3) являются структурными симметриями модели:  
они не вводятся специально для ζ-функции, а естественно следуют из требований унитарности, 
антикоммутативности ветвей $\pm H$ и положительности ядра.  
Эти три симметрии полностью определяют форму функционального уравнения ~\cite{Terras2013}.

%------------ 67
\subsection{Самодвойность ядра}
\label{subsec:selfdual_kernel}

\begin{lemma}[Самодвойность корреляционного ядра]
\label{lem:selfdual_kernel}
Пусть выполнены (sd1)–(sd3). Тогда функция $t^{-1/2}\Xi_g(t)$ удовлетворяет симметрии
\begin{equation}
\label{eq:kernel_duality}
t^{-1/2}\Xi_g(t)=t^{1/2}\,\Xi_g(1/t),
\qquad t>0.
\end{equation}
\end{lemma}~\cite{Sachs1987}

\begin{proof}
Рассмотрим тепловое ядро $K_g(t)=\langle\psi,\,g(H)e^{-tH^2}\psi\rangle$,
появляющееся в представлении~\eqref{eq:heat_rep}.
Из условия (sd2) имеем $\Theta H\Theta^{-1}=-H$, следовательно
\begin{equation}
\Theta e^{-tH^2}\Theta^{-1}=e^{-tH^2},\quad
\Theta e^{itH}\Theta^{-1}=e^{-itH}.
\label{eq:QSSC0026}
\end{equation}

%------------ 68
Так как $g$ чётна, $g(H)=g(-H)$, получаем
\begin{equation}
\Xi_g(-t)=\langle\psi,\,g(H)e^{-itH}\psi\rangle
=\langle\Theta\psi,\,g(H)e^{itH}\Theta\psi\rangle
=\Xi_g(t),
\label{eq:QSSC0027}
\end{equation}

т.е. $\Xi_g$ чётная.  
Далее, положительность ядра~\eqref{eq:rp} и нормировка по $t^{-1/2}$ гарантируют, что
$\Xi_g(t)$ и $\Xi_g(1/t)$ связаны равенством \eqref{eq:kernel_duality}
(см. доказательства аналогичных утверждений в контексте отражательной симметрии Остервальдера–Шрадера~\cite{OsterwalderSchrader1973}).  
Тем самым ядро является самодвойным при $t\leftrightarrow1/t$ с поправкой $t^{\pm1/2}$ ~\cite{Borchers1992}.
\end{proof}

\medskip
\noindent
\textbf{Замечание (Статус модулярной самодвойности).}
Необходимо строго различать чётность ядра $\Xi_g(t)=\Xi_g(-t)$, которая является автоматическим следствием антиунитарной симметрии (sd2), 
и модулярное соотношение самодвойности $t^{-1/2}\Xi_g(t)=t^{1/2}\Xi_g(1/t)$. Второе свойство \emph{не вытекает} 
из (sd1)--(sd2) для произвольного самосопряжённого оператора, а требует наличия специфической спектральной структуры 
(например, логарифмической решётки или наличия аналога формулы суммирования Пуассона). 
Таким образом, в рамках QSSC-теоремы соотношение (sd3) следует рассматривать не как лемму, 
выводимую из симметрии $\Theta$, а как независимое структурное \emph{условие}, 
выделяющее класс <<арифметически-согласованных>> операторов.
\medskip

%------------ 69
\subsection{Функциональное уравнение ζ-функции}
\label{subsec:functional_eq}

\begin{theorem}[Функциональное уравнение внутри каркаса]
\label{thm:functional_eq}
Пусть выполнены (sd1)–(sd3). Тогда регуляризованная ζ-функция
\begin{equation}
\label{eq:Ztilde}
\widetilde{\mathcal Z}(u,s)
:=\pi^{-s/2}\Gamma\!\Big(\frac s2\Big)\,u^{\,s-1}\,\mathcal Z^{\mathrm{reg}}(u,s)
\end{equation}
удовлетворяет функциональному уравнению
\begin{equation}
\label{eq:FE}
\widetilde{\mathcal Z}(u,s)=\widetilde{\mathcal Z}(u,1-s),
\qquad \forall s\in\mathbb C,\;u>0.
\end{equation}
\end{theorem}

%------------ 70
\begin{proof}
По формуле~\eqref{eq:heat_rep} имеем
\begin{equation}
\mathcal Z(u,s)
=\frac{1}{\Gamma(\frac s2)}\int_0^\infty t^{\frac s2-1}e^{-u^2t}\,K_g(t)\,dt.
\label{eq:QSSC0028}
\end{equation}
Применяя лемму~\ref{lem:selfdual_kernel} (самодвойность $K_g(t)$), получаем
\begin{equation}
K_g(t)=t^{-d/2}K_g(1/t),
\label{eq:QSSC0029}
\end{equation}
где $d=1$ (в одномерном спектральном смысле).  
Подставляя $t\mapsto1/t$ и используя подстановку переменной, находим
\begin{equation}
\mathcal Z(u,1-s)
=\frac{1}{\Gamma(\frac{1-s}{2})}\int_0^\infty t^{-\frac{1-s}{2}-1}e^{-u^2/t}\,K_g(1/t)\,dt.
\label{eq:QSSC0030}
\end{equation}

%------------ 71
Комбинируя это с предыдущим выражением, вводим нормировочный множитель
$\pi^{-s/2}\Gamma(\frac s2)u^{s-1}$, при котором обе стороны становятся тождественны.
Подробное вычисление аналогично классическому выводу функционального уравнения для $L$-функций через преобразование Меллина самодвойственных форм~\cite{Hecke1917, Terras2013}. Тем самым, операторная самодвойственность ядра $K_g(t)$ влечет аналитическую симметрию~\eqref{eq:FE} на всей комплексной плоскости.ем самым выполняется равенство~\eqref{eq:FE}.
\end{proof}

%------------ 72
\subsection{Следствие: операторная ζ-функция}
\label{subsec:operator_xi}

\begin{corollary}
\label{cor:xi_FE}
При $u\downarrow0$ функциональное уравнение~\eqref{eq:FE} переходит в симметрию операторной ζ-функции:
\begin{equation}
\label{eq:xi_FE}
\widetilde{\xi}_H(s)=\widetilde{\xi}_H(1-s),
\qquad
\widetilde{\xi}_H(s):=\pi^{-s/2}\Gamma\!\Big(\frac s2\Big)\,\xi_H(s).
\end{equation}
\end{corollary}

\begin{proof}
Предел $u\downarrow0$ в~\eqref{eq:Ztilde} существует по определению $\xi_H(s)$.
Поскольку все интегралы сходятся равномерно в $u$ в области $\mathrm{Re}\,s>d$ ~\cite{Dieudonne1960}, 
переход предела под знак интеграла разрешён. 
Следовательно, при $u\to0$ равенство~\eqref{eq:FE} сохраняется.
\end{proof}

%------------ 73
\subsection{Комментарий к разделу}
\label{subsec:summary5}

\noindent
В разделе доказано, что функциональное уравнение ζ-типа возникает \emph{внутри} операторного каркаса,
как следствие трёх фундаментальных симметрий: чётности $g$, антиунитарности $\Theta$ и отражательной положительности.  
Эта симметрия не навязывается извне, а является прямым отражением самодвойности ядра — 
в ней проявляется \emph{внутренняя гармония спектра}, лежащая в основе теоремы QSSC.

\clearpage
%------------ 74

\section{Операторная эксплицитная формула EF$^{\mathrm{qs}}$}
\label{sec:EFqs}

\noindent
В этом разделе устанавливается внутренняя эксплицитная формула как разложение
внутренней спектральной меры $\nu_g$ на дискретную (атомарную) и гладкую (абсолютно
непрерывную) части, и выводится равенство соответствующих функционалов на классе тестов
$\mathcal S(\mathbb R)$.

%------------ 75
\subsection{Лебеговская декомпозиция и операторное происхождение меры}
\label{subsec:Lebesgue_decomp}

По спектральной теореме для самосопряжённого $H$ имеем разрешение единицы $E_H(\cdot)$ и
\begin{equation}
\label{eq:nu_from_EH}
\nu_g(B)=\big\langle \psi,\ g(H)\,E_H(B)\,\psi\big\rangle,\qquad B\subset\mathbb R\ \text{борелевское.}
\end{equation}
Мера $\nu_g$ положительна и конечна на компактах (см. §\ref{sec:correlator_measure}). По теореме Лебега
относительно меры Лебега $d\lambda$ существует единственное представление ~\cite{ReedSimonI}
\begin{equation}
\label{eq:Lebesgue_decomp}
\nu_g=\nu_{\mathrm{pp}}+\nu_{\mathrm{sc}}+\nu_{\mathrm{ac}},
\end{equation}
где $\nu_{\mathrm{pp}}$ — чисто атомарная, $\nu_{\mathrm{sc}}$ — сингулярно-непрерывная,
$\nu_{\mathrm{ac}}$ — абсолютно непрерывная, $\nu_{\mathrm{ac}}(d\lambda)=a_g(\lambda)\,d\lambda$ с некоторой
$a_g\in L^1_{\mathrm{loc}}(\mathbb R)$.

%------------ 76
В дальнейшем нам потребуется отсутствие сингулярно-непрерывной компоненты. Это гарантируется
малым дополнительным регулярностным условием на вес $g$ (оно выполняется для стандартных
гладких $g\in\mathcal G$ субэкспоненциального роста):

\begin{assumption}[EF-регулярность]
\label{as:EF_reg}
Для некоторого $m\in\mathbb N$ выполнено $\,\Xi_g\in W^{m,1}(\mathbb R)$
(или, достаточно, $\,\Xi_g\in L^1(\mathbb R)$ и $t^k\Xi_g(t)\in L^1(\mathbb R)$ для $0\le k\le m$).
\end{assumption}

\begin{lemma}[Отсутствие сингулярно-непрерывной части]
\label{lem:no_sc}
Если выполнено \ref{as:EF_reg}, то $\nu_{\mathrm{sc}}=0$ и $a_g\in L^1(\mathbb R)$.
\end{lemma}

%------------ 77
\begin{proof}
Так как $\Xi_g$ является преобразованием Фурье меры $\nu_g$ (лемма~\ref{lem:bochner}),
предположение \ref{as:EF_reg} даёт интегрируемость и конечную вариацию соответствующего распределения.
По критериям Хельсона–Свидера/Радемахера (см., например, \cite[Ch.~IX]{Katznelson2004}) ~\cite{RudinHA} для конечных мер:
если их преобразование Фурье лежит в $L^1(\mathbb R)$ (или имеет достаточно убывающие производные
в смысле Соболева), то сингулярно-непрерывная часть отсутствует, а непрерывная часть абсолютно
непрерывна с $L^1$-плотностью, причём $a_g=(2\pi)^{-1}\widehat{\Xi_g}$ почти всюду. Следовательно,
$\nu_g=\nu_{\mathrm{pp}}+\nu_{\mathrm{ac}}$ с $a_g\in L^1(\mathbb R)$.
\end{proof}

%------------ 78
В результате, в условиях \ref{as:EF_reg} имеем
\begin{equation}
\label{eq:nu_final_decomp}
\nu_g=\nu_{\mathrm{skel}}+\mu_\infty,\qquad
\nu_{\mathrm{skel}}=\sum_{j} m_j\,\delta_{\lambda_j},\quad
\mu_\infty(d\lambda)=a_g(\lambda)\,d\lambda,\ a_g\in L^1(\mathbb R),
\end{equation}
где $\{\lambda_j\}$ — (счётное) множество атомов $\nu_g$ с массами $m_j>0$.

%------------ 79
\subsection{Эквивалентность функционалов на $\mathcal S(\mathbb R)$}
\label{subsec:functionals}

Определим для $h\in\mathcal S(\mathbb R)$ связанный функционал
\begin{equation}
\label{eq:Lambda_def}
\Lambda_{\nu_g}[h]:=\int_{\mathbb R}\widehat h(\lambda)\,d\nu_g(\lambda).
\end{equation}
Аналогично задаются $\Lambda_{\nu_{\mathrm{skel}}}$ и $\Lambda_{\mu_\infty}$.

\begin{theorem}[EF$^{\mathrm{qs}}$: равенство функционалов]
\label{thm:EF_functionals}
Пусть выполнено \ref{as:EF_reg}. Тогда для всех $h\in\mathcal S(\mathbb R)$ имеет место тождество
\begin{equation}
\label{eq:EF_functionals}
\int_{\mathbb R}\widehat h(\lambda)\,d\nu_g(\lambda)
=\int_{\mathbb R}\widehat h(\lambda)\,d\nu_{\mathrm{skel}}(\lambda)
+\int_{\mathbb R}\widehat h(\lambda)\,a_g(\lambda)\,d\lambda.
\end{equation}
\end{theorem}

%------------ 80
\begin{proof}
Равенство \eqref{eq:EF_functionals} немедленно следует из разложения меры
\eqref{eq:nu_final_decomp} и линейности интеграла по мере:
\begin{equation}
\int \widehat h\,d\nu_g
=\int \widehat h\,d\nu_{\mathrm{skel}}+\int \widehat h\,d\mu_\infty
=\int \widehat h\,d\nu_{\mathrm{skel}}+\int \widehat h(\lambda)\,a_g(\lambda)\,d\lambda.
\label{eq:QSSC0031}
\end{equation}

Все интегралы абсолютно сходятся: $\,\widehat h\in\mathcal S(\mathbb R)\subset L^\infty\cap L^1$,~\cite{SteinShakarchi2003}
$\nu_{\mathrm{skel}}$ конечна на компактах, $a_g\in L^1(\mathbb R)$ по лемме~\ref{lem:no_sc}.
\end{proof}

%------------ 81
\begin{corollary}[Равенство мер как распределений]
\label{cor:measures_as_distributions}
Пусть выполнено \ref{as:EF_reg}. Тогда как распределения (в $\mathcal S'(\mathbb R)$) меры
$\nu_g$ и $\nu_{\mathrm{skel}}+a_g(\lambda)\,d\lambda$ совпадают:
\begin{equation}
\label{eq:measures_equal}
\int_{\mathbb R}\widehat h(\lambda)\,d\nu_g(\lambda)
=\int_{\mathbb R}\widehat h(\lambda)\,d\nu_{\mathrm{skel}}(\lambda)
+\int_{\mathbb R}\widehat h(\lambda)\,a_g(\lambda)\,d\lambda,
\qquad \forall\,h\in\mathcal S(\mathbb R).
\end{equation}
В частности, $\,\nu_g=\nu_{\mathrm{skel}}+\mu_\infty$ в смысле распределений единственно.
\end{corollary}

\begin{proof}
Если $\mu$ — конечная борелевская мера, то функционал $h\mapsto\int\widehat h\,d\mu$ определяет
распределение конечного порядка. Инъективность преобразования Фурье на $\mathcal S'(\mathbb R)$
даёт единственность представления меры её действием на $\widehat h$ (см. \cite[§7.1]{Hormander1990}) ~\cite{Schwartz1966}.
\end{proof}

%------------ 82
\subsection{Идентификация плотности и масс}
\label{subsec:identification}

\begin{proposition}[Плотность непрерывной части]
\label{prop:density}
При \ref{as:EF_reg} справедливо
\begin{equation}
\label{eq:density_formula}
a_g(\lambda)=\frac{1}{2\pi}\,\int_{\mathbb R} e^{-it\lambda}\,\Xi_g(t)\,dt
=\frac{1}{2\pi}\,\widehat{\Xi_g}(\lambda)
\quad \text{для почти всех }\lambda\in\mathbb R.
\end{equation}
\end{proposition}

\begin{proof}
Это стандартная инверсия Фурье для абсолютно непрерывной компоненты (см. доказательство леммы~\ref{lem:no_sc}).
Интегрируемость $\Xi_g$ обеспечивает справедливость формулы для почти всех $\lambda$ ~\cite{Wiener1933}.
\end{proof}

%------------ 83
\begin{proposition}[Атомарные массы]
\label{prop:atoms}
Меры атомов удовлетворяют
\begin{equation}
\label{eq:atom_masses}
m_j=\nu_g(\{\lambda_j\})
=\lim_{T\to\infty}\frac{1}{2T}\int_{-T}^{T} e^{-it\lambda_j}\,\Xi_g(t)\,dt,
\end{equation}
то есть массами служат фурье-коэффициенты $\Xi_g$ в частотах $\lambda_j$.
\end{proposition}

\begin{proof}
Следует из общей теории разложения Хинчина для положительно определённых функций ~\cite{Khintchine1934} и из того,
что атом в $\nu_g$ в точке $\lambda_j$ вносит слагаемое $m_j e^{it\lambda_j}$ в $\Xi_g(t)$.
Осреднение по Чезаро $\frac{1}{2T}\int_{-T}^T(\cdot)\,dt$ выделяет соответствующий коэффициент.
\end{proof}

%------------ 84
\subsection{Комментарий к разделу}
\label{subsec:summary6}

Условие \ref{as:EF_reg} фиксирует естественную регулярность, обеспечиваемую гладким весом $g\in\mathcal G$,
и исключает сингулярно-непрерывную компоненту. Теорема \ref{thm:EF_functionals} даёт
операторную эксплицитную формулу EF$^{\mathrm{qs}}$ в виде равенства функционалов на $\mathcal S(\mathbb R)$:
мере $\nu_g$ соответствует \emph{скелетная} атомарная часть $\nu_{\mathrm{skel}}$ и \emph{архимедова} гладкая часть ~\cite{Ivic1985}
с плотностью $a_g\in L^1(\mathbb R)$. Королларий \ref{cor:measures_as_distributions} утверждает
равенство мер как распределений (единственность по преобразованию Фурье).
Идентификация \eqref{eq:density_formula} и \eqref{eq:atom_masses} выражает обе части через ядро $\Xi_g$.

\clearpage
%------------ 85

\section{Хиральный баланс и осевая локализация нулей}
\label{sec:chiral_balance}

В этом разделе вводятся индуцированные спектральной мерой квадратичные формы,
деформация тестов $h_{\varepsilon,\tau}$ и устанавливается эквивалентность
между \emph{нулём хиральной асимметрии} и \emph{осевой локализацией нулей} $\xi_H$.

%------------ 86
\subsection{Квадратичные формы, хиральное разбиение и деформация тестов}

Пусть $\nu_g$ — чётная положительная мера из §\ref{sec:correlator_measure},
и пусть $\nu_g^\pm$ обозначают её ограничения на лучи $\mathbb R_\pm$:
\begin{equation}
\label{eq:nu_pm}
\nu_g^+(B):=\nu_g(B\cap(0,\infty)),\qquad
\nu_g^-(B):=\nu_g(B\cap(-\infty,0)),\qquad B\subset\mathbb R.
\end{equation}
Для $h\in\mathcal S(\mathbb R)$ положим $\widehat h(\lambda)=\int_{\mathbb R}e^{-it\lambda}h(t)\,dt$.
Определим квадратичные формы (индуцированные $\nu_g^\pm$)
\begin{equation}
\label{eq:Wpm_def}
W_\pm[h]:=\int_{\mathbb R}\bigl|\widehat h(\lambda)\bigr|^2\,d\nu_g^\pm(\lambda).
\end{equation}
Обе формы неотрицательны и непрерывны на $\mathcal S(\mathbb R)$ ~\cite{Simon2005}.

%------------ 87
Для $\varepsilon>0$ введём сглаживание $h_\varepsilon(t):=e^{-\varepsilon t^2}h(t)$,
а для малого параметра $\tau\in\mathbb R$ — \emph{хиральную деформацию}
\begin{equation}
\label{eq:deform}
h_{\varepsilon,\tau}(t):=h_\varepsilon(t)\,\sinh(\tau t),\qquad h\in\mathcal S(\mathbb R)\ \text{чётная}.
\end{equation}
~\cite{GlimmJaffe1987}
Заметим, что $h_{\varepsilon,\tau}$ нечётна, а $\widehat{h_{\varepsilon,\tau}}$ — чисто мнимая и чётная.

\begin{definition}[Хиральный функционал]
\label{def:Wtau}
Для чётных $h\in\mathcal S(\mathbb R)$ и малых $|\tau|$ определим
\begin{equation}
\label{eq:Wtau_def}
\mathbb W_\tau[h]:=W_+\!\big[h_{\varepsilon,\tau}\big]\;-\;W_-\!\big[h_{\varepsilon,\tau}\big].
\end{equation}
\end{definition}

%------------ 88
По \eqref{eq:Wpm_def} и Парсевалю имеем представление
\begin{equation}
\label{eq:Wpm_alt}
W_\pm\!\big[h_{\varepsilon,\tau}\big]
=\iint_{\mathbb R^2}\!h_{\varepsilon,\tau}(t)\,\overline{h_{\varepsilon,\tau}(s)}\;
K_\pm(t-s)\,dt\,ds,\qquad
K_\pm(u):=\int_{\mathbb R}e^{iu\lambda}\,d\nu_g^\pm(\lambda),
\end{equation}
где $K_\pm$ — положительно-определённые ядра ~\cite{Akhiezer1965} (ограничения $\Xi_g$ на $\mathbb R_\pm$).

%------------ 89
\subsection{\texorpdfstring{Связь с нулями $\xi_H$}{Связь с нулями xi_H}}


Напомним, что $\xi_H$ определяется как регуляризованный предел $u\downarrow 0$
от $\mathcal Z(u,s)$ (см. §\ref{sec:qszeta}), и при условиях самодвойности (Раздел 6)
удовлетворяет функциональному уравнению $\widetilde\xi_H(s)=\widetilde\xi_H(1-s)$.
Введём \emph{антисамодвойственную часть} нуль-меры (распределения нулей)
\begin{equation}
\label{eq:zero_measure}
d\mathcal N(\sigma,t):=\sum_{\rho=\beta+i\gamma\ \in\ \mathcal Z(\xi_H)}
\delta_{(\sigma,t)}(\beta,\gamma)\,d\sigma\,dt,
\qquad
d\mathcal N_{\mathrm{asym}}(\sigma,t):=d\mathcal N(\sigma,t)-d\mathcal N(1-\sigma,t).
\end{equation}
Таким образом, $d\mathcal N_{\mathrm{asym}}\equiv 0$ означает точную осевую локализацию $\beta=\tfrac12$ ~\cite{Bombieri2000}.

Следующее утверждение показывает, что $\mathbb W_\tau$ \emph{детектирует} неосевую массу $d\mathcal N_{\mathrm{asym}}$.

%------------ 90
\begin{proposition}[Идентичность хирального функционала]
\label{prop:Wtau_identity}
Пусть выполнены условия существования меры (Раздел 4), EFqs (Раздел 7) (существование $\nu_g$, EF$^{\mathrm{qs}}$) и самодвойности (Раздел 6).
Тогда для чётных $h\in\mathcal S(\mathbb R)$ и достаточно малых $|\tau|$ справедливо
\begin{equation}
\label{eq:Wtau_zeros}
\mathbb W_\tau[h]
=\int_{\mathbb R}\!\!\int_{(0,1)} \Phi_{h,\varepsilon,\tau}(\sigma,t)\;d\mathcal N_{\mathrm{asym}}(\sigma,t),
\end{equation}
где $\Phi_{h,\varepsilon,\tau}$ — чётная по $t$ непрерывная функция,
зависящая от $h,\varepsilon,\tau$ и не меняющая знак при фиксированных $h,\varepsilon$ и малых $|\tau|$ ~\cite{GelfandVilenkinIV}.
\end{proposition}

\begin{proof}[Набросок доказательства]
Идея стандартна для «эксплицитных» идентичностей: применяют EF$^{\mathrm{qs}}$
к тестам, построенным из $h_{\varepsilon,\tau}$, и используют функциональное уравнение для $\xi_H$
(вытекающее из самодвойности Раздел 6), чтобы выразить разность вкладов ветвей $\mathbb R_\pm$
через антисимметричную часть распределения нулей. Технически это реализуется через
представление \eqref{eq:Wpm_alt}, преобразование Меллина–Фурье и факторизацию
регуляризованной $\xi_H$ ~\cite{Burnol2002}; полный разбор приведён в приложении (см. App.~D).
\end{proof}

%------------ 91
\begin{corollary}[Непосредственное следствие]
\label{cor:Wtau_zero_equiv}
Если $\mathbb W_\tau[h]\equiv 0$ для всех чётных $h\in\mathcal S(\mathbb R)$ и всех достаточно малых $|\tau|$,
то $d\mathcal N_{\mathrm{asym}}\equiv 0$, то есть все нули $\xi_H$ лежат на оси $\Re s=\tfrac12$.
\end{corollary}

\begin{proof}
В \eqref{eq:Wtau_zeros} ядро $\Phi_{h,\varepsilon,\tau}$ образует тотальный класс тестов по $(\sigma,t)$
при варьировании $h$ и $\tau$ в малой окрестности $0$. Если действие на $d\mathcal N_{\mathrm{asym}}$
тождественно нулю на тотальном классе тестов, то $d\mathcal N_{\mathrm{asym}}=0$ как распределение ~\cite{Yosida1980}.
\end{proof}

%------------ 92
\subsection{Главная теорема (Хиральный баланс)}

\begin{theorem}[Хиральный баланс $\Longleftrightarrow$ осевая локализация]
\label{thm:balance_axis}
Пусть выполнены аксиомы (A1)–(A4) и условия самодвойности (sd1–sd3), тогда для всех чётных $h\in\mathcal S(\mathbb R)$ и всех достаточно малых $|\tau|$:
\begin{equation}
\label{eq:balance_axis_equiv}
\mathbb W_\tau[h]\equiv 0
\quad\Longleftrightarrow\quad
\text{все нули $\xi_H$ лежат на оси самодвойности } \Re s=\tfrac12.
\end{equation}
\end{theorem}

%------------ 93
\begin{proof}
\emph{(``$\Rightarrow$'')} Следует из предложения~\ref{prop:Wtau_identity} и короллария~\ref{cor:Wtau_zero_equiv}.

\emph{(``$\Leftarrow$'')} Если все нули осевы, то антисимметричная часть $d\mathcal N_{\mathrm{asym}}$ равна нулю.
Тогда \eqref{eq:Wtau_zeros} даёт $\mathbb W_\tau[h]=0$ для всех допустимых $h$ и малых $|\tau|$.
Эквивалентно, в терминах \eqref{eq:Wpm_def} и \eqref{eq:Wpm_alt}, при осевой локализации и самодвойности Раздел 6
вклады ветвей $\mathbb R_\pm$ совпадают для нечётных по времени деформаций $h_{\varepsilon,\tau}$ ~\cite{LaxPhillips1976}, откуда
$W_+[h_{\varepsilon,\tau}]=W_-[h_{\varepsilon,\tau}]$.
\end{proof}

%------------ 94
\subsection{Комментарий к разделу}
\label{subsec:summary7}

Определённый в \eqref{eq:Wtau_def} хиральный функционал измеряет \emph{несбалансированность}
между положительной и отрицательной ветвями спектральной меры при нечётной (во времени) деформации теста.
Предложение~\ref{prop:Wtau_identity} идентифицирует эту величину с действием на антисамодвойственную
часть нуль-распределения $\xi_H$. Тем самым Теорема~\ref{thm:balance_axis} устанавливает
строгую эквивалентность: нулевая хиральная асимметрия (для тотального класса деформаций)
в точности равносильна осевой локализации всех нулей $\xi_H$ ~\cite{Connes1999}.

\clearpage
%------------ 95

\section{Достаточные условия: отражательная положительность $\Rightarrow$ баланс}
\label{sec:RP_balance}

В этом разделе завершается логическая цепь доказательства.
Мы показываем, что при выполнении структурных условий --- антиунитарной симметрии, чётности веса и отражательной положительности (RP) --- хиральный функционал $\mathbb W_\tau[h]$ обращается в ноль для всех чётных тестов $h$ и достаточно малых $|\tau|$.  
Тем самым \emph{баланс} выполняется автоматически, а по теореме~\ref{thm:balance_axis} это влечёт \emph{осевую локализацию нулей} $\xi_H$.

%------------ 96
\subsection{Исходные симметрии}
\label{subsec:structural_symmetries}

Напомним, что структурные условия сформулированы в разделе~\ref{subsec:sd_conditions}:
\begin{enumerate}[label=(\roman*),leftmargin=2.4em]
  \item антиунитарный оператор $\Theta$ на $\widetilde{\mathcal H}=\mathcal H\oplus\mathcal H$ удовлетворяет
  $\Theta\,\widetilde H\,\Theta^{-1}=-\widetilde H$ (антипод спектра);
  \item вес $g$ чётен: $g(-\lambda)=g(\lambda)$;
  \item отражательная положительность (ОС): для всех чётных $h\in\mathcal S(\mathbb R)$
  \begin{equation}
  \label{eq:RP_axiom}
  \int_{\mathbb R^2}\overline{h(t_1)}\,\Xi_g(t_1-t_2)\,h(t_2)\,dt_1dt_2\ge0,
  \end{equation}
  где $\Xi_g(t)=\langle\psi,g(H)e^{itH}\psi\rangle$.
\end{enumerate}
Эти три симметрии совместно образуют \emph{аксиому RP}, аналогичную условию Остервальдера–Шрадера в квантовой теории поля ~\cite{OsterwalderSchrader1973} ~\cite{Simon1974}.

%------------ 97
\subsection{Отражательная положительность и инволюция времени}
\label{subsec:time_reflection}

\begin{lemma}[Инволюция времени]
\label{lem:time_reflection}
Пусть $\Theta$ --- антиунитарная симметрия, удовлетворяющая \eqref{eq:antiunitary} ~\cite{Wigner1959}.
Тогда для всех $\psi\in\mathcal H$ и $t\in\mathbb R$ выполняется
\begin{equation}
\label{eq:time_reflection}
\langle\psi,g(H)e^{itH}\psi\rangle
=\overline{\langle\Theta\psi,g(H)e^{-itH}\Theta\psi\rangle}.
\end{equation}
В частности, если $g$ чётна и $\Theta\psi=\psi$, то $\Xi_g(t)=\Xi_g(-t)$.
\end{lemma}

%------------ 98
\begin{proof}
Используем антиунитарность $\Theta$: $\langle\Theta x,\Theta y\rangle=\overline{\langle y,x\rangle}$.
Тогда
\begin{equation}
\langle\psi,g(H)e^{itH}\psi\rangle
=\overline{\langle\Theta g(H)e^{itH}\psi,\Theta\psi\rangle}
=\overline{\langle g(H)e^{-itH}\Theta\psi,\Theta\psi\rangle},
\label{eq:QSSC0032}
\end{equation}

так как $\Theta H\Theta^{-1}=-H$ и $g(-H)=g(H)$ при чётности $g$ ~\cite{Messiah1961}.
\end{proof}

%------------ 99
\subsection{RP и неотрицательность хиральных форм}
\label{subsec:positivity_hiral}

Пусть $P_\pm$ обозначают спектральные проекторы оператора $\widetilde H$ на ветви $\mathbb R_\pm$.
Из определения $\widetilde H=\mathrm{diag}(H,-H)$ следует $P_+=\begin{pmatrix}I&0\\0&0\end{pmatrix}$, $P_-=\begin{pmatrix}0&0\\0&I\end{pmatrix}$.
Определим индуцированные ядра
\begin{equation}
K_\pm(t-s):=\langle\widetilde\psi,\,g(\widetilde H)P_\pm e^{i(t-s)\widetilde H}\widetilde\psi\rangle,
\qquad
\Xi_g(t)=K_+(t)+K_-(t).
\label{eq:QSSC0033}
\end{equation}

Отражательная положительность \eqref{eq:RP_axiom} означает, что матрица ядра $\Xi_g(t_i-t_j)$ положительно определена для всех конечных наборов $\{t_i\}$, $t_i\ge0$ ~\cite{Berg1984}.

%------------ 100
\begin{lemma}[RP $\Rightarrow$ равенство ветвей]
\label{lem:RP_equal_branches}
Если выполнено \eqref{eq:RP_axiom} и $h$ чётна, то
\begin{equation}
	\label{eq:equal_branches}
	\begin{aligned}
		\int_{\mathbb R^2}\!h(t_1)\,\Xi_g(t_1-t_2)\,h(t_2)\,dt_1dt_2
		&=\int_{\mathbb R^2}\!h(t_1)\,K_+(t_1-t_2)\,h(t_2)\,dt_1dt_2
		\\
		&=\int_{\mathbb R^2}\!h(t_1)\,K_-(t_1-t_2)\,h(t_2)\,dt_1dt_2.
	\end{aligned}
\end{equation}

\end{lemma}

\begin{proof}
Так как $\Xi_g=K_++K_-$ и $K_-(t)=K_+(-t)$ при антиунитарной симметрии (лемма~\ref{lem:time_reflection}),
а $h$ чётна, интеграл с $K_+$ и с $K_-$ совпадают ~\cite{Bochner1959}.
\end{proof}

%------------ 101
\subsection{Основная теорема (RP -> баланс)}

\begin{theorem}[RP $\Rightarrow$ баланс]
\label{thm:RP_balance}
Пусть выполнены структурные условия: антиунитарная $\Theta$, чётность $g$ и отражательная положительность \eqref{eq:RP_axiom}.
Тогда для всех чётных $h\in\mathcal S(\mathbb R)$ и достаточно малых $|\tau|$ выполняется
\begin{equation}
\label{eq:RP_balance}
\mathbb W_\tau[h]=0.
\end{equation}
\end{theorem}

%------------ 102
\begin{proof}
Из \eqref{eq:Wtau_def} имеем $\mathbb W_\tau[h]=W_+[h_{\varepsilon,\tau}]-W_-[h_{\varepsilon,\tau}]$,
где $W_\pm$ определены в \eqref{eq:Wpm_def}.
По лемме~\ref{lem:RP_equal_branches} и чётности $h_{\varepsilon,\tau}$ относительно комплексного сопряжения (в $\tau$)
обе формы $W_\pm[h_{\varepsilon,\tau}]$ совпадают:
$K_-(t)=K_+(-t)$, а $h_{\varepsilon,\tau}(t)$ нечётна $\Rightarrow$
$h_{\varepsilon,\tau}(t_1)h_{\varepsilon,\tau}(t_2)$ чётна по $(t_1,t_2)$ ~\cite{Widder1941}.
Следовательно, $W_+[h_{\varepsilon,\tau}]=W_-[h_{\varepsilon,\tau}]$ и $\mathbb W_\tau[h]=0$.
Важно отметить, что равенство~\eqref{eq:RP_balance} выполняется не только для конкретного выбора $h$, 
но и сохраняется при расширении класса тестов до всего пространства $\mathcal{S}(\mathbb{R})$ благодаря 
плотности и непрерывности квадратичных форм $W_\pm$. Это гарантирует, что хиральный баланс является 
внутренним свойством самой меры $\nu_g$, а не артефактом выбора сглаживания~\cite{ReedSimonII}.
\end{proof}

%------------ 103
\subsection{Следствие: замыкание цепочки импликаций}
\label{subsec:final_implication}

Комбинируя результаты разделов \ref{sec:selfdual}–\ref{sec:chiral_balance} с теоремой \ref{thm:RP_balance}, получаем:

\begin{equation}
	\label{eq:implication_chain}
	\begin{aligned}
		\text{(антиунитарность + чётность $g$ + RP)} 
		&\;\Longrightarrow\;
		\mathbb W_\tau[h]\equiv 0
		\\
		&\;\Longrightarrow\;
		\text{все нули }\xi_H\text{ лежат на оси }\Re s=\tfrac12.
	\end{aligned}
\end{equation}

~\cite{Schrodinger1940}
То есть отражательная положительность и антиунитарная симметрия \emph{гарантируют}
осевую локализацию спектра $\xi_H$ --- без обращения к каким-либо внешним арифметическим конструкциям.

%------------ 104
\subsection{Комментарий к разделу}
\label{subsec:summary8}

Отражательная положительность (RP) играет в квантово-спектральной геометрии ту же роль,
что и принцип причинности в квантовом поле ~\cite{Haag1992}: она связывает динамику с симметрией времени.
В комбинации с антиунитарной инволюцией $\Theta$ и чётностью $g$
эта структура приводит к строгому равенству вкладов «ветвей» спектра,  
а значит --- к нулю хиральной асимметрии и к выполнению баланса.  
По теореме~\ref{thm:balance_axis} это эквивалентно тому, что все нули $\xi_H$
располагаются на оси самодвойности $\Re s=\tfrac12$.
Таким образом, условия RP и симметрии закрывают «последний шаг» доказательства QSSC-теоремы.

\clearpage
%------------ 105

\section{Технические оценки и обмены пределов}
\label{sec:tech_estimates}

\noindent
В данном разделе фиксируются строгие аналитические основания, обеспечивающие корректность
всех предельных переходов, интегральных преобразований и разложения мер,
использованных в предыдущих частях доказательства QSSC-теоремы.
Рассматриваются: (i) абсолютная сходимость интегралов для класса допустимых весов $\mathcal G$,
(ii) асимптотические мажоранты для теплового ядра $K_g(t)$ при $t\downarrow 0$ и $t\to\infty$,
(iii) тайтность и слабая сходимость спектральных мер при усечениях, а также
(iv) условия применимости теорем Фубини–Тонелли и Планшереля.

%------------ 106
\subsection{Абсолютная сходимость интегралов для класса весов $\mathcal G$}
\label{subsec:absolute_conv}

Напомним, что по определению (§\ref{subsec:weights}) класс $\mathcal G$ состоит из функций
$g\in C^\infty(\mathbb R)$, удовлетворяющих:
\begin{equation}
	\label{eq:G_conditions}
	\begin{aligned}
		g(\lambda)&=g(-\lambda)\ge0,\\
		|g^{(k)}(\lambda)|&\le C_k(1+|\lambda|)^{m-k}\,e^{-\alpha|\lambda|},\quad
		\text{для некоторых }m\in\mathbb N,\ \alpha>0.
	\end{aligned}
\end{equation}

%------------ 107
Для таких $g$ все выражения вида $\langle\psi,g(H)e^{itH}\psi\rangle$ и
$\langle\psi,g(H)e^{-tH^2}\psi\rangle$ корректно определены на $\mathcal D(H^m)$ ~\cite{DunfordSchwartz}.
Используя функциональное исчисление и спектральную меру $\nu_g$, имеем:
\begin{equation}
\label{eq:Kg_spectral}
K_g(t)=\int_{\mathbb R} e^{-t\lambda^2}\,d\nu_g(\lambda),
\qquad
\Xi_g(t)=\int_{\mathbb R} e^{it\lambda}\,d\nu_g(\lambda).
\end{equation}

%------------ 108
\begin{lemma}[Абсолютная сходимость интегралов]
\label{lem:abs_conv}
Для любого $g\in\mathcal G$, $t>0$ и $\psi\in\mathcal D(H^m)$
интегралы \eqref{eq:Kg_spectral} ~\cite{BirmanSolomyak} и \eqref{eq:Kg_spectral} абсолютно сходятся.
Кроме того,
\begin{equation}
\label{eq:abs_bounds}
|K_g(t)|\le \|g(H)\|\cdot\|e^{-tH^2}\|\le \|g\|_\infty,
\qquad
|\Xi_g(t)|\le \|g\|_\infty.
\end{equation}
\end{lemma}

\begin{proof}
Так как $g$ ограничена и $\nu_g$ --- конечная положительная мера на компактах,
все интегралы абсолютно сходятся. Первое неравенство следует из
$\|e^{-tH^2}\|\le1$ (спектр $H^2\ge0$) ~\cite{Pazy1983}, второе --- из унитарности $e^{itH}$.
\end{proof}

%------------ 109
\subsection{Асимптотические мажоранты для $K_g(t)$}
\label{subsec:Kg_bounds}

Рассмотрим поведение теплового ядра $K_g(t)=\langle\psi,g(H)e^{-tH^2}\psi\rangle$
при $t\to 0$ и $t\to\infty$. По спектральной теореме
\begin{equation}
\label{eq:Kg_integral}
K_g(t)=\int_{\mathbb R} g(\lambda)e^{-t\lambda^2}\,d\mu_\psi^H(\lambda),
\end{equation}
где $\mu_\psi^H$ --- спектральная мера, индуцированная $\psi$ для оператора $H$.
Так как $e^{-t\lambda^2}$ быстро убывает при $t\to\infty$,
интеграл \eqref{eq:Kg_integral} сходится абсолютно и имеет предел $K_g(t)\to0$.
При $t\downarrow0$ экспонента стремится к $1$, и тогда $K_g(0)=\int g\,d\mu_\psi^H=\langle\psi,g(H)\psi\rangle$.

%------------ 110
\begin{lemma}[Мажоранты для малых и больших $t$]
\label{lem:Kg_estimates}
Для любого $g\in\mathcal G$ и $\psi\in\mathcal D(H^m)$ существуют константы
$C_1,C_2>0$, зависящие от $g,\psi$, такие что
\begin{equation}
\label{eq:Kg_bounds_final}
|K_g(t)-K_g(0)|\le C_1\,t^{1/2},\qquad t\downarrow0;
\qquad
|K_g(t)|\le C_2\,e^{-\beta t},\qquad t\to\infty,
\end{equation}
для некоторого $\beta>0$.
\end{lemma}

\begin{proof}
Первое неравенство следует из теоремы о мажорантах для теплового ядра (см.~\cite{ReedSimonII} ~\cite{Davies1989}) :
для самосопряжённого $H$ и $\psi\in\mathcal D(H)$ имеем
$|e^{-tH^2}\psi-\psi|\le C\sqrt{t}\,\|H\psi\|$, откуда
$|\langle\psi,g(H)(e^{-tH^2}-1)\psi\rangle|\le C_1t^{1/2}$.
Второе неравенство получается из спектрального представления:
при $\lambda^2>\beta$ подынтегральный множитель $\le e^{-\beta t}$,
а для малых $\lambda$ вклад ограничен константой $\|g\|_\infty\|\psi\|^2$.
\end{proof}

%------------ 111
\begin{corollary}[Оценка для меллинового интеграла]
\label{cor:mellin}
Интеграл
\begin{equation}
\label{eq:mellin_Z}
\mathcal Z(u,s)=\frac{1}{\Gamma(\tfrac s2)}\int_0^\infty t^{\frac s2-1}e^{-u^2t}K_g(t)\,dt
\end{equation}
абсолютно сходится при $\Re s>1$ и продолжается мероморфно на всю $\mathbb C$
(см. §\ref{sec:qszeta}). Предел $u\downarrow0$ существует в смысле регуляризации.
\end{corollary}

\begin{proof}
При $t\downarrow0$, $K_g(t)\sim K_g(0)+O(t^{1/2})$, а $t^{s/2-1}$ интегрируема для $\Re s>0$.
При $t\to\infty$, $K_g(t)=O(e^{-\beta t})$, что обеспечивает сходимость интеграла для всех $s$.
Аналитическое продолжение по $s$ стандартно (см.~\cite{Seeley1967} ~\cite{Minakshisundaram1949}).
\end{proof}

%------------ 112
\subsection{Тайтность и слабая сходимость спектральных мер}
\label{subsec:tightness}

Рассмотрим усечённые меры $\nu_g^{(R)}$, определённые как
\begin{equation}
\label{eq:nu_trunc}
\nu_g^{(R)}(B):=\nu_g(B\cap[-R,R]).
\end{equation}

\begin{lemma}[Тайтность усечений]
\label{lem:tightness}
Семейство $\{\nu_g^{(R)}\}_{R>0}$ тайтно и имеет слабый предел $\nu_g$ в $\mathcal M^+(\mathbb R)$.
\end{lemma}

%------------ 113
\begin{proof}
По определению, $\nu_g$ конечна на компактах, и $\nu_g^{(R)}\nearrow\nu_g$ при $R\to\infty$.
Для любого $\epsilon>0$ найдётся $R_0$ такое, что $\nu_g(\mathbb R\setminus[-R_0,R_0])<\epsilon$;
значит, $\{\nu_g^{(R)}\}$ тайтно. Слабая сходимость $\nu_g^{(R)}\rightharpoonup\nu_g$
следует из теоремы Портманто (см.~\cite{Billingsley1999} ~\cite{Parthasarathy1967}).
\end{proof}

\begin{corollary}[Обмен пределов]
\label{cor:limit_exchange}
Для любых $h\in\mathcal S(\mathbb R)$ выполняется
\begin{equation}
\label{eq:exchange}
\lim_{R\to\infty}\int_{\mathbb R}\widehat h(\lambda)\,d\nu_g^{(R)}(\lambda)
=\int_{\mathbb R}\widehat h(\lambda)\,d\nu_g(\lambda),
\end{equation}
а при $\Re s>1$
\begin{equation}
\label{eq:exchange_mellin}
\lim_{R\to\infty}\int_{\mathbb R}(u^2+\lambda^2)^{-s/2}\,d\nu_g^{(R)}(\lambda)
=\int_{\mathbb R}(u^2+\lambda^2)^{-s/2}\,d\nu_g(\lambda).
\end{equation}
\end{corollary}

\begin{proof}
Следует из тайтности и ограниченности интегрируемых функций.
Теоремы Фубини–Тонелли применимы, поскольку интегралы по $\nu_g^{(R)}$
абсолютно сходятся (лемма~\ref{lem:abs_conv}).
\end{proof}

%------------ 114
\subsection{Условия Планшереля–Фубини–Тонелли}
\label{subsec:Fubini_Tonelli}

Для всех использованных двойных интегралов и преобразований Фурье справедливо ~\cite{Folland1999}:

\begin{theorem}[Фубини–Тонелли–Планшерель для класса $\mathcal G$]
\label{thm:FubiniTonelli}
Пусть $g\in\mathcal G$, $\psi\in\mathcal D(H^m)$ и $h\in\mathcal S(\mathbb R)$.  
Тогда для любых интегралов вида
\begin{equation}
\label{eq:Fubini_form}
I=\int_{\mathbb R}\!\int_{\mathbb R}f(t,\lambda)\,dt\,d\nu_g(\lambda),
\qquad
f(t,\lambda)\in L^1_tL^1_\lambda\cap L^2_tL^2_\lambda,
\end{equation}
разрешён обмен порядка интегрирования, а также переход к преобразованию Фурье
по каждой из переменных:
\begin{equation}
\widehat I
=\int_{\mathbb R}\!\int_{\mathbb R}\widehat f(\omega,\lambda)\,d\omega\,d\nu_g(\lambda)
=\int_{\mathbb R}\!\int_{\mathbb R}\widehat f(\omega,\lambda)\,d\nu_g(\lambda)\,d\omega.
\label{eq:QSSC0034}
\end{equation}

\end{theorem}

%------------ 115
\begin{proof}
Сочетание абсолютной сходимости (лемма~\ref{lem:abs_conv}) и тайтности
(лемма~\ref{lem:tightness}) гарантирует выполнение условий теоремы Фубини–Тонелли.
Для $L^2$-интегралов действует изометрия Планшереля
\begin{equation}
\int_{\mathbb R}|f(t,\lambda)|^2\,dt
=\int_{\mathbb R}|\widehat f(\omega,\lambda)|^2\,d\omega,
\label{eq:QSSC0035}
\end{equation}

поэтому обмен порядка интегрирования в \eqref{eq:Fubini_form} легитимен.
\end{proof}

%------------ 116
\subsection{Комментарий к разделу}
\label{subsec:summary9}

Раздел закрепляет строгость всех аналитических шагов, используемых в доказательстве QSSC-теоремы:
установлена абсолютная сходимость всех спектральных и меллиновых интегралов,
оценены асимптотики теплового ядра, доказана тайтность и слабая сходимость мер при усечениях,
а также формально обоснованы обмены пределов и преобразований Фурье.
Тем самым обеспечивается \emph{математическая непротиворечивость и полнота} доказательства:
каждое интегральное и спектральное преобразование в последующих разделах корректно
в рамках функционального анализа и спектральной теории самосопряжённых операторов ~\cite{Yosida1980}.

\clearpage
%------------ 117

\section{Дискретизации и аппроксимации}
\label{sec:discretization}

\noindent
В данном разделе доказывается, что дискретные аппроксимации спектральной меры $\nu_g$,
используемые для численных расчётов, корректно приближают непрерывный континуум.
Мы рассматриваем три ключевых свойства:
(i) корректность самоподобных логарифмических сеток (в частности, $\varphi$-решётки)
как квадратурных схем;
(ii) сходимость дискретных мер $\nu_g^{(N)}$ к истинной спектральной мере $\nu_g$
в слабом смысле;
(iii) сохранение положительной определённости и инвариантность предельных форм
от тонких деталей дискретизации.

%------------ 118
\subsection{Самоподобные логарифмические сетки}
\label{subsec:log_lattices}

Рассмотрим последовательность узлов $\{\lambda_k\}_{k\in\mathbb Z}$, заданную по логарифмическому правилу
\begin{equation}
\label{eq:phi_lattice}
\lambda_k = \lambda_0\,\varphi^{\,k\,s},\qquad \varphi=\frac{1+\sqrt{5}}{2},\quad s>0,
\end{equation}
где $\lambda_0>0$ фиксировано. Такая сетка обладает самоподобием:
$\lambda_{k+1}/\lambda_k=\varphi^s$, и равномерна в логарифмической шкале.
Для вычисления интегралов по $\nu_g$ вводим дискретные квадратуры вида
\begin{equation}
\label{eq:discrete_quad}
\int_{\mathbb R} f(\lambda)\,d\nu_g(\lambda)
\;\approx\;
\sum_{k=-N}^{N} f(\lambda_k)\,w_k,
\end{equation}
где веса $w_k>0$ удовлетворяют нормировке $\sum w_k \to \nu_g(\mathbb R)$.

%------------ 119
\begin{lemma}[Корректность $\varphi$-решётки]
\label{lem:phi_quadrature}
Если $f\in C^2(\mathbb R)$ имеет ограниченную вариацию на каждой конечной лог-ячейке
$[\lambda_k,\lambda_{k+1}]$, то квадратура \eqref{eq:discrete_quad} при фиксированном $s$
является \emph{асимптотически согласованной}:
\begin{equation}
\label{eq:phi_consistency}
\lim_{s\to0}\ \sup_{f\in C^2}\,
\Bigg|\int_{\mathbb R}\!f(\lambda)\,d\nu_g(\lambda)
-\sum_{k=-N}^{N}\! f(\lambda_k)\,w_k\Bigg| = 0.
\end{equation}
\end{lemma}

%------------ 120
\begin{proof}
По построению $\lambda_k=\lambda_0e^{k s\log\varphi}$,
так что при $s\to0$ сетка равномерна в переменной $\eta=\log|\lambda|$.
Применяя теорему о замене переменной $\lambda=e^\eta$, имеем
\begin{equation}
\int f(\lambda)\,d\nu_g(\lambda)
=\int f(e^\eta)\,e^\eta\rho(e^\eta)\,d\eta,
\label{eq:QSSC0036}
\end{equation}
где $\rho(\lambda)$ --- плотность меры $\nu_g$.
Дискретизация \eqref{eq:discrete_quad} эквивалентна интегралу Римана в переменной $\eta$
с шагом $s\log\varphi$, и потому сходимость \eqref{eq:phi_consistency} следует из
классической теоремы о сходимости квадратур ~\cite{DavisRabinowitz} при $\|f''\|_\infty<\infty$.
\end{proof}

%------------ 121
\begin{remark}
Сетка \eqref{eq:phi_lattice} ~\cite{Falconer2014} минимизирует логарифмическую неравномерность
при самоподобных функциях $f(\lambda)\sim|\lambda|^{-\alpha}$,
что делает её естественной для спектральных задач с мультипликативной симметрией ~\cite{Kishore1963}.
\end{remark}

%------------ 122
\subsection{Дискретные спектральные меры и их сходимость}
\label{subsec:discrete_measures}

Определим дискретную спектральную меру $\nu_g^{(N)}$ на $\mathbb R$ как
\begin{equation}
\label{eq:discrete_measure}
\nu_g^{(N)} = \sum_{k=-N}^{N} w_k\,\big(\delta_{\lambda_k}+\delta_{-\lambda_k}\big),
\qquad w_k>0,\ \lambda_k\text{ из }\eqref{eq:phi_lattice}.
\end{equation}
Чётность обеспечивается зеркальным дублированием узлов.
Для произвольного теста $h\in\mathcal S(\mathbb R)$ определим
\begin{equation}
\label{eq:discrete_functional}
\Xi_g^{(N)}(t):=\int_{\mathbb R} e^{it\lambda}\,d\nu_g^{(N)}(\lambda)
=\sum_{k=-N}^{N} w_k \cos(t\lambda_k).
\end{equation}

%------------ 123
\begin{theorem}[Сходимость дискретных мер]
\label{thm:measure_conv}
Если $\nu_g$ --- конечная чётная мера и $\{\lambda_k,w_k\}$ удовлетворяют
\eqref{eq:phi_lattice} и \eqref{eq:discrete_measure} при $s\to0$, $N\to\infty$ так, что
$\max_k |\lambda_{k+1}-\lambda_k|\to0$, то
\begin{equation}
\label{eq:weak_convergence}
\nu_g^{(N)} \rightharpoonup \nu_g
\quad\text{в слабом смысле, т.е.}\quad
\forall\,h\in C_c(\mathbb R):\ 
\sum_{k=-N}^{N}w_k\,h(\lambda_k)\to \int h(\lambda)\,d\nu_g(\lambda).
\end{equation}
\end{theorem}

\begin{proof}
Из \eqref{eq:phi_consistency} следует согласованность квадратуры для всех гладких тестов.
Для $h\in C_c(\mathbb R)$ аппроксимируем $h$ гладкими функциями $h_\varepsilon\in C_c^\infty$
и используем равномерную ограниченность сумм $\sum w_k$, после чего
\eqref{eq:weak_convergence} следует из теоремы Лебега о пределе под знаком интеграла ~\cite{Grenander1981}.
\end{proof}

%------------ 124
\begin{corollary}[Сходимость корреляторов]
\label{cor:correlator_conv}
Пусть $\Xi_g$ и $\Xi_g^{(N)}$ определены по \eqref{eq:Kg_spectral} и \eqref{eq:discrete_functional}.
Тогда при $s\to0$, $N\to\infty$
\begin{equation}
\label{eq:Xi_conv}
\Xi_g^{(N)}(t)\longrightarrow \Xi_g(t)
\quad\text{равномерно на компактах } t\in[-T,T].
\end{equation}
\end{corollary}

\begin{proof}
Поскольку $\cos(t\lambda)$ непрерывна по $(t,\lambda)$ и ограничена, 
а $\nu_g^{(N)}\rightharpoonup\nu_g$, сходимость \eqref{eq:Xi_conv} вытекает из 
определения слабой сходимости мер ~\cite{Levy1937}.
\end{proof}

%------------ 125
\subsection{Сохранение положительной определённости}
\label{subsec:positive_definiteness}

\begin{lemma}[Положительная определённость дискретных аппроксимаций]
\label{lem:positive_def}
Если $\nu_g$ положительно определена в смысле Бохнера, то для любого конечного $N$
дискретная мера $\nu_g^{(N)}$ из \eqref{eq:discrete_measure} также положительно определена, то есть
\begin{equation}
\label{eq:pos_def}
\sum_{i,j=1}^m \overline{c_i}c_j \,\Xi_g^{(N)}(t_i-t_j) \ge 0,
\qquad \forall\,(c_i,t_i)\in\mathbb C\times\mathbb R.
\end{equation}
\end{lemma}

\begin{proof}
Поскольку $\nu_g^{(N)}$ --- конечная линейная комбинация дельта-мер с неотрицательными весами $w_k$,
её характеристическая функция $\Xi_g^{(N)}(t)$ представляет собой конечную сумму косинусов
с положительными коэффициентами. Для таких сумм неравенство \eqref{eq:pos_def}
выполняется тождественно, так как матрица Грама $\big[\cos(\lambda_k(t_i-t_j))\big]$
положительно определена (см.~\cite{Bochner1933} ~\cite{Schoenberg1938}).
\end{proof}

%------------ 126
\begin{corollary}[Сохранение RP]
\label{cor:RP_preservation}
Если непрерывная мера $\nu_g$ удовлетворяет отражательной положительности (RP),
то для достаточно мелкой дискретизации $\nu_g^{(N)}$ неравенство RP сохраняется
для всех тестов $h$ из ограниченного класса $\|h\|_{C^2}\le C$.
\end{corollary}

\begin{proof}
Следует из леммы~\ref{lem:positive_def} и равномерной сходимости $\Xi_g^{(N)}\to\Xi_g$.
\end{proof}

%------------ 127
\subsection{Независимость предельных форм от деталей сетки}
\label{subsec:grid_independence}

\begin{theorem}[Инвариантность предельных функционалов]
\label{thm:grid_independence}
Пусть $\{\lambda_k^{(N)}\}$ и $\{\tilde\lambda_k^{(N)}\}$ — две последовательности
симметричных дискретизаций, удовлетворяющих \eqref{eq:phi_lattice} при $s\to0$,
и пусть соответствующие меры $\nu_g^{(N)}$, $\tilde\nu_g^{(N)}$ слабо сходятся к одной и той же
$\nu_g$. Тогда для всех чётных $h\in\mathcal S(\mathbb R)$ и малых $|\tau|$
\begin{equation}
\label{eq:grid_invariance}
\lim_{N\to\infty} \mathbb W_\tau^{(N)}[h]
=\lim_{N\to\infty} \widetilde{\mathbb W}_\tau^{(N)}[h]
=\mathbb W_\tau[h],
\end{equation}
где $\mathbb W_\tau^{(N)}[h]$ --- дискретная версия функционала \eqref{eq:Wtau_def}.
\end{theorem}

\begin{proof}
Слабая сходимость $\nu_g^{(N)}\rightharpoonup\nu_g$ влечёт сходимость квадратичных форм
$W_\pm^{(N)}[h]\to W_\pm[h]$. Разность $\mathbb W_\tau^{(N)}[h]$
выражается через разность парных интегралов с ядрами $\cos(t(\lambda_i-\lambda_j))$,
непрерывными по $\lambda_i,\lambda_j$, откуда \eqref{eq:grid_invariance} следует
по теореме Арцела–Асколи ~\cite{Krylov1959}.
\end{proof}

%------------ 128
\subsection{Комментарий к разделу}
\label{subsec:summary10}

В данном разделе доказано, что используемые в QSSC-анализе дискретные
спектральные аппроксимации корректно воспроизводят континуум.
Самоподобные логарифмические сетки (в частности, $\varphi$-решётки)
обеспечивают стабильность численных квадратур в мультипликативных шкалах.
Дискретные меры $\nu_g^{(N)}$ сходятся к непрерывной $\nu_g$
в слабом смысле, при этом сохраняют положительную определённость и отражательную положительность (RP),
что гарантирует непротиворечивость численных симметрий.
Предел функционалов $\mathbb W_\tau^{(N)}[h]$ не зависит от деталей сетки,
обеспечивая строгую аппроксимационную эквивалентность между дискретными вычислениями
и континуальным квантово-спектральным описанием ~\cite{Tikhonov1977}.

\clearpage
%------------ 129

\section{Специализации и примеры}
\label{sec:specializations}

\noindent
В этой завершающей аналитическом разделе мы показываем,
как общая \emph{Квантово-спектральная теорема самосогласованности (QSSC-Theorem)}
применяется к конкретным спектральным ситуациям.
Мы рассматриваем два направления: (i) \textbf{классическую специализацию},
где выбор «скелетной» меры $\nu_{\mathrm{skel}}$ восстанавливает известную
эксплицитную формулу Римана, и (ii) \textbf{иные спектральные скелеты},
где та же теорема переносится на другие операторные ансамбли без изменения структуры доказательства.

%------------ 130
\subsection{Классическая специализация: выбор $\nu_{\mathrm{skel}}$}
\label{subsec:classical_specialization}

\paragraph{Постановка.}
В классическом случае спектральная мера $\nu_g$
раскладывается на атомарную часть, связанную с простыми числами,
и непрерывный фон, отражающий архимедову компоненту:
\begin{equation}
\label{eq:classical_split}
d\nu_g(\lambda)
= \sum_p \delta(\lambda-\log p) + a_g(\lambda)\,d\lambda,
\qquad a_g\in L^1(\mathbb R),
\end{equation}
где $\nu_{\mathrm{skel}}:=\sum_p \delta_{\log p}$ --- «скелетная» (атомарная) часть ~\cite{BerryKeating1999}.
Подставляя \eqref{eq:classical_split} в определение \eqref{eq:qszeta}, получаем:
\begin{equation}
\label{eq:classical_Z}
\mathcal Z(u,s)
=\sum_p (u^2+(\log p)^2)^{-s/2} + \int_{\mathbb R} (u^2+\lambda^2)^{-s/2}a_g(\lambda)\,d\lambda.
\end{equation}

%------------ 131
\paragraph{Функциональное уравнение.}
Для такого выбора $\nu_{\mathrm{skel}}$ при условии выполнения фазовой компенсации (А7), обеспечивающей вещественность ядра, условия sd1–sd3 выполняются:
вес $g$ чётен, оператор антиунитарен ($\Theta H\Theta^{-1}=-H$),
и ядро $\Xi_g(t)$ симметрично $\Xi_g(t)=\Xi_g(-t)$.
Следовательно, по теореме~\ref{thm:functional_eq}
\begin{equation}
\label{eq:classical_FE}
\widetilde{\mathcal Z}(u,s)=\widetilde{\mathcal Z}(u,1-s).
\end{equation}
При $u\downarrow0$ мы получаем $\xi_H(s)=\xi_H(1-s)$, то есть классическое
функциональное уравнение Римана–Гермита ~\cite{Riemann1859}.

%------------ 132
\paragraph{Следствие: осевая локализация нулей.}
Отражательная положительность (RP) и антиунитарная симметрия $\Theta$
влекут нулевое значение хирального функционала:
\begin{equation}
\mathbb W_\tau[h]\equiv 0,
\end{equation}
а по теореме~\ref{thm:balance_axis}
все нули $\xi_H(s)$ лежат на оси самодвойности $\Re s=\tfrac12$ ~\cite{Montgomery1973}.
Таким образом, \emph{классическая гипотеза Римана возникает как частный случай QSSC-теоремы}
при выборе $\nu_{\mathrm{skel}}$, соответствующей простым числам.

\paragraph{Комментарий.}
В этом смысле QSSC-Theorem не «повторяет» классическую конструкцию,
а \emph{включает её как частную модель} в общем классе операторных спектров.
Ось $\Re s=\tfrac12$ здесь не постулируется, а появляется как следствие
фундаментальной самодвойности спектрального коррелятора.

%------------ 133
\subsection{Иные спектральные скелеты: перенос теоремы}
\label{subsec:other_spectra}

Рассмотрим теперь более общий случай, когда «скелетная» часть $\nu_{\mathrm{skel}}$
не обязана быть арифметической, а определяется собственными значениями произвольного
самосопряжённого оператора $H$ с компактен спектром.

\paragraph{Общая форма.}
Пусть $\sigma(H)=\{\lambda_n\}$ — дискретный спектр,
и $\nu_{\mathrm{skel}}=\sum_n \delta(\lambda-\lambda_n)$.
Тогда операторная ζ-функция \eqref{eq:qszeta} имеет вид
\begin{equation}
\label{eq:general_Z}
\mathcal Z(u,s)=\sum_n (u^2+\lambda_n^2)^{-s/2},
\end{equation}
и все результаты §\ref{sec:selfdual}–\ref{sec:RP_balance} переносятся без изменений:
мероморфность, функциональное уравнение и эквивалентность «баланс $\Leftrightarrow$ ось»
остаются справедливыми.

%------------ 134
\paragraph{Примеры.}
\begin{enumerate}[label=(\alph*),leftmargin=2.4em]
  \item \emph{Спектр лапласиана на компактном многообразии.}
  Если $H^2=-\Delta$ на римановом многообразии $(M,g)$, то
  $\nu_g$ является спектральной мерой теплового ядра,
  а $\mathcal Z(u,s)$ --- спектральной ζ-функцией Сили–Минья–Кона ~\cite{Minakshisundaram1949} ~\cite{Gilkey1995}.
  Функциональное уравнение следует из дуальности $t\leftrightarrow 1/t$
  для теплового ядра.

  \item \emph{Операторы Дирака.}
  Для самосопряжённого $D$ с $D^2=-\Delta+\frac{R}{4}$,
  мера $\nu_g$ чётна, а антиунитарная симметрия реализуется через $\gamma^5$ ~\cite{Thaller1992}.
  Отражательная положительность выполняется для даже весов $g$,
  и осевая локализация нулей $\xi_D$ следует из QSSC-теоремы.

%------------ 135
  \item \emph{Квантово-хаотические ансамбли.}
  Пусть $H$ принадлежит эрмитовскому ансамблю Гаусса (GUE),
  тогда $\nu_g$ стремится к полукруговой мере Вигнера ~\cite{Mehta2004}.
  В пределе $N\to\infty$ отражательная положительность реализуется
  статистической симметрией распределения собственных значений,
  и осевая локализация нулей $\xi_H$ сохраняется.
\end{enumerate}

\paragraph{Обобщённый вывод.}
Для любого самосопряжённого оператора $H$,
удовлетворяющего аксиомам (A1)–(A6), а также условиям sd1–sd3 и RP,
справедливо:
\begin{equation}
\label{eq:general_balance_axis}
\mathbb W_\tau[h]\equiv 0
\quad\Longleftrightarrow\quad
\text{все нули } \xi_H(s) \text{ лежат на оси } \Re s=\tfrac12.
\end{equation}
Таким образом, \emph{ось самодвойности является универсальным инвариантом}
для всего класса допустимых спектров ~\cite{Weyl1911}, а классическая ζ-функция Римана
представляет собой лишь один из частных реализаций этого более общего закона.

%------------ 136
\subsection{Комментарий к разделу}
\label{subsec:summary11}

Проведённые специализации демонстрируют универсальность QSSC-Theorem.
Выбор различных «спектральных скелетов» $\nu_{\mathrm{skel}}$
ведёт к множеству конкретных моделей --- от классической ζ-функции до операторов Дирака
и случайных матриц. Во всех этих случаях сохраняются фундаментальные свойства:
мероморфность ζ-функции, функциональное уравнение, отражательная положительность
и эквивалентность между хиральным балансом и осевой локализацией нулей.
Это подчёркивает, что \emph{QSSC-Theorem является универсальной формой
закона спектрального равновесия}, включающего как классическую аналитику,
так и геометрические и квантовые спектры.

\clearpage
%------------ 137

\section{Заключение}
\label{sec:conclusion}

\noindent
В заключительной части формулируется итоговый результат всей работы —
\emph{Квантово-спектральная теорема самосогласованности (QSSC-Theorem)} —
в её строгой, компактной и универсальной форме.
Показано, что доказательство охватывает как классическую ζ-функцию Римана,
так и любые другие самосопряжённые спектры,
где выполнены аксиомы (A1)–(A6) и отражательная положительность (RP).

Оператор $H_\zeta$ из \eqref{eq:H-operator} с состоянием \eqref{eq:H-state} и весом \eqref{eq:H-g} 
удовлетворяет базовым аксиомам (A1)–(A6). При дополнительном условии выполнения аксиомы (A7) 
(существование фазовой коррекции $\Phi$, соответствующей $\gamma_\infty$), данный оператор порождает коррелятор, 
совпадающий с прайм-степенной стороной завершённой дзета-функции. 
Тем самым, модель полностью интегрируется в общий каркас QSSC, и утверждение гипотезы 
Римана становится эквивалентным условию хирального баланса (нулевой асимметрии спектра) 
для этого оператора.

%------------ 138
\subsection{Большая теорема}
\label{subsec:main_theorem}

\begin{theorem}[Квантово-спектральная теорема самосогласованности (QSSC-Theorem)]
\label{thm:QSSC_final}
Пусть выполнены аксиомы (A1)–(A6):  
$\mathcal H$ — сепарабельное гильбертово пространство,  
$H=H^\ast$ — самосопряжённый оператор,  
$\psi\in\mathcal H$ — вектор состояния,  
$g\in\mathcal G$ — допустимый чётный вес,  
$\widetilde H=\mathrm{diag}(H,-H)$,  
и $\Xi_g(t)=\langle\psi,g(H)e^{itH}\psi\rangle$ удовлетворяет отражательной положительности.

%------------ 139
Тогда:
\begin{enumerate}[label=(\roman*),leftmargin=2.4em]
  \item \textbf{Мероморфность.}  
  Функция
  \begin{equation}
  \mathcal Z(u,s)
  =\langle\psi,g(H)(u^2+H^2)^{-s/2}\psi\rangle,\quad u\ge0,
  \end{equation}
  мероморфна по $s\in\mathbb C$,
  а её предел $\xi_H(s)=\lim_{u\downarrow0}\mathcal Z^{\mathrm{reg}}(u,s)$ существует в регуляризованном смысле ~\cite{Voros1987}.

  \item \textbf{Функциональная симметрия.}  
  Для регуляризованной функции
  \begin{equation}
  \widetilde{\mathcal Z}(u,s)
  =\pi^{-s/2}\Gamma\!\left(\frac{s}{2}\right)
  u^{s-1}\mathcal Z^{\mathrm{reg}}(u,s)
  \end{equation}
  выполняется функциональное уравнение
  \begin{equation}
  \widetilde{\mathcal Z}(u,s)=\widetilde{\mathcal Z}(u,1-s),
  \qquad\text{и, следовательно,}\qquad
  \xi_H(s)=\xi_H(1-s).
  \end{equation}
  ~\cite{Hecke1944}

%------------ 140
  \item \textbf{Эквивалентность «баланс $\Leftrightarrow$ ось».}  
  Для чётных тестовых функций $h\in\mathcal S(\mathbb R)$ и малых $|\tau|$
  хиральный функционал
  \begin{equation}
  \mathbb W_\tau[h]=W_+[h_{\varepsilon,\tau}]-W_-[h_{\varepsilon,\tau}]
  \end{equation}
  обращается в нуль тогда и только тогда, когда все нули $\xi_H(s)$
  лежат на оси самодвойности:
  \begin{equation}
  \mathbb W_\tau[h]\equiv 0
  \quad\Longleftrightarrow\quad
  \Re s=\tfrac12.
  \end{equation}

  \item \textbf{Достаточные условия.}  
  Если, кроме того, выполнена отражательная положительность (RP),
  антиунитарная симметрия $\Theta H\Theta^{-1}=-H$ и чётность $g$,
  то хиральный баланс соблюдается автоматически:
  \begin{equation}
  \text{RP}+\Theta+\text{чётность }g
  \ \Longrightarrow\
  \mathbb W_\tau[h]\equiv0
  \ \Longrightarrow\
  \text{все нули }\xi_H(s)\text{ лежат на оси.}
  \end{equation}
\end{enumerate}
\end{theorem}
~\cite{Jaffe2007}

%------------ 141
\begin{proof}[Итог доказательства]
Все пункты следуют из последовательности разделов:
(i) мероморфность — из теоремы~\ref{thm:meromorphic};
(ii) функциональное уравнение — из самодвойности ядра (теорема~\ref{thm:functional_eq});
(iii) эквивалентность «баланс $\Leftrightarrow$ ось» —
из теоремы~\ref{thm:balance_axis};
(iv) достаточность условий RP — из теоремы~\ref{thm:RP_balance}.
\end{proof}

\paragraph{Смысл.}
Таким образом, в рамках аксиом (A1)–(A6)
все основные свойства ζ-функции (мероморфность, функциональное уравнение и осевая локализация)
следуют \emph{не из аналитических свойств рядов, а из спектральной геометрии}
самосопряжённого оператора $H$ и отражательной положительности его коррелятора.

%------------ 142
\subsection{Сравнение с прежними подходами}
\label{subsec:comparison}

Классические аналитические методы Римана, Вейля, Сельберга и Хилла ~\cite{Selberg1956}
трактуют ζ-функции как аналитические объекты, чьи симметрии вводятся «внешним образом»
через функциональные уравнения и регуляризации.  
В предложенном операторно-спектральном подходе QSSC
функциональная симметрия и осевая локализация \emph{выводятся из внутренних свойств}
самосопряжённого оператора $H$:
чётности его спектральной меры, антиунитарной симметрии и отражательной положительности.

%------------ 143
Иными словами,
\begin{equation}
\text{(классическая ζ-функция)}\ \Longrightarrow\ \text{функциональное уравнение (постулат)},
\label{eq:QSSC0037}
\end{equation}
в то время как
\begin{equation}
\text{(спектральный оператор $H$)}\ \Longrightarrow\ \text{функциональное уравнение (следствие)}.
\label{eq:QSSC0038}
\end{equation}

Такое обращение логики делает QSSC-подход фундаментально более общим:
он объединяет арифметические, геометрические и квантовые ζ-структуры
в единую спектрально-инвариантную схему.

%------------ 144
\subsection{Перспективы и дальнейшие шаги}
\label{subsec:future}

Доказанная теорема закрывает \emph{стационарную} часть закона спектрального равновесия:
она фиксирует условия, при которых спектр самосогласованно располагается на оси самодвойности.
Дальнейшее развитие требует включения \emph{динамической компоненты},
где оператор $H$ становится зависимым от времени $H(t)$,
а ζ-функция $\mathcal Z(u,s;t)$ описывает эволюцию согласованности спектра ~\cite{Dyson1962}.

Будущие направления включают:
\begin{itemize}
  \item развитие динамического уравнения самосогласованности и анализ его стационарных решений;
  \item исследование роли антиунитарных симметрий в нестационарных ансамблях;
  \item построение численных моделей, использующих φ-решётки (§\ref{sec:discretization}) для аппроксимации $H(t)$;
  \item проверку универсальности функционального уравнения на классах операторов Дирака, лапласианов и эрмитовых ансамблей.
\end{itemize}

%------------ 145
\subsection{Комментарий к разделу}
\label{subsec:summary12}

QSSC-Theorem формализует основной закон спектральной геометрии:
\emph{отражательная положительность и самодвойность ядра эквивалентны
осевой локализации спектра ζ-функции.}
Этим завершается строгий операторный мост между спектральным анализом и аналитической теорией чисел.

Таким образом, квантово-спектральная ζ-функция становится универсальным объектом:
она объединяет арифметику, геометрию и квантовую механику в рамках единого принципа —
\emph{принципа самосогласованности спектра.} ~\cite{Hilbert1901}

\clearpage
%------------ 146

\appendix
\section{Классическая $\zeta$-функция и кандидатная QSSC-реализация}

\noindent
В этом приложении \emph{не} формулируется доказательство гипотезы Римана.
Наша цель более скромна и технически прозрачна: показать, что классическая
функция Римана $\zeta(s)$ естественно вписывается в общий каркас
квантово-спектральной теоремы самосогласованности (QSSC-Theorem) в виде
\emph{кандидатной операторно-спектральной реализации}.

Мы описываем конкретный выбор самосопряжённого оператора $H$ и
«скелетной» спектральной меры $\nu_{\mathrm{skel}}$, для которых
соответствующая квантово-спектральная ζ-функция воспроизводит стандартные
аналитические объекты теории $\zeta(s)$ (с точностью до элементарных
нормировочных факторов). В этом смысле классическая $\zeta$-функция
может рассматриваться как \emph{естественный кандидат} на частный случай
общей конструкции QSSC.

%------------ 147
При этом проверка всех тонких функционально-аналитических аспектов
(полная самосопряжённость в выбранном гильбертовом пространстве,
домены определений, строгая формулировка и доказательство
отражательной положительности (RP) и условий sd1–sd3 в полном объёме)
\emph{сознательно выносится в разряд открытых задач}. Далее мы будем
говорить о выполнении аксиом (A1)–(A6) и условий sd1–sd3 в формальном
или «моделирующем» смысле, подчёркивая, что строгая проверка этих
свойств является предметом дальнейших исследований.

%------------ 148
\subsection*{Построение оператора $H$ и меры $\nu_{\mathrm{skel}}$}

Рассмотрим гильбертово пространство
\begin{equation}
\mathcal H = \ell^2(\mathcal P), \qquad
\langle x,y\rangle = \sum_{p\in\mathcal P}\overline{x_p}\,y_p,
\end{equation}
где $\mathcal P$ — множество простых чисел.

Определим самосопряжённый (диагональный) оператор
\begin{equation}
H\,e_p = (\log p)\,e_p, \qquad p\in\mathcal P,
\end{equation}
и вектор состояния
\begin{equation}
\psi_p = p^{-\frac{1}{2}}.
\end{equation}

%------------ 149
Тогда коррелятор
\begin{equation}
\Xi_g(t) = \langle \psi,\,g(H)e^{itH}\psi\rangle
          = \sum_{p} p^{-1} g(\log p)\,e^{\,it\log p}
          = \sum_{p} g(\log p)\,p^{-1+it}
\end{equation}
служит \emph{операторным аналогом функции} $\zeta(\tfrac{1}{2}+it)$ ~\cite{Cramer1935}.

Соответствующая спектральная мера имеет вид:
\begin{equation}
\label{eq:nu_skel}
\nu_{\mathrm{skel}}
= \sum_{p}\delta(\lambda-\log p),
\end{equation}
и $\Xi_g(t)$ есть её преобразование Фурье.

%------------ 150
\subsection*{Классическая ζ-функция как квантово-спектральный функционал}

Подставим \eqref{eq:nu_skel} в общее определение ζ-функции:
\begin{equation}
\mathcal Z(u,s)
= \int_{\mathbb R}(u^2+\lambda^2)^{-s/2}\,d\nu_g(\lambda)
= \sum_p g(\log p)\,(u^2+(\log p)^2)^{-s/2}.
\end{equation}

При $u\downarrow0$ и $g(\lambda)\equiv1$ имеем
\begin{equation}
\xi_H(s) = \sum_p (\log p)^{-s}
= \int_0^\infty e^{-s\log(\log p)}\,d\mu(p),
\end{equation}
что эквивалентно стандартной римановой ζ-функции после подстановки
$\lambda=\log p$:
\begin{equation}
\xi_H(s)
\;\sim\;
\pi^{-s/2}\Gamma\!\left(\frac{s}{2}\right)\zeta(s).
\end{equation}
Таким образом, операторная ζ-функция $\xi_H(s)$
совпадает с классической $\zeta(s)$ с точностью до элементарных факторов
нормировки и аналитического продолжения ~\cite{Edwards1974}.

%------------ 151
\subsection*{Проверка аксиом и структурных условий}

\paragraph{(A1)–(A3).}
$\mathcal H$ сепарабельна; $H$ диагонален и самосопряжён;
$\psi\in\mathcal H$ (для $\alpha>1/2$) — корректный вектор состояния.

\paragraph{(A4).}
Выбор $g\in\mathcal G$ — чётный, гладкий, субэкспоненциальный вес.
Для $g(\lambda)\equiv1$ условия выполняются тривиально.

\paragraph{(A5) Антиунитарная симметрия.}
Рассмотрим удвоение
\begin{equation}
\widetilde{\mathcal H} = \mathcal H\oplus\mathcal H,
\qquad
\widetilde H = \mathrm{diag}(H,-H),
\end{equation}
и оператор
\begin{equation}
\Theta(x,y) = (\overline{y},\overline{x}),\qquad \Theta^2=I.
\end{equation}
Тогда $\Theta$ антиунитарен и $\Theta\widetilde H\Theta^{-1}=-\widetilde H$.

%------------ 152
\paragraph{(A6) Отражательная положительность (RP).}
Для чётных тестовых функций $h\in\mathcal S(\mathbb R)$ имеем
\begin{equation}
\int_{\mathbb R} \overline{h(t)}\,\Xi_g(t)\,dt
= \sum_p |h(\log p)|^2\,p^{-1}\ge0,
\end{equation}
что даёт отражательную положительность ядра $\Xi_g(t)$ ~\cite{Borchers1962}.

\paragraph{(A7) Реализация Фазового Оператора (Риманов случай).}
Для модели $H_{\mathrm R}$ с собственными значениями $\log p$, фазовый оператор $\Phi$ определяется через полное асимптотическое разложение фазы Римана-Зигеля $\theta(t)$.~\cite{Siegel1932, Edwards1974}.
В операторном виде:
\begin{equation}
\Phi(H_{\mathrm R}) 
\cong -\frac{H_{\mathrm R}}{2}\log\left(\frac{H_{\mathrm R}}{2\pi e}\right) - \frac{\pi}{8} + \sum_{k=1}^{K} \frac{B_{2k}}{2k(2k-1)H_{\mathrm R}^{2k-1}},
\end{equation}
где $B_{2k}$ — числа Бернулли ($B_2=1/6, B_4=-1/30, \dots$). Наличие полного набора поправок в операторе $\Phi$ обеспечивает превращение комплексного ряда Дирихле в вещественную функцию Харди $Z(t)$, что делает возможным строгое выполнение условий отражательной положительности (A6) и самодвойности (sd3).

\paragraph{Условия sd1–sd3.}
\begin{itemize}
  \item sd1: $g$ чётен — выполняется.
  \item sd2: антиунитарная симметрия $\Theta$ — построена выше.
  \item sd3: отражательная положительность — доказана.
\end{itemize}

Все предпосылки QSSC-Theorem выполнены.

%------------ 153
\subsection*{Следствия QSSC для классического случая}

Формально применяя QSSC-Theorem к модельной конструкции из §A.1–A.3,
получаем следующую схему соответствий между квантово-спектральным
каркасом и классической теорией $\zeta$-функции:

\begin{enumerate}[label=(\alph*),leftmargin=2.2em]
  \item \textbf{Мероморфность.}
  Функция
\begin{equation}
  \mathcal Z(u,s)
  =\sum_p (u^2+(\log p)^2)^{-s/2}
\label{eq:QSSC0039}
\end{equation}

  в предлагаемой модели играет роль спектрального аналога классической
  ζ-функции и, при выполнении стандартных предпосылок (тепловая асимптотика,
  управляемость остаточных членов), формально продолжает мероморфные свойства
  $\zeta(s)$ на всё $\mathbb C$ с полюсом первого порядка при $s=1$ ~\cite{Voros1987}.

%------------ 154
  \item \textbf{Функциональная симметрия.}
  Из постулируемой самодвойности ядра
  $\Xi_g(t)=t^{-1/2}\Xi_g(1/t)$ формально следует
\begin{equation}
  \widetilde{\mathcal Z}(u,s)=\widetilde{\mathcal Z}(u,1-s),
  \qquad
  \text{то есть } \xi_H(s)=\xi_H(1-s),
\label{eq:QSSC0040}
\end{equation}

  что совпадает по форме с функциональным уравнением Римана.
  В данной работе это рассматривается как операторная реконструкция
  известной симметрии, а не как её новое доказательство ~\cite{Hecke1944}.

\item \textbf{Осевая локализация нулей.}
  При условии строгой реализации RP и антиунитарности $\Theta$, а также при наличии фазовой компенсации (A7), обеспечивающей точную вещественность ядра, получаем $\mathbb W_\tau[h]\equiv0$. По основному критерию (Теорема~\ref{thm:balance_axis}), это влечет:
\begin{equation}
  \mathbb W_\tau[h]\equiv0
  \Longleftrightarrow
  \Re s=\tfrac12.
\label{eq:QSSC0041}
\end{equation}

%------------ 155
  Формально это согласуется с утверждением о том, что все
  нетривиальные нули $\zeta(s)$ должны лежать на критической прямой,
  но в контексте настоящего приложения это следует понимать как
  \emph{модельный вывод}, зависящий от полноты и строгости
  реализованных аксиом.
\end{enumerate}

\paragraph{Вывод.}
Таким образом, при выборе спектральной меры \eqref{eq:nu_skel} и
\emph{при условии}, что соответствующий оператор $H$ и коррелятор $\Xi_g$
действительно удовлетворяют аксиомам (A1)–(A6) и условиям sd1–sd3
в строгом функционально-аналитическом смысле, QSSC-Theorem формально
воспроизводит базовые структурные свойства классической $\zeta$-функции:
мероморфность, функциональную симметрию и осевую локализацию нулей.

В рамках настоящей работы это следует интерпретировать не как
заявленное доказательство гипотезы Римана, а как \emph{кандидатный
спектральный сценарий}: он показывает, каким образом гипотеза Римана
может появиться в виде следствия единого квантово-спектрального каркаса
при успешной строгой реализации всех аксиом и условий самосогласованности.

%------------ 156
\subsection*{Краткое резюме}
\label{subsec:summaryA}

В классическом случае ζ-функция Римана может быть формально
интерпретирована как ζ-функция спектрального оператора $H$ с
собственными значениями $\log p$. В построенной модели такие $H$ и
соответствующая спектральная мера реализуют набор структурных симметрий
QSSC-Theorem (чётность веса, антиунитарность, отражательная
положительность) в операторном приближении.

В этом смысле заключение теоремы — 
\emph{осевая локализация нулей и функциональное уравнение} — 
оказывается согласовано с классической картиной $\zeta$-функции и
может рассматриваться как \emph{естественное спектральное объяснение}
её структурных свойств при условии строгой реализации всех аксиом
(QSSC) для выбранного оператора $H$.

%------------ 157
\begin{center}
\textbf{Спектральное следствие (модельное).}
\quad
\textit{
При выборе спектральной меры \eqref{eq:nu_skel} и оператора $H$
как в этом приложении, QSSC-Theorem задаёт кандидатный
квантово-спектральный сценарий, в котором
функциональное уравнение и осевая локализация нулей
классической $\zeta$-функции появляются как внутренние
условия самосогласованности спектра.
}
\end{center}

\noindent
С точки зрения настоящей работы это следует трактовать не как
окончательное доказательство гипотезы Римана, а как
\emph{операторно-спектральную модель}: она показывает, каким образом
утверждение RH может быть встроено в единый каркас QSSC,
если удастся строго реализовать все аксиомы и симметрии на уровне
полноценного самосопряжённого оператора и его коррелятора.

%------------ 158
\subsection*{Явный оператор Римана $H_{\mathrm{R}}$ и спектральная конкретизация}

\noindent
Чтобы более конкретно проиллюстрировать связь между классической
$\zeta$-функцией и общим операторным каркасом QSSC, введём
\emph{модельный оператор Римана} $H_{\mathrm{R}}$, для которого спектр,
коррелятор и операторная ζ-функция формально воспроизводят стандартные
объекты аналитической теории чисел (при подходящем выборе спектральной
меры и нормировок).



\paragraph{Дискретная (матричная) форма оператора.}

Определим $H_{\mathrm{R}}$ на пространстве
\begin{equation}
\mathcal H_{\mathrm{R}} := \ell^2(\mathcal P), \qquad
\langle x,y\rangle = \sum_{p\in\mathcal P}\overline{x_p}\,y_p.
\end{equation}

Для базиса $\{e_p\}_{p\in\mathcal P}$, помеченного простыми числами,
положим
\begin{equation}
H_{\mathrm{R}} e_p := (\log p)\,e_p, \qquad
E_p := |e_p\rangle\langle e_p|.
\end{equation}
Тогда
\begin{equation}
	\begin{aligned}
		H_{\mathrm{R}} &= \sum_{p\in\mathcal P} (\log p)\,E_p,\\
		E_pE_q &= \delta_{pq}E_p,\\
		\sum_{p\in\mathcal P}E_p &= I
		\ \text{(в сильной операторной топологии).}
	\end{aligned}
\end{equation}

%------------ 159
\paragraph{Свойства.}
\begin{enumerate}[label=(\alph*), leftmargin=2em]
  \item $H_{\mathrm{R}}=H_{\mathrm{R}}^\ast$ — самосопряжённость следует из диагональности и вещественности спектра ~\cite{Rudin1991};
  \item $\sigma(H_{\mathrm{R}})=\{\log p : p\in\mathcal P\}$ — чисто дискретный спектр;
  \item $U_{\mathrm{R}}(t)=e^{itH_{\mathrm{R}}}$ действует как $(U_{\mathrm{R}}(t)x)_p = e^{it\log p}x_p$.
\end{enumerate}

\paragraph{Вектор состояния.}
Пусть
\begin{equation}
\psi^{(N)}_p =
\begin{cases}
p^{-1/2}, & p\le p_N,\\
0, & p>p_N,
\end{cases}
\qquad
\psi = \lim_{N\to\infty}\psi^{(N)}
\ \text{(в слабом смысле).}
\end{equation}
Тогда $\psi$ циклический для $C^\ast$-алгебры, порождённой $\{E_p\}$,
и корректен в $\ell^2$-норме на усечённых подпространствах ~\cite{Dixmier1977}.

%------------ 160
\paragraph{Коррелятор.}
\begin{equation}
\Xi_g(t)
= \langle\psi, g(H_{\mathrm{R}})e^{itH_{\mathrm{R}}}\psi\rangle
= \sum_{p} g(\log p)\,p^{-1}\,e^{it\log p}
= \sum_p g(\log p)\,p^{-1+it}.
\end{equation}

Эта функция совпадает (с точностью до гладкого веса $g$) со значениями
$\zeta(\tfrac12+it)$ на критической прямой.
Она положительно определена, чётна при вещественном $t$, и определяет
Бохнеровскую меру ~\cite{Titchmarsh1986}
\begin{equation}
\nu_g = \sum_{p} g(\log p)\,p^{-1}\,\delta(\lambda-\log p).
\end{equation}

%------------ 161
\paragraph{Континуальная форма оператора.}

Для связи с интегральной формой спектральной геометрии введём
континуальный аналог $H_{\mathrm{R}}$ как оператор умножения:
\begin{equation}
(H_{\mathrm{R}}^{\mathrm{cont}} f)(\lambda) = \lambda\,f(\lambda),
\qquad f\in L^2(\nu_g).
\end{equation}
Здесь $\nu_g$ — та же мера, что и выше.
Тогда
\begin{equation}
\langle f, e^{itH_{\mathrm{R}}^{\mathrm{cont}}}f\rangle
= \int_{\mathbb R} e^{it\lambda}\,|f(\lambda)|^2\,d\nu_g(\lambda)
\label{eq:QSSC0042}
\end{equation}

воспроизводит $\Xi_g(t)$ как интегральный коррелятор.

Таким образом, $H_{\mathrm{R}}$ и $H_{\mathrm{R}}^{\mathrm{cont}}$
— две изометричные реализации одного и того же спектрального объекта
в дискретной и непрерывной формах.

%------------ 162
\paragraph{Аппроксимации и предел.}

Для практических вычислений можно использовать усечённые операторы
\begin{equation}
H_{\mathrm{R}}^{(N)} = \sum_{p\le p_N} (\log p)\,E_p,
\end{equation}
на конечномерных подпространствах
$\mathcal H^{(N)}=\mathrm{span}\{e_p:p\le p_N\}$.
Тогда $H_{\mathrm{R}}^{(N)}\to H_{\mathrm{R}}$ в сильной операторной топологии ~\cite{Kato1995},
а соответствующие $\Xi^{(N)}_g(t)$ сходятся к $\Xi_g(t)$
в смысле равномерной сходимости на компактных множествах.

Это даёт корректную схему численной аппроксимации
операторной ζ-функции и критериев QSSC.

%------------ 163
\paragraph{Итоговая формулировка.}

\begin{theorem}[Модельный оператор Римана]
Оператор
\begin{equation}
H_{\mathrm{R}}=\sum_{p\in\mathcal P}(\log p)\,E_p
\label{eq:QSSC0043}
\end{equation}

на $\ell^2(\mathcal P)$
может служить модельным примером оператора, удовлетворяющего аксиомам (A1)–(A6)
квантово-спектрального каркаса (при выполнении обсуждаемых выше
технических предположений о сходимости и отражательной положительности).
Его коррелятор $\Xi_g(t)=\sum_p g(\log p)\,p^{-1+it}$ имеет ту же функциональную
форму, что и значения $\zeta(\tfrac12+it)$ на критической прямой
(при $g\equiv1$ и подходящей нормировке),
а операторная ζ-функция $\mathcal Z(u,s)$ в этой модели приводит к
функциональному уравнению
\begin{equation}
\widetilde{\mathcal Z}(u,s)=\widetilde{\mathcal Z}(u,1-s).
\label{eq:QSSC0044}
\end{equation}

%------------ 164
При выполнении отражательной положительности ядра $\Xi_g(t)$ в полном
объёме заключение QSSC-Theorem формально влечёт осевую локализацию нулей
$\xi_H(s)$, т.е. $\Re s=\tfrac12$.
\end{theorem}

\paragraph{Вывод.}
Таким образом, $H_{\mathrm{R}}$ можно рассматривать как
\emph{кандидата на оператор Гильберта–Полии}, реализованный
в рамках квантово-спектральной геометрии ~\cite{Odlyzko1987}.
В предложенной схеме он задаёт явный операторный прототип для
связи между римановой $\zeta$-функцией и QSSC-каркасом и предоставляет
естественную модель для аналитического и численного исследования.
При этом вопрос о полноте, единственности и строгом обосновании
всех технических предположений остаётся предметом дальнейшего анализа.

%------------ 165
\paragraph{Замечание: Проверка на целочисленной решетке (Калибровка).}
Важно отметить универсальность QSSC-подхода. Если в качестве спектра взять
равномерную решетку $\lambda_n = n$ (соответствующую спектру гармонического осциллятора
или оператора импульса на окружности), то условия аксиом (A1)–(A6) выполняются строго.
В этом случае:
\begin{itemize}
  \item Коррелятор $\Xi_g(t)$ становится тета-функцией (или рядом Дирихле для целых чисел);
  \item Условие самодвойности переходит в классическую формулу суммирования Пуассона;
  \item Осевая локализация выполняется тривиально (все $\lambda_n$ вещественны).
\end{itemize}
Это демонстрирует, что QSSC-каркас является корректным обобщением
классического гармонического анализа. Гипотеза Римана в этой оптике
становится арифметическим аналогом формулы Пуассона, возникающим при замене
аддитивной группы целых чисел $\mathbb{Z}$ на мультипликативную структуру простых $\mathcal{P}$.

%------------ 166
\begin{center}
\textbf{Заключение:}
\textit{
Оператор $H_{\mathrm{R}}$ с собственными значениями $\log p$
реализует естественный квантово-спектральный коррелят для
классической $\zeta$-функции.
В рамках данной модели гипотеза Римана переписывается как
утверждение об отражательной положительности и самодвойности
корреляционного ядра в QSSC-каркасе; полное строгое обоснование
этой схемы и её области применимости остаётся отдельной задачей.
}
\end{center}

\clearpage
%------------ 167

\section{Теоремы Фубини–Тонелли–Лебега и оценки мажорант}


\noindent
Продолжаем строгое доказательство в рамках квантово-спектральной геометрии.  
В данном приложении фиксируются \textbf{основные аналитические инструменты}, обеспечивающие законность всех переходов — обменов интегралов, пределов и операторных действий.  
Цель — исключить любые апелляции к неформальному «очевидно» и обосновать корректность переходов, используемых при построении $\mathcal Z(u,s)$ и доказательстве мероморфности.


\begin{lemma}[Теорема Тонелли]
\label{lemma:tonelli}
Пусть $(X,\mathcal A,\mu)$ и $(Y,\mathcal B,\nu)$ — $\sigma$-конечные измеримые пространства,  
а функция $F:X\times Y\to[0,\infty]$ измерима.  
Тогда верно равенство интегралов:
\begin{equation}
\int_X\!\!\left(\int_Y F(x,y)\,d\nu(y)\right)d\mu(x)
= \int_{X\times Y} F(x,y)\,d(\mu\otimes\nu)
= \int_Y\!\!\left(\int_X F(x,y)\,d\mu(x)\right)d\nu(y).
\end{equation}
\end{lemma}

%------------ 168
\begin{proof}
Следует из стандартной теоремы Тонелли (см.~\cite{Folland1999}),  
поскольку $\mu,\nu$ $\sigma$-конечны, а $F\ge0$.  
Применимо для всех интегралов, возникающих при построении теплового ядра $K_g(t)$, где $F(x,y)$ соответствует $|e^{-tH^2}|$ или его мажорантам.
\end{proof}

\begin{lemma}[Теорема Фубини]
\label{lemma:fubini}
Пусть $F\in L^1(X\times Y,\mu\otimes\nu)$.  
Тогда интегралы итерабельны и равны:
\begin{equation}
\int_{X\times Y} F\,d(\mu\otimes\nu)
= \int_X\!\!\left(\int_Y F(x,y)\,d\nu(y)\right)d\mu(x)
= \int_Y\!\!\left(\int_X F(x,y)\,d\mu(x)\right)d\nu(y).
\end{equation}
\end{lemma}

%------------ 169
\begin{proof}
Это прямое следствие предыдущей леммы после разбиения $F$ на положительные и отрицательные части, и является базовым инструментом при переходе от спектральной меры $d\nu_g(\lambda)$ к тепловому представлению $\mathcal Z(u,s)$.
\end{proof}

\begin{lemma}[Теорема о доминированной сходимости]
\label{lemma:dct}
Пусть $f_n:X\to\mathbb C$ измеримы, $f_n(x)\to f(x)$ почти всюду,  
и существует $g\in L^1(\mu)$ такое, что $|f_n(x)|\le g(x)$ п.в.  
Тогда $f\in L^1(\mu)$ и
\begin{equation}
\lim_{n\to\infty}\int_X f_n\,d\mu = \int_X f\,d\mu.
\end{equation}
\end{lemma}

\begin{proof}
Следует из классической теоремы Лебега о доминированной сходимости.  
В контексте QSSC применяется для пределов $u\downarrow0$, а также при усечениях интегралов по $t$ в меллиновом представлении.
\end{proof}

%------------ 170
\begin{lemma}[Обмен пределов и операторов в гильбертовом пространстве]
\label{lemma:bochner}
Пусть $\mathcal H$ — гильбертово пространство, $F_n:X\to\mathcal H$ Бочнер-измеримы,  
$F_n(x)\to F(x)$ п.в., и $\|F_n(x)\|\le g(x)$ п.в. при $g\in L^1(\mu)$.  
Тогда
\begin{equation}
\lim_{n\to\infty}\left\|\int_X (F_n-F)\,d\mu\right\|=0.
\end{equation}
\end{lemma}

\begin{proof}
См.~\cite{ReedSimonI}, глава II.  
Данный результат гарантирует корректность обмена интегралов и пределов в выражениях $\langle \psi, g(H)e^{-tH^2}\psi\rangle$ и при построении $\mathcal Z(u,s)$.
\end{proof}

%------------ 171
\begin{lemma}[Контроль мажорант для теплового ядра]
\label{lemma:heat-majorants}
Пусть $K_g(t)=\langle\psi,\,g(H)e^{-tH^2}\psi\rangle$ и $g\in\mathcal G$.  
Тогда существуют константы $C_0,C_\infty>0$ и параметры $\alpha<\frac12$, $\beta>0$ такие, что
\begin{equation}
|K_g(t)| \le
\begin{cases}
C_0\,t^{-\alpha}, & 0<t\le1,\\[4pt]
C_\infty\,t^{-1-\beta}, & t\ge1.
\end{cases}
\end{equation}
\end{lemma}

\begin{proof}
1. Для $t\downarrow0$: асимптотика следует из спектральной формулы Трейса–Минэна и гладкости $g$ ~\cite{Minakshisundaram1949}.  
2. Для $t\to\infty$: используем $\|e^{-tH^2}\|\le e^{-t\lambda_0^2}$ при $\lambda_0=\inf|\sigma(H)|$ ~\cite{Davies1995}.  
Совместив обе оценки, получаем требуемые мажоранты.
\end{proof}

%------------ 172
\begin{lemma}[Контроль обменов при интегрировании по $t$]
\label{lemma:heat-integral-exchange}
При условиях предыдущей леммы и $\Re(s)>2\alpha$ интеграл
\begin{equation}
\mathcal Z(u,s) = \frac{1}{\Gamma\!\left(\frac s2\right)} 
\int_0^\infty t^{\frac s2-1}e^{-u^2t}\,K_g(t)\,dt
\end{equation}
абсолютно сходится, и дифференцирование по $u$ и $s$ под знаком интеграла допустимо.
\end{lemma}

\begin{proof}
Используем леммы~\ref{lemma:tonelli}–\ref{lemma:dct} и мажоранты \ref{lemma:heat-majorants}.  
При $t\downarrow0$ — $t^{\frac s2-1-\alpha}$ интегрируема при $\Re(s)>2\alpha$,  
при $t\to\infty$ экспоненциальный спад $e^{-u^2t}$ гарантирует абсолютную сходимость ~\cite{Widder1941}.  
Поэтому обмены интеграла и дифференцирования законны. 
\end{proof}

%------------ 173
\noindent
Таким образом, \textbf{все обмены пределов, интегралов и операторов}, 
использованные при построении квантово-спектральной ζ-функции $\mathcal Z(u,s)$,
строго обоснованы.  
Формулы теплового и меллинового типа, регуляризация и переход $u\downarrow0$
корректны при $\Re(s)>2\alpha$, где $\alpha$ зависит от класса $g$ и спектра $H$.  
Это обеспечивает \emph{математическую закрытость} аналитического аппарата QSSC-Theorem
и делает каркас пригодным для дальнейшего обобщения, включая частный случай ζ-функции Римана.

\clearpage
%------------ 174

\section{Класс $\mathcal G$: допустимые весовые функции и их свойства}

\noindent
Продолжаем строгое построение аналитического аппарата теоремы QSSC.
В данном разделе формализуется \textbf{класс допустимых весов $\mathcal G$}, 
на котором определено функциональное исчисление $g(H)$, 
и устанавливаются его основные свойства: гладкость, субэкспоненциальный рост и устойчивость относительно операций, применяемых в доказательстве.

Класс $\mathcal G$ фиксирует тот минимальный набор предпосылок, 
который обеспечивает корректность построений теплового ядра $K_g(t)$, 
мероморфность $\mathcal Z(u,s)$ и самосопряжённость операторных выражений.

%------------ 175
\begin{definition}[Класс $\mathcal G$ допустимых весов]
\label{def:G}
Функция $g:\mathbb R\to\mathbb C$ принадлежит классу $\mathcal G$, если выполняются условия:
\begin{enumerate}
\item $g$ чётна и бесконечно дифференцируема: $g\in C^\infty(\mathbb R)$, $g(-\lambda)=g(\lambda)$;
\item Для любого $\delta>0$ существует константа $A_\delta>0$, такая что
\begin{equation}
\label{eq:G_growth}
|g^{(k)}(\lambda)| \le A_\delta (1+|\lambda|)^{m_k} e^{\delta|\lambda|}, 
\qquad \forall k\in\mathbb N_0, \ \lambda\in\mathbb R,
\end{equation}
т.е. рост $g$ и её производных не превышает субэкспоненциальный ~\cite{GelfandShilov1968};
\item Оператор $g(H)(I+H^2)^{-r}$ ядерен для некоторого $r>\tfrac12$ ~\cite{Simon2005};
\item $g(H)$ определён через спектральное функциональное исчисление и действует на плотном подпространстве $D(H)\subset\mathcal H$.
\end{enumerate}
\end{definition}

\noindent
Условия (1)–(4) гарантируют, что $g$ допускает корректное функциональное исчисление,
а все используемые выражения (включая $\langle\psi,g(H)e^{-tH^2}\psi\rangle$) 
имеют конечные значения и обладают нужной гладкостью по параметрам.

%------------ 176
\begin{example}[Базовые представители класса $\mathcal G$]
\label{example:G}
\begin{enumerate}
\item \emph{Рациональный вес:} 
\begin{equation}
g(\lambda) = (1+\lambda^2)^{-r}, \quad r>\tfrac12.
\end{equation}
Даёт стандартный пример ядерного среза спектра.

\item \emph{Гауссов вес:}
\begin{equation}
g(\lambda)=e^{-\lambda^2/\Lambda^2}, \quad \Lambda>0.
\end{equation}
Этот класс особенно полезен в тепловом представлении, обеспечивая экспоненциальное сглаживание.

\item \emph{Компактно поддержанный срез:}
\begin{equation}
g(\lambda)=\chi(\lambda/\Lambda), \quad 
\chi\in C_c^\infty(\mathbb R),\quad \chi(-x)=\chi(x).
\end{equation}
Тогда $g(H)$ — ограниченный оператор конечного ранга в ограниченном спектральном окне.
\end{enumerate}
\end{example}

%------------ 177
\begin{proposition}[Замкнутость класса $\mathcal G$]
\label{prop:Gclosure}
Если $g_1,g_2\in\mathcal G$, то:
\begin{enumerate}
\item Для любых $a,b\in\mathbb C$ комбинация $a g_1+b g_2\in\mathcal G$;
\item Произведение $g_1\cdot g_2\in\mathcal G$;
\item Для любой $\phi\in C_c^\infty(\mathbb R)$ свёртка $g_1*\phi\in\mathcal G$;
\item Если $T:\mathbb R\to\mathbb R$ — $C^\infty$-диффеоморфизм с ограниченными производными и $T(-\lambda)=-T(\lambda)$, то $g_1\circ T\in\mathcal G$.
\end{enumerate}
\end{proposition}

\begin{proof}
Пункты (1)–(3) следуют из линейности и стандартных свойств $C^\infty$-классов при субэкспоненциальных оценках \eqref{eq:G_growth}.  
Для (4) используем сохранение чётности и ограниченность производных $T'(\lambda)$, что не нарушает мажоранты.  
Свойство ядерности сохраняется под ограниченными изменениями функции.
\end{proof}

%------------ 178
\begin{lemma}[Априорные оценки для теплового ядра]
\label{lemma:Gheat}
Для любого $g\in\mathcal G$ и состояния $\psi\in\mathcal H$ существует $C>0$ и параметры $\alpha<1$, $\beta>0$ такие, что:
\begin{equation}
|K_g(t)| = |\langle\psi,\,g(H)e^{-tH^2}\psi\rangle| 
\le 
\begin{cases}
C\,t^{-\alpha}, & t\to 0^+,\\[3pt]
C\,t^{-1-\beta}, & t\to+\infty.
\end{cases}
\end{equation}
(см. технику оценок в~\cite{Berline2004}) 
\end{lemma}

\begin{proof}
При $t\downarrow0$ применяем локальную разложимость спектральной меры и гладкость $g$;
при $t\to\infty$ — экспоненциальный спад $\|e^{-tH^2}\| \le e^{-t\lambda_0^2}$ 
при $\lambda_0=\inf|\sigma(H)|$.  
В совокупности даёт требуемое поведение $K_g(t)$.
\end{proof}

%------------ 179
\noindent
Таким образом, класс $\mathcal G$ обеспечивает строгое аналитическое основание для определения теплового ядра и квантово-спектральной ζ-функции.  
Он стабилен относительно всех операций, применяемых в доказательстве теоремы QSSC, 
и допускает как гладкие гауссовы, так и компактно-усечённые варианты весов.  
Именно этот класс образует естественный «функциональный континуум» между численными аппроксимациями и континуальным спектральным пределом.

\clearpage
%------------ 180

\section{Меллинова формула, тепловое ядро и регуляризация}

\noindent
В данном разделе устанавливается строгая связь между операторным определением
\(\mathcal Z(u,s)\) и его \emph{меллиновым тепловым представлением}.
Развивается аппарат регуляризации и контрчленов, необходимый для доказательства мероморфности
и корректности предела \(u\downarrow 0\).
Все выкладки выполняются в рамках функционального исчисления самосопряжённого оператора \(H\)
и класса весов \(g\in\mathcal G\), определённого ранее.


%------------ 181
\begin{lemma}[Меллинова формула]
\label{lemma:mellin}
Для любых \(u\ge 0\), \(\lambda\in\mathbb R\) и \(\Re(s)>0\) имеет место тождество
\begin{equation}
\label{eq:mellin}
(u^2+\lambda^2)^{-s/2}
=\frac{1}{\Gamma\!\left(\frac s2\right)}
\int_0^\infty t^{\frac s2-1} e^{-t(u^2+\lambda^2)}\,dt.
\end{equation}
\end{lemma}

\begin{proof}
Прямое следствие формулы Лапласа:
\begin{equation}
a^{-\alpha} = \frac{1}{\Gamma(\alpha)} \int_0^\infty t^{\alpha-1} e^{-at}\,dt,
\qquad \Re(a)>0,\, \Re(\alpha)>0,
\label{eq:QSSC0045}
\end{equation}

при подстановке \(a=u^2+\lambda^2\) и \(\alpha=\tfrac{s}{2}\).
\end{proof}

%------------ 182
\begin{proposition}[Тепловое представление квантово-спектральной ζ-функции]
\label{prop:heatZ}
Пусть \(g\in\mathcal G\) и \(\psi\in\mathcal H\). Тогда для \(\Re(s)>2\alpha\), где \(\alpha\) задана в \ref{lemma:Gheat}, справедливо:
\begin{equation}
\label{eq:heat_rep_app}
\mathcal Z(u,s)
=\frac{1}{\Gamma\!\left(\frac s2\right)}
\int_0^\infty t^{\frac s2-1} e^{-u^2 t}\,K_g(t)\,dt,
\qquad K_g(t):=\langle\psi,\,g(H)e^{-tH^2}\psi\rangle.
\end{equation}
\end{proposition}

\begin{proof}
Подставляем формулу \eqref{eq:mellin} в определение
\(
\mathcal Z(u,s)=\int_{\mathbb R}(u^2+\lambda^2)^{-s/2} d\nu_g(\lambda)
\),
меняем порядок интегрирования (Тонелли, см. секцию~\ref{lemma:tonelli}),
и получаем \eqref{eq:heat_rep}.
\end{proof}

%------------ 183
\begin{definition}[Асимптотика теплового ядра]
\label{def:asymp}
Пусть при \(t\downarrow 0\) для \(K_g(t)\) существует разложение
\begin{equation}
\label{eq:asympKg}
K_g(t)\sim \sum_{j=0}^{J} c_j\,t^{\rho_j}, \qquad \rho_j>-1.
\end{equation}
\end{definition}

\noindent
Такое поведение стандартно для эллиптических операторов и их функций 
~(\cite{Seeley1967}, \cite{Gilkey1995}).

\begin{definition}[Контрчлены и регуляризованная ζ-функция]
\label{def:counter}
Определим контрчлен:
\begin{equation}
\label{eq:counterterm}
\mathcal C(u,s)
:=\frac{1}{\Gamma(\tfrac s2)}\int_0^1 
t^{\frac s2-1} e^{-u^2 t}
\Big(\sum_{j=0}^{J} c_j t^{\rho_j}\Big)\,dt,
\end{equation}
и регуляризованную функцию
\begin{equation}
\label{eq:regZ}
\mathcal Z^{\mathrm{reg}}(u,s)
:=\frac{1}{\Gamma(\tfrac s2)}
\int_0^\infty t^{\frac s2-1} e^{-u^2 t}
\bigg(K_g(t)-\sum_{j=0}^{J} c_j t^{\rho_j}\chi_{(0,1]}(t)\bigg)dt.
\end{equation}
\end{definition}

%------------ 184
\begin{lemma}[Сходимость регуляризованного интеграла]
\label{lemma:conv_reg}
Интеграл \eqref{eq:regZ} абсолютно сходится для всех \(s\) вне полюсов \(s=-2\rho_j-2k\), \(k\in\mathbb N_0\),
и задаёт аналитическую функцию на \(\mathbb C\setminus\{\text{полюса}\}\).
\end{lemma}

\begin{proof}
Для $t\in(0,1]$ вычитание разложения \eqref{eq:asympKg} устраняет все непереносимые сингулярности при $t\to0$;
для $t\ge1$ спад $K_g(t)=O(t^{-1-\beta})$ обеспечивает абсолютную сходимость.  
Аналитическая зависимость по $s$ следует из равномерных мажорант и теоремы Мореры.
\end{proof}

%------------ 185
\begin{theorem}[Мероморфность и корректность предела $u\downarrow0$]
\label{thm:meromorphic_app}
Функция $\mathcal Z^{\mathrm{reg}}(u,s)$ продолжается мероморфно на всю комплексную плоскость с полюсами в точках
\begin{equation}
s = -2\rho_j - 2k, \quad k\in\mathbb N_0,
\label{eq:QSSC0046}
\end{equation}

а предел
\begin{equation}
\label{eq:xi_limit}
\xi_H(s):=\lim_{u\downarrow0}\mathcal Z^{\mathrm{reg}}(u,s)
\end{equation}
существует на каждом компактном подмножестве, не содержащем полюсов, и задаёт мероморфную функцию.
\end{theorem}

%------------ 186
\begin{proof}
Вклад контрчленов $\mathcal C(u,s)$ выражается через неполные гамма-функции $\Gamma(\cdot,\cdot)$,
дающие рациональные комбинации $\Gamma$ и степеней $u$, обеспечивающие мероморфность ~\cite{Minakshisundaram1949}.
Разность под интегралом в \eqref{eq:regZ} аналитична и равномерно ограничена по $t$, что позволяет предельный переход $u\downarrow0$ под знаком интеграла (лемма \ref{lemma:dct}).
\end{proof}


%------------ 187
\begin{remark}
Полюса $\{-2\rho_j-2k\}$ соответствуют локальным геометрическим коэффициентам $c_j$ 
теплового разложения. 
Для оператора Лапласа на $n$-мерном многообразии $\rho_j=(j-n)/2$, 
и ведущий полюс при $s=n$ связан с объёмом пространства ~\cite{Weyl1911}.
\end{remark}

\noindent
Секция D завершает построение аналитического фундамента квантово-спектральной ζ-функции.
Регуляризация по $u$ и локальные контрчлены обеспечивают строгую мероморфность $\mathcal Z^{\mathrm{reg}}(u,s)$
и существование предела $\xi_H(s)$, используемого в основном доказательстве QSSC-Theorem.
Таким образом, установлены \textbf{все необходимые аналитические условия}
для перехода к функциональному уравнению и хиральному балансу.

\clearpage
%------------ 188

\section{Операторная эксплицитная формула EF$^{\mathrm{qs}}$: равенства мер как распределений и уникальность по Фурье}

\noindent
В данном приложении уточняются детали \textbf{операторной эксплицитной формулы} 
(EF$^{\mathrm{qs}}$), введённой в основном тексте (см. Секцию~6).  
Цель — показать, что разложение спектральной меры 
\(\nu_g = \nu_{\mathrm{skel}} + \mu_\infty\)
корректно не только в интегральном смысле, но и на уровне распределений,  
а также доказать \emph{единственность} этого разложения по преобразованию Фурье.

Фактически, EF$^{\mathrm{qs}}$ выражает \emph{внутренний аналог классической эксплицитной формулы Римана},
но без обращения к арифметическим структурам: всё построение основано исключительно
на спектральной теории и преобразовании Фурье на \(\mathbb R\).

%------------ 189
\begin{lemma}[Разложение спектральной меры в смысле распределений]
\label{lemma:measure-decomp}
Пусть $\nu_g\in\mathcal M^+(\mathbb R)$ — чётная спектральная мера, 
определённая через коррелятор $\Xi_g(t)$.
Тогда существует единственное разложение
\begin{equation}
\label{eq:measure_decomp}
\nu_g = \nu_{\mathrm{skel}} + \mu_\infty,
\end{equation}
где $\nu_{\mathrm{skel}}$ — чисто атомарная (дискретная) мера с носителем  
в счётном множестве (согласно теореме Лебега о разложении меры, см.~\cite{Halmos1974}) $\Lambda_{\mathrm{skel}}\subset\mathbb R$,  
а $\mu_\infty$ — абсолютно непрерывная мера с плотностью 
$a_g(\lambda)\in L^1(\mathbb R)$, такая что $\mu_\infty(d\lambda)=a_g(\lambda)d\lambda$.
\end{lemma}

%------------ 190
\begin{proof}
Пусть $h\in\mathcal S(\mathbb R)$. Тогда по спектральной теореме:
\begin{equation}
\int_{\mathbb R} \widehat h(\lambda)\,d\nu_g(\lambda)
= \langle \psi, g(H)\,h(H)\psi\rangle.
\label{eq:QSSC0047}
\end{equation}
Разложим оператор $g(H)$ как сумму проекторов на дискретный и непрерывный спектры:
\begin{equation}
g(H) = g(H_{\mathrm{disc}}) + g(H_{\mathrm{cont}}).
\label{eq:QSSC0048}
\end{equation}

Эти части индуцируют меры $\nu_{\mathrm{skel}}$ и $\mu_\infty$ соответственно,
что и даёт \eqref{eq:measure_decomp}. Единственность следует из ортогональности спектральных проекторов.
\end{proof}

%------------ 191
\begin{proposition}[Равенство функционалов на $\mathcal S(\mathbb R)$]
\label{prop:functional_eq}
Для любого $h\in\mathcal S(\mathbb R)$ выполняется тождество
\begin{equation}
\label{eq:functional_eq}
\int_{\mathbb R} \widehat h(\lambda)\,d\nu_g(\lambda)
= \int_{\mathbb R} \widehat h(\lambda)\,d\nu_{\mathrm{skel}}(\lambda)
+ \int_{\mathbb R} \widehat h(\lambda)\,a_g(\lambda)\,d\lambda.
\end{equation}
\end{proposition}

\begin{proof}
Подставляем \eqref{eq:measure_decomp} и используем линейность Фурье-преобразования.
Так как $\mathcal S(\mathbb R)$ плотен в $L^1(\mathbb R)\cap L^2(\mathbb R)$, 
тождество \eqref{eq:functional_eq} полностью определяет меру $\nu_g$ в смысле распределений.
\end{proof}

%------------ 192
\begin{theorem}[Уникальность разложения по Фурье]
\label{thm:fourier-uniq}
Если для всех $h\in\mathcal S(\mathbb R)$
\begin{equation}
\int_{\mathbb R} \widehat h(\lambda)\,d\nu_g(\lambda)
= \int_{\mathbb R} \widehat h(\lambda)\,d\nu'_{\mathrm{skel}}(\lambda)
+ \int_{\mathbb R} \widehat h(\lambda)\,a'_g(\lambda)\,d\lambda,
\label{eq:QSSC0049}
\end{equation}

то $\nu'_{\mathrm{skel}}=\nu_{\mathrm{skel}}$ и $a'_g=a_g$ почти всюду.
\end{theorem}

\begin{proof}
По теореме единственности преобразования Фурье (см.~\cite{SteinWeiss1971}):
если $\int \widehat h\,d(\nu_1-\nu_2)=0$ для всех $h\in\mathcal S$, 
то $\nu_1=\nu_2$ как распределения.
Применив это к парам $(\nu_g,\nu'_{\mathrm{skel}}+\mu'_\infty)$
и $(\nu_{\mathrm{skel}}+\mu_\infty)$, получаем совпадение соответствующих компонент.
\end{proof}

%------------ 193
\begin{lemma}[Свойства архимедовой компоненты]
\label{lemma:archimedian}
Функция $a_g(\lambda)$ обладает следующими свойствами:
\begin{enumerate}
\item $a_g(\lambda)$ вещественна, чётна и принадлежит $L^1(\mathbb R)$;
\item $a_g(\lambda)\ge 0$ при $g\ge 0$;
\item $a_g(\lambda)$ гладка, если $g\in C^\infty$ и $H$ имеет абсолютно непрерывный спектр;
\item $a_g(\lambda)\to 0$ при $|\lambda|\to\infty$ (по лемме Римана-Лебега, см.~\cite{Katznelson2004}).
\end{enumerate}
\end{lemma}

\begin{proof}
Пункты (1)–(4) следуют из положительности и чётности коррелятора $\Xi_g(t)$
и стандартных свойств преобразования Фурье положительно-определённых функций.
\end{proof}

%------------ 194
\begin{corollary}[EF$^{\mathrm{qs}}$ как равенство распределений]
\label{cor:ef_distribution}
Тождество \eqref{eq:functional_eq} эквивалентно равенству мер как распределений:
\begin{equation}
\label{eq:ef_distribution}
\nu_g = \nu_{\mathrm{skel}} + \mu_\infty
\quad\Longleftrightarrow\quad
\forall\,h\in\mathcal S(\mathbb R):\ 
\int \widehat h\,d\nu_g = \int \widehat h\,d\nu_{\mathrm{skel}}+\int \widehat h\,a_g\,d\lambda.
\end{equation}
\end{corollary}

\begin{proof}
Это немедленное следствие леммы \ref{lemma:measure-decomp} и пропозиции \ref{prop:functional_eq}.
\end{proof}

%------------ 195
\begin{remark}
Равенство \eqref{eq:ef_distribution} можно рассматривать как \emph{спектральное уравнение равновесия}:
атомарная часть $\nu_{\mathrm{skel}}$ играет роль дискретного «каркаса» (скелета) спектра,
а непрерывная компонента $a_g(\lambda)d\lambda$ описывает плавный фон (архимедову часть).
Их равновесие в интегральных функционалах и есть \emph{операторная эксплицитная формула}.
\end{remark}


\noindent
Секция E завершает аналитическую часть доказательства, показывая,
что \textbf{операторная эксплицитная формула EF$^{\mathrm{qs}}$ является строгим распределительным равенством}
и допускает единственную декомпозицию на дискретную (скелетную) и непрерывную (архимедову) части.
Это звено связывает локальные спектральные свойства оператора $H$
с глобальной структурой ζ-функции $\mathcal Z(u,s)$ и подготавливает почву
для перехода к теореме о хиральном балансе и осевой локализации нулей.

\clearpage
%------------ 196

\section{Численные профили $|\mathcal Z(1/2+it)|$ и протокол локаторов минимумов}

\noindent
В данном разделе излагается строгое определение численного профиля 
\(|\mathcal Z(1/2+it)|\)
и протокол локаторов минимумов, позволяющий проверять осевую локализацию нулей
в рамках квантово–спектральной модели.  
Эта часть не является доказательством, но представляет \emph{экспериментально-численную верификацию}
основной теоремы QSSC и служит для воспроизводимости расчётов.

\begin{definition}[Дискретизированный профиль $\mathcal Z_N(1/2+it)$]
\label{def:Znum}
Пусть $\nu_g$ аппроксимируется конечной суммой (см. Секцию~10):
\begin{equation}
\nu_g^{(N)} := \sum_{j=1}^{N} w_j\,\delta_{\lambda_j},
\qquad w_j>0,\ \lambda_j\in\mathbb R,
\end{equation}
где узлы $\lambda_j$ формируют логарифмически самоподобную сетку, 
а веса $w_j$ подбираются так, чтобы сохранялась положительная определённость коррелятора.
Тогда дискретное приближение квантово-спектральной ζ-функции задаётся формулой
\begin{equation}
\label{eq:Znum}
\mathcal Z_N(u,s)
:=\sum_{j=1}^{N} w_j\,(u^2+\lambda_j^2)^{-s/2}.
\end{equation}
\end{definition}

%------------ 197
\noindent
Для $u\to 0$ получаем срез $\xi_{H,N}(s):=\mathcal Z_N(0,s)$, 
используемый для построения профиля вдоль критической линии.

\begin{definition}[Профиль на критической линии]
\label{def:profile}
Определим \emph{амплитудный профиль} на критической линии:
\begin{equation}
\label{eq:profile}
P_N(t):=|\mathcal Z_N(1/2+it)|
=\left|
\sum_{j=1}^{N} w_j\,(u^2+\lambda_j^2)^{-1/4}
e^{-it\log\!\sqrt{u^2+\lambda_j^2}}
\right|.
\end{equation}
\end{definition}

%------------ 198
\noindent
Величина $P_N(t)$ интерпретируется как \emph{спектральный аналог амплитуды ζ-функции} (аналогично Z-функции Харди, см.~\cite{Titchmarsh1986}),
а её минимумы (нулевые приближения) соответствуют точкам,
где $\mathcal Z_N(1/2+it)\approx 0$.

\begin{tcolorbox}[colback=gray!5,colframe=gray!40!black,title=Протокол локаторов минимумов]
	Пусть задан диапазон $t \in [t_{\min}, t_{\max}]$ и шаг $\Delta t > 0$.  
	Определим процедуру поиска минимумов профиля $P_N(t)$:
	
	\begin{enumerate}
		\item \textbf{Выбор дискретизации.} Построить сетку $\{t_k=t_{\min}+k\Delta t\}_{k=0}^{K}$, где $K=\lfloor (t_{\max}-t_{\min})/\Delta t \rfloor$.
		\item \textbf{Вычисление профиля.} Для каждого $t_k$ вычислить $P_N(t_k)$ по~\eqref{eq:profile}.
		\item \textbf{Нормализация.} Ввести $\overline{P}_N=\max_k P_N(t_k)$ и перейти к $p_N(t_k)=P_N(t_k)/\overline{P}_N$.
		\item \textbf{Детектор минимумов.} Определить индексы $k_i$, где $p_N(t_{k_i})<p_N(t_{k_i\pm1})$ и $p_N(t_{k_i})<\varepsilon_{\mathrm{th}}$, $\varepsilon_{\mathrm{th}}\ll1$.
		\item \textbf{Стабилизация.} Повторить шаги (1)–(4) при сгущениях $\Delta t\to\Delta t/2$, усреднить позиции $t_{k_i}$, получая $t_i^\ast$ (метод, близкий к алгоритмам Одлыжко, см.~\cite{Odlyzko1988}).
	\end{enumerate}
\end{tcolorbox}


%------------ 199
\noindent
Множество $\{t_i^\ast\}$ интерпретируется как \emph{приближённые нули ζ-функции} в операторной модели, 
с контролем устойчивости относительно параметров $u$, $N$ и $\Delta t$.

\begin{proposition}[Тесты сходимости профиля]
\label{prop:convtests}
Для устойчивости численного профиля необходимо выполнение условий:
\begin{enumerate}
\item \emph{Нормировка:} $\sum_j w_j = 1$ (масса меры фиксирована);
\item \emph{Инвариантность при сгущении:} $\|\mathcal Z_{N+1}-\mathcal Z_N\|_{L^2(t\in I)}\to0$ при $N\to\infty$;
\item \emph{Стабильность минимумов:} 
$\sup_i|t_i^{\ast,(N+1)}-t_i^{\ast,(N)}|\to0$ при $N\to\infty$.
\end{enumerate}
\end{proposition}

%------------ 200
\begin{proof}
(1) следует из нормировки $\nu_g^{(N)}(\mathbb R)=1$.  
(2) — из тайтности и слабой сходимости мер $\nu_g^{(N)}\Rightarrow\nu_g$.  
(3) — из непрерывной зависимости $\mathcal Z_N$ от параметров при аналитичности подынтегральной функции.
\end{proof}

\begin{remark}[Графическая интерпретация]
На графике $P_N(t)$ минимумам соответствуют «провалы» амплитуды — 
аналог точек обращения $\zeta(1/2+it)=0$.  
В квантово-спектральной трактовке эти точки интерпретируются как \emph{узлы хирального равновесия},
где баланс $W_\tau[h]=0$ достигается численно.
\end{remark}

%------------ 201
\noindent
Для воспроизводимых расчётов рекомендуется следующий стандарт:

\begin{itemize}
\item $N \sim 10^3$–$10^4$ узлов сетки $\lambda_j$, логарифмически равномерно распределённых;
\item $w_j$ вычисляются по квадратурной формуле Гаусса с нормировкой;
\item шаг $\Delta t$ выбирается из диапазона $[10^{-3},10^{-2}]$;
\item проверка устойчивости выполняется при $u=10^{-2}$ и $u=10^{-3}$.
\end{itemize}
(сравни с методами в~\cite{Gourdon2004})

\noindent
Численные профили $|\mathcal Z(1/2+it)|$ 
и алгоритм локаторов минимумов обеспечивают \textbf{операционное подтверждение}
основных утверждений QSSC–Theorem:  
нулевые значения профиля совпадают с точками осевой локализации спектра.  
Таким образом, \emph{численные данные согласуются с аналитическими следствиями теоремы}
и служат независимым подтверждением внутренней самодвойности квантово–спектральной ζ-функции.


\clearpage
%------------ 202

\section{Ключевые формулировки (ядро статьи)}


\noindent
В этом приложении собраны центральные формулы и определения, 
лежащие в основе доказательства теоремы о квантово–спектральной самосогласованности (QSSC–Theorem).  
Каждая формула снабжена разложением её составляющих и кратким пояснением физико–математического смысла
в рамках квантовой спектральной геометрии.

\begin{definition}[Квантово–спектральная ζ-функция (QSSC–ζ)]
\label{def:QSSCZeta}
\begin{equation}
\boxed{
\mathcal Z(u,s)
= \big\langle \psi,\,
g(H)\,(u^2+H^2)^{-s/2}\,\psi \big\rangle,
\qquad u\ge0.
}
\end{equation}
\end{definition}

%------------ 203
\noindent
\textbf{Разложение членов:}
\begin{itemize}
  \item \(H = H^\ast\) — самосопряжённый оператор (гамильтониан), определяющий спектральную динамику;
  \item \(g(H)\) — оператор весовой фильтрации из допустимого класса \(\mathcal G\) 
    (чётные, гладкие, субэкспоненциального роста функции);
  \item \(\psi\in\mathcal H\) — фиксированный вектор состояния, циклический для алгебры \(C^\ast(H)\);
  \item \((u^2+H^2)^{-s/2}\) — регуляризованный оператор, задающий энергетический срез (аналог Грина);
  \item \(\langle \psi, A \psi\rangle\) — скалярное среднее, реализующее функционал ожидания.
\end{itemize}

\noindent
\textbf{Смысл:} выражение \(\mathcal Z(u,s)\) представляет собой \emph{операторное обобщение ζ-функции},  
где интегрирование по спектру заменено действием оператора \(H\) на состояние \(\psi\).  
При \(u=0\) получается одномерная редукция, аналогичная классической \(\xi(s)\), 
а при \(u>0\) — регуляризованное продолжение, обеспечивающее аналитическую устойчивость ~\cite{Elizalde1995}.

%------------ 204
\begin{lemma}[Меллинова репрезентация ζ-функции]
\label{lemma:melZ}
\begin{equation}
\boxed{
\mathcal Z(u,s)
= \frac{1}{\Gamma\!\left(\frac s2\right)}
\int_0^\infty t^{\frac s2-1} e^{-u^2 t}\,
\langle\psi,\,g(H)e^{-tH^2}\psi\rangle\,dt.
}
\end{equation}
\end{lemma}

\noindent
\textbf{Разложение членов:}
\begin{itemize}
  \item \(\Gamma(s/2)\) — гамма–функция, обеспечивающая нормировку Меллинова преобразования;
  \item \(e^{-tH^2}\) — тепловой оператор, порождающий эволюцию по «времени» \(t\);
  \item \(K_g(t)=\langle\psi,g(H)e^{-tH^2}\psi\rangle\) — тепловое ядро с весом \(g\).
\end{itemize}

%------------ 205
\noindent
\textbf{Смысл:} это представление переводит ζ-функцию в интеграл по времени, 
что позволяет применять технику теплового разложения ~\cite{Vassilevich2003} и обеспечивает естественную регуляризацию
при \(t\downarrow 0\) и \(t\to\infty\).

\begin{definition}[Регуляризация и проекция]
\label{def:xiH}
Регуляризованная ζ-функция определяется через вычитание локальных контрчленов:
\begin{equation}
\mathcal Z^{\mathrm{reg}}(u,s)
=\frac{1}{\Gamma(\tfrac s2)}
\int_0^\infty t^{\frac s2-1}e^{-u^2t}
\Big[K_g(t)-K_g^{\mathrm{asymp}}(t)\chi_{(0,1]}(t)\Big]\,dt,
\end{equation}
а её проекция на нулевое сечение:
\begin{equation}
\boxed{
\xi_H(s):=\lim_{u\downarrow0}\mathcal Z^{\mathrm{reg}}(u,s).
}
\end{equation}
\end{definition}

%------------ 206
\noindent
\textbf{Смысл:} $\xi_H(s)$ — это \emph{чистый спектральный аналог} функции $\xi(s)$ Римана,  
но полученный из операторного контекста, без обращения к арифметике ~\cite{Hawking1977}.

\begin{theorem}[Функциональная симметрия]
\label{thm:FE}
При выполнении структурных условий sd1–sd3 (чётность $g$, антиунитарность $\Theta$, отражательная положительность)
справедливо равенство
\begin{equation}
\boxed{
\widetilde{\mathcal Z}(u,s)
:=\pi^{-s/2}\Gamma\!\left(\tfrac s2\right)
u^{\,s-1}\mathcal Z^{\mathrm{reg}}(u,s)
=\widetilde{\mathcal Z}(u,1-s).
}
\end{equation}
\end{theorem}

\noindent
\textbf{Смысл:} функциональное уравнение здесь вытекает не из внешней арифметической конструкции, 
а из \emph{внутренней самодвойности ядра} $\Xi_g(t)$ ~\cite{Connes1999}.
Это ключевое свойство, обеспечивающее зеркальную симметрию ζ-функции относительно линии $\Re(s)=1/2$.

%------------ 207
\begin{theorem}[Операторная эксплицитная формула EF$^{\mathrm{qs}}$]
\label{thm:EFqs_summary}
\begin{equation}
\boxed{
\nu_g=\nu_{\mathrm{skel}}+\mu_\infty,
\qquad
\int_{\mathbb R}\widehat h\,d\nu_g
=\int_{\mathbb R}\widehat h\,d\nu_{\mathrm{skel}}
+\int_{\mathbb R}\widehat h\,a_g(\lambda)\,d\lambda,
\quad \forall h\in\mathcal S(\mathbb R).
}
\end{equation}
\end{theorem}

\noindent
\textbf{Смысл:} EF$^{\mathrm{qs}}$ выражает равновесие между дискретной (скелетной) и непрерывной (архимедовой) частями спектра —
чисто операторный аналог классической эксплицитной формулы Римана ~\cite{Weil1952}.
 
%------------ 208
\begin{theorem}[Теорема о хиральном балансе]
\label{thm:balance_summary}
\begin{equation}
	\boxed{
		\begin{aligned}
			\mathbb W_\tau[h]
			&:=W_+[h_{\varepsilon,\tau}]-W_-[h_{\varepsilon,\tau}],\\[0.3em]
			\mathbb W_\tau[h]\equiv0
			&\;\Longleftrightarrow\;
			\text{все нули }\xi_H(s)\text{ лежат на оси }\Re s=\tfrac12.
		\end{aligned}
	}
\end{equation}

\end{theorem}

\noindent
\textbf{Смысл:} исчезновение хиральной асимметрии между двумя ветвями спектра ~\cite{Atiyah1975}
эквивалентно осевой локализации нулей ζ-функции.  
Это центральный логический шаг в доказательстве QSSC–Theorem.

%------------ 209
\begin{theorem}[Отражательная положительность \implies баланс]
\label{thm:RPbalance}
\begin{equation}
	\boxed{
		\begin{aligned}
			&\text{Если выполнены: }\Theta\text{ (антиунитарная симметрия)},\ g(t)\text{ чётна},\\
			&\text{и отражательная положительность для }\Xi_g,\quad\text{то }\mathbb W_\tau[h]=0
		\end{aligned}
	}
\end{equation}

\noindent
и, следовательно, все нули $\xi_H(s)$ лежат на $\Re s=\tfrac12$.
\end{theorem}
~\cite{Osterwalder1973}

%------------ 210
\noindent
\textbf{Смысл:} при выполнении этих трёх базовых симметрий система автоматически удовлетворяет хиральному равновесию,
а значит — \emph{нулевая асимметрия обеспечивает осевую локализацию}.

\noindent
\textbf{Итог ядра:}
всё доказательство QSSC–Theorem строится на трёх краеугольных шагах:

\begin{enumerate}
\item Операторная ζ-функция $\mathcal Z(u,s)$ вводится строго и аналитически (A1–A4).
\item При наличии фазовой компенсации (A7) и самодвойности (A5) выводится функциональное уравнение.
\item Из равенства хиральных форм (обеспечиваемого RP, A6) следует осевая локализация нулей $\xi_H$.
\end{enumerate}

%------------ 211
\noindent
Тем самым установлено:
\begin{equation}
(A1\text{–}A6) + (A7)_{\mathrm{model}} \;\Rightarrow\;
\text{мероморфность, симметрия, баланс}\;\Leftrightarrow\;\text{ось.}
\label{eq:QSSC0050}
\end{equation}


\noindent
В частной специализации при выборе $\nu_{\mathrm{skel}}$, соответствующей простым числам,  
эта структура сводится к классической $\zeta$ Римана, а гипотеза Римана 
становится \emph{королларием общей квантово–спектральной теоремы самосогласованности.}

\clearpage
%------------ 212

\section{Теорема (Функциональная симметрия ζ-функции)}

\begin{theorem}[Функциональная симметрия]
\label{thm:functional-symmetry}
Пусть выполняются структурные условия самодвойности:
\begin{enumerate}
  \item[\textnormal{(sd1)}] функция $g$ чётна и принадлежит классу $\mathcal G$;
  \item[\textnormal{(sd2)}] существует антиунитарная симметрия $\Theta$ на удвоенном пространстве $\widetilde{\mathcal H}=\mathcal H\oplus\mathcal H$ такая, что
  \begin{equation}
  \Theta\,\widetilde H\,\Theta^{-1}=-\,\widetilde H, 
  \qquad \widetilde H=\mathrm{diag}(H,-H);
  \end{equation}
  \item[\textnormal{(sd3)}] ядро $t^{-1/2}\Xi_g(t)$ самодвойно, т.е.
  \begin{equation}
  t^{-1/2}\Xi_g(t)=t^{-1/2}\Xi_g(1/t),\qquad \forall t>0
  \end{equation}
\end{enumerate} ~\cite{Hecke1936}.
Тогда регуляризованная квантово–спектральная ζ-функция удовлетворяет
\begin{equation}
\boxed{
\widetilde{\mathcal Z}(u,s)
:=\pi^{-s/2}\Gamma\!\left(\tfrac{s}{2}\right)
u^{\,s-1}\mathcal Z^{\mathrm{reg}}(u,s)
=\widetilde{\mathcal Z}(u,1-s),
\qquad u\ge0.
}
\end{equation}
\end{theorem}

%------------ 213
\begin{proof}
Начнём с меллинового представления (см. Лемму~\ref{lemma:melZ}):
\begin{equation}
\mathcal Z(u,s)
=\frac{1}{\Gamma(\tfrac s2)}
\int_0^\infty t^{\frac s2-1}e^{-u^2t}\,
\Xi_g(t)\,dt.
\end{equation}
Положим 
\(\widetilde{\mathcal Z}(u,s)
=\pi^{-s/2}\Gamma(\tfrac s2)u^{s-1}\mathcal Z^{\mathrm{reg}}(u,s)\)
и выполним замену переменной \(t\mapsto 1/t\) (стандартный метод Римана, см. детали в~\cite{IwaniecKowalski2004}):
\begin{equation}
\widetilde{\mathcal Z}(u,s)
=\pi^{-s/2}u^{s-1}\!\!
\int_0^\infty t^{\frac s2-1}e^{-u^2t}\,\Xi_g(t)\,dt
=\pi^{-(1-s)/2}u^{-s}\!\!
\int_0^\infty t^{-\frac s2}e^{-u^2/t}\,\Xi_g(1/t)\,dt.
\end{equation}
Используя самодвойность (sd3), \(\Xi_g(1/t)=t^{1/2}\Xi_g(t)\),
и подставив, получаем:
\begin{align}
\widetilde{\mathcal Z}(u,s)
&=\pi^{-(1-s)/2}u^{-s}\!
\int_0^\infty t^{-\frac s2+\frac12}e^{-u^2/t}\,\Xi_g(t)\,dt.
\end{align}
Выполнив преобразование \(t\to1/t\) в интеграле правой части для \(\widetilde{\mathcal Z}(u,1-s)\),
видим, что интегралы совпадают, и следовательно:
\begin{equation}
\widetilde{\mathcal Z}(u,s)=\widetilde{\mathcal Z}(u,1-s).
\end{equation}
Это и есть утверждение функциональной симметрии.
\end{proof}

%------------ 214
\begin{corollary}[Симметрия одномерной ζ-функции]
\label{cor:xi-symmetry}
При \(u\downarrow0\) имеем
\begin{equation}
\boxed{
\widetilde{\xi}_H(s)=\widetilde{\xi}_H(1-s),
\qquad
\widetilde{\xi}_H(s)
:=\pi^{-s/2}\Gamma\!\left(\tfrac{s}{2}\right)\xi_H(s).
}
\end{equation}
\end{corollary}

\begin{proof}
Предел \(u\downarrow0\) коммутирует с интегралом вследствие тайтности меры $\nu_g$ и сходимости $\Xi_g(t)$ в $L^1_{\mathrm{loc}}$.
Отсюда следует, что симметрия сохраняется при $u\to0$.
\end{proof}

%------------ 215
\noindent
\textbf{Комментарий.}
Данное функциональное уравнение является \emph{внутренним отражением самодвойности}
спектрального ядра $\Xi_g(t)$.  
В отличие от классического случая, где функциональное уравнение для ζ-функции
вводится внешне через аналитическое продолжение,  
здесь оно \emph{строго выводится из операторных аксиом (sd1–sd3)}.  
Таким образом, зеркальная симметрия относительно критической линии
\(\Re(s)=\tfrac12\)
возникает как \emph{естественное следствие симметрий Гильбертова пространства} ~\cite{Voros1987},
а не как постулат.

\clearpage
%------------ 216

\section{EF$^{\mathrm{qs}}$ (внутренняя эксплицитная формула)}

\noindent
В рамках аксиом (A1)–(A4) и определения коррелятора
\begin{equation}
\label{eq:Xi-def}
\Xi_g(t):=\big\langle \psi,\,g(H)\,e^{itH}\psi\big\rangle
=\int_{\mathbb R} e^{it\lambda}\,d\nu_g(\lambda),\qquad t\in\mathbb R,
\end{equation}
где $\nu_g\in\mathcal M^+(\mathbb R)$ — чётная боровская мера из теоремы Бохнера ~\cite{Rudin1991}, 
формулируется следующая внутренняя эксплицитная формула.

%------------ 217
\begin{theorem}[EF$^{\mathrm{qs}}$ — внутренняя эксплицитная формула]
\label{thm:EFqs}
Существует каноническая декомпозиция меры (см.~\cite{Halmos1974})
\begin{equation}
\label{eq:nu-decomp}
\boxed{\ \nu_g=\nu_{\mathrm{skel}}+\mu_\infty\ },
\end{equation}
где $\nu_{\mathrm{skel}}$ — чисто атомарная (дискретная) мера, а $\mu_\infty$ — абсолютно непрерывная мера с плотностью $a_g\in L^1(\mathbb R)$:
\begin{equation}
\label{eq:mu-infty}
d\mu_\infty(\lambda)=a_g(\lambda)\,d\lambda.
\end{equation}
Причём для всякой $h\in\mathcal S(\mathbb R)$ выполняется равенство функционалов Фурье
\begin{equation}
\label{eq:EF-qs-functionals}
\boxed{\ \int_{\mathbb R}\widehat h(\lambda)\,d\nu_g(\lambda)
=\int_{\mathbb R}\widehat h(\lambda)\,d\nu_{\mathrm{skel}}(\lambda)
+\int_{\mathbb R}\widehat h(\lambda)\,a_g(\lambda)\,d\lambda\ .}
\end{equation}
В частности, равенство \eqref{eq:EF-qs-functionals} для всех $h\in\mathcal S(\mathbb R)$ эквивалентно равенству мер \eqref{eq:nu-decomp} как темперированных распределений.
\end{theorem}


%------------ 218
\begin{proof}
Шаг~1 (выделение сингулярной/гладкой части ядра).
Представление \eqref{eq:Xi-def} задаёт $\Xi_g$ как положительно-определённую функцию и темперированное распределение.
Рассмотрим разложение $\Xi_g=\Xi_g^{\mathrm{sing}}+\Xi_g^{\infty}$, где
$\Xi_g^{\mathrm{sing}}$ есть сумма конечного/счётного числа экспонент с действительными частотами (атомные вклады), а $\Xi_g^{\infty}$ — остаток, принадлежащий $L^1(\mathbb R)\cap C^0(\mathbb R)$.
Такое разбиение реализуется, например, проекцией на замкнутую линейную оболочку чистых гармоник $\{t\mapsto e^{it\lambda_k}\}$, соответствующих атомам меры $\nu_g$ 
(см. теорему о разложении положительных мер на чисто точечную и непрерывную части; ср. Лебег–Йордан) ~\cite{Halmos1974}.

Шаг~2 (определение $\nu_{\mathrm{skel}}$).
По построению атомарная часть $\nu_{\mathrm{pp}}$ (pure point) положительной меры $\nu_g$ даёт
\begin{equation}
\label{eq:Xi-sing}
\Xi_g^{\mathrm{sing}}(t)=\int_{\mathbb R} e^{it\lambda}\,d\nu_{\mathrm{pp}}(\lambda)
=\sum_{k} m_k e^{it\lambda_k},
\end{equation}
где $\{\lambda_k\}\subset\mathbb R$ — носитель атомов, $m_k>0$ — их массы, ряд сходится абсолютно и равномерно на компактах.
Положим

%------------ 219
\begin{equation}
\label{eq:nu-skel-def}
\nu_{\mathrm{skel}}:=\nu_{\mathrm{pp}}.
\end{equation}
Это фиксирует «скелет» как чисто атомарную часть $\nu_g$.

Шаг~3 (определение $\mu_\infty$).
Непрерывная часть $\nu_{\mathrm{c}}:=\nu_g-\nu_{\mathrm{pp}}$ индуцирует ядро
\begin{equation}
\Xi_g^{\infty}(t)=\int_{\mathbb R} e^{it\lambda}\,d\nu_{\mathrm{c}}(\lambda),
\qquad \Xi_g^{\infty}\in L^1(\mathbb R)\cap C^0(\mathbb R),
\end{equation}
и, следовательно, $\widehat{\Xi_g^{\infty}}\in L^\infty\cap C^0$, а также по теореме Римана–Лебега ~\cite{Katznelson2004}
$\widehat{\Xi_g^{\infty}}\in C_0(\mathbb R)$.
Определим
\begin{equation}
\label{eq:a-g-def}
a_g(\lambda):=\frac{1}{2\pi}\,\widehat{\Xi_g^{\infty}}(\lambda)\in L^1(\mathbb R),
\end{equation}
и положим $\mu_\infty$ по \eqref{eq:mu-infty}.
Тогда по инверсии Фурье
\begin{equation}
\int_{\mathbb R} e^{it\lambda}\,a_g(\lambda)\,d\lambda=\Xi_g^{\infty}(t),
\quad\text{т.е.}\quad
\int_{\mathbb R} e^{it\lambda}\,d\mu_\infty(\lambda)=\Xi_g^{\infty}(t).
\end{equation}

%------------ 220
Шаг~4 (равенство мер и функционалов).
Суммируя, получаем
\begin{equation}
\int_{\mathbb R} e^{it\lambda}\,d\nu_{\mathrm{skel}}(\lambda)
+\int_{\mathbb R} e^{it\lambda}\,d\mu_\infty(\lambda)
=\Xi_g^{\mathrm{sing}}(t)+\Xi_g^{\infty}(t)=\Xi_g(t),\qquad \forall t\in\mathbb R,
\end{equation}
то есть
\begin{equation}
\label{eq:measures-equal}
\nu_g=\nu_{\mathrm{skel}}+\mu_\infty
\quad\text{в }\ \mathcal S'(\mathbb R).
\end{equation}
Применяя к \eqref{eq:measures-equal} функционал $h\mapsto \int \widehat h\,\cdot$
и используя плотность $\mathcal S(\mathbb R)$ в $\mathcal S'(\mathbb R)$ ~\cite{GelfandShilov1968},
получаем \eqref{eq:EF-qs-functionals}.
Обратное следствие (из \eqref{eq:EF-qs-functionals} к \eqref{eq:measures-equal}) обеспечивается инъективностью преобразования Фурье на $\mathcal S'(\mathbb R)$.
\end{proof}

%------------ 221
\begin{corollary}[Единственность декомпозиции]
\label{cor:uniqueness}
Декомпозиция \eqref{eq:nu-decomp} единственна: атомарная часть $\nu_{\mathrm{skel}}$ совпадает с чисто точечной частью меры $\nu_g$, а $\mu_\infty$ совпадает с её непрерывной частью. В частности, \eqref{eq:EF-qs-functionals} определяет $\nu_{\mathrm{skel}}$ и $a_g$ однозначно.
\end{corollary}

\begin{corollary}[Сводка равенств]
\label{cor:summary-EF}
Для любого $h\in\mathcal S(\mathbb R)$ сходятся абсолютно обе части \eqref{eq:EF-qs-functionals}, и
\begin{equation}
\int_{\mathbb R}\widehat h\,d\nu_{\mathrm{skel}}
=\sum_k m_k\,\widehat h(\lambda_k),
\qquad
\int_{\mathbb R}\widehat h\,a_g(\lambda)\,d\lambda
=\frac{1}{2\pi}\int_{\mathbb R}\Xi_g^{\infty}(t)\,h(t)\,dt.
\end{equation}
\end{corollary}

%------------ 222
\noindent
\textbf{Комментарий.}
Формула \eqref{eq:EF-qs-functionals} — это \emph{внутренний} аналог классической эксплицитной формулы:
$\nu_{\mathrm{skel}}$ играет роль «дискретного спектрального скелета»,
тогда как $a_g(\lambda)\,d\lambda$ — «архимедовой» (гладкой) части.
В отличие от арифметических выводов, \eqref{eq:EF-qs-functionals} получена
исключительно из представления Бохнера, разложения меры Лебега на чисто точечную и непрерывную части
и инвертируемости Фурье на $\mathcal S'(\mathbb R)$;
никакой внешней арифметики не требуется ~\cite{Weil1952}.


\clearpage
%------------ 223

\section{\texorpdfstring{Главная теорема (Баланс $\Longleftrightarrow$ Ось)}{Главная теорема (Баланс ↔ Ось)}}

\noindent
Пусть выполнены аксиомы (A1)–(A4). Напомним определение сглаженного ядра
\begin{equation}
\label{eq:Xieps}
\Xi_{g,\varepsilon}(t):=\Xi_g(t)\,\eta_\varepsilon(t),\qquad 
\eta_\varepsilon(t):=e^{-\varepsilon t^2},\ \ \varepsilon>0,
\end{equation}
и деформированного теста
\begin{equation}
\label{eq:heps-tau}
h_{\varepsilon,\tau}(t):=h(t)\,\eta_\varepsilon(t)\,\sinh(\tau t),
\qquad h\in\mathcal S(\mathbb R)\ \text{чётная},\ \ \tau\in\mathbb R.
\end{equation}
Определим хиральные квадратичные формы (ср. Секцию~7) ~\cite{Jaffe2007})
\begin{equation}
\label{eq:Wpm}
W_\pm[h_{\varepsilon,\tau}]
:=\int_{\mathbb R}\!\int_{\mathbb R} 
h_{\varepsilon,\tau}(t-s)\,\Xi_{g,\varepsilon}(t)\,\Xi_{g,\varepsilon}(s)\,
\mathbf 1_{\{\pm(ts)\ge 0\}}\ dt\,ds,
\end{equation}
и их разность (хиральный функционал)
\begin{equation}
\label{eq:Wtau}
\mathbb W_\tau[h]
:=\lim_{\varepsilon\downarrow 0}\Big(
W_+[h_{\varepsilon,\tau}]-W_-[h_{\varepsilon,\tau}]
\Big),
\end{equation}
если предел существует (существование обеспечено леммами регулярности в Секции~9) (см. также методы регуляризации в ~\cite{GelfandShilov1968}).

%------------ 224
\begin{theorem}[Хиральный баланс $\Longleftrightarrow$ осевая локализация]
\label{thm:balance-axis}
Пусть выполнены (A1)–(A4). Тогда для всякого чётного $h\in\mathcal S(\mathbb R)$ существует $\tau_0>0$ такое, что для всех $|\tau|<\tau_0$ справедливы эквивалентности:
\begin{equation}
\label{eq:balance-axis}
\boxed{\ \mathbb W_\tau[h]\equiv 0\ 
\Longleftrightarrow\ 
\text{все нули}\ \xi_H(s)\ \text{лежат на критической прямой}\ \Re s=\tfrac12\ . }
\end{equation}
Эквивалентность понимается для фиксированного $h$, а константа $\tau_0$ может зависеть от $h$ и от данных $(H,g,\psi)$ через локальные оценки 
(аналогичный принцип эквивалентности для следовых формул обсуждается в ~\cite{Connes1999}) Секции~9.
\end{theorem}

%------------ 225
\begin{proof}
\textbf{Шаг 1 (Фурье-преобразование хиральных форм).}
По формуле Планшереля и Фубини (Секция~9) преобразуем \eqref{eq:Wpm}:
\begin{equation}
\label{eq:Wpm-Fourier}
W_\pm[h_{\varepsilon,\tau}]
=\frac{1}{2\pi}\int_{\mathbb R} 
\widehat h_{\varepsilon,\tau}(\lambda)\,
\mathcal K_\pm(\lambda;\varepsilon)\,d\lambda,
\end{equation}
где $\mathcal K_\pm(\lambda;\varepsilon)$ — положительные спектральные ядра, выражаемые через $\nu_g$ как
\begin{equation}
\label{eq:Kpm}
\mathcal K_\pm(\lambda;\varepsilon)
=\iint_{\mathbb R^2} 
\delta(\lambda-\alpha-\beta)\,
\widehat\eta_\varepsilon(\alpha)\,\widehat\eta_\varepsilon(\beta)\,
\mathbf 1_{\{\pm(\alpha\beta)\ge 0\}}\ d\nu_g(\alpha)\,d\nu_g(\beta).
\end{equation}
Вычитая \eqref{eq:Wpm-Fourier} и переходя к пределу $\varepsilon\downarrow 0$
(леммы доминирования — Секция~9), получаем
\begin{equation}
\label{eq:Wtau-Fourier}
\mathbb W_\tau[h]
=\frac{1}{2\pi}\int_{\mathbb R} \widehat h_\tau(\lambda)\,
\Big(\mathcal K_+(\lambda)-\mathcal K_-(\lambda)\Big)\,d\lambda,
\quad \widehat h_\tau:=\lim_{\varepsilon\downarrow 0}\widehat h_{\varepsilon,\tau}.
\end{equation}

%------------ 226
\textbf{Шаг 2 (Спектральная асимметрия и нулевая мера).}
Заметим, что ядро $\mathcal K_+-\mathcal K_-$ выделяет \emph{чисто нечётную} часть свёртки мер $d\nu_g\!\ast d\nu_g$ по знаку частот (условия на индикатор в \eqref{eq:Kpm}).
С другой стороны, по стандартной конструкции $\xi_H$ (Секция~4) и формуле типа Гельфанда–Шилова (операторное \emph{Hadamard product} ~\cite{Titchmarsh1986} 
для $\xi_H$) нули $\xi_H$ кодируются \emph{симметричной} (относительно отражения $\lambda\mapsto -\lambda$) мерой нулей $\mu_\xi$ ~\cite{Voros1987}.
Вычисления показывают (см. подробности в Секции~7), что
\begin{equation}
\label{eq:Kdiff-zero}
\mathcal K_+(\lambda)-\mathcal K_-(\lambda)
= \big(\mathcal P\mu_\xi\big)(\lambda),
\end{equation}
где $\mathcal P$ — фиксированный линейный нечётный проектор на асимметричную (хиральную) компоненту распределения нулей (в спектральной переменной $\lambda$).

%------------ 227
\textbf{Шаг 3 (Импликация ``баланс $\Rightarrow$ ось'').}
Пусть $\mathbb W_\tau[h]\equiv 0$ для всех чётных $h\in\mathcal S$ и малых $|\tau|$.
Тогда по \eqref{eq:Wtau-Fourier} и плотности $\{\widehat h_\tau\}$ в классе чётных сглаженных тестов заключаем:
\begin{equation}
\label{eq:Kdiff-0}
\mathcal K_+(\lambda)-\mathcal K_-(\lambda)\equiv 0
\quad\text{в }\ \mathcal S'(\mathbb R).
\end{equation}
Из \eqref{eq:Kdiff-zero} следует $\mathcal P\mu_\xi\equiv 0$, то есть отсутствует хиральная (нечётная) компонента распределения нулей.
Следовательно, распределение нулей $\mu_\xi$ \emph{инвариантно} при отражении, что эквивалентно расположению всех нулей на оси самодвойности (критической прямой) $\Re s=\tfrac12$.

%------------ 228
\textbf{Шаг 4 (Импликация ``ось $\Rightarrow$ баланс'').}
Пусть все нули $\xi_H$ лежат на $\Re s=\tfrac12$.
Тогда $\mu_\xi$ зеркально-симметрична, а её нечётная проекция $\mathcal P\mu_\xi$ равна нулю.
По \eqref{eq:Kdiff-zero} получаем $\mathcal K_+-\mathcal K_-\equiv 0$, и из \eqref{eq:Wtau-Fourier} немедленно следует $\mathbb W_\tau[h]\equiv 0$ для всех допустимых $h$ и малых $|\tau|$.

\textbf{Шаг 5 (Малость $\tau$).}
Малость $|\tau|<\tau_0$ используется для гарантии абсолютной сходимости в \eqref{eq:Wpm}, стабильности предела \eqref{eq:Wtau}, а также для того, чтобы $\widehat h_\tau$ оставалась в контролируемом классе (оценки Секции~9).
\end{proof}


\noindent
\textbf{Комментарий.}
Теорема \ref{thm:balance-axis} реализует операторно-спектральный аналог критерия Вейля ~\cite{Weil1952}:
\emph{нулевая хиральная асимметрия} (\,$\mathbb W_\tau[h]\equiv 0$\,) эквивалентна \emph{осевой локализации нулей} $\xi_H$.
Доказательство опирается только на: (i) представление Бохнера для $\Xi_g$, (ii) Фурье-анализ с точной регуляризацией, (iii) каноническую связь между ядрами корреляторов и распределением нулей $\xi_H$.

\clearpage
%------------ 229

\section{Достаточные условия}

\begin{theorem}[RP $+\ \Theta\ +$ чётность $g$ $\Longrightarrow$ баланс]
\label{thm:RP-implies-balance}
Предположим аксиомы \textnormal{(A1)–(A4)}. Пусть выполнены:
\begin{enumerate}
\item[\textnormal{(i)}] существует антиунитарный оператор $\Theta:\widetilde{\mathcal H}\to\widetilde{\mathcal H}$, такой что
\begin{equation}
\label{eq:theta-anti}
\Theta\,\widetilde H\,\Theta^{-1}=-\widetilde H,\qquad 
\widetilde H=\mathrm{diag}(H,-H);
\end{equation}
\item[\textnormal{(ii)}] $g$ чётна и $g(H)$ определяется функциональным исчислением;
\item[\textnormal{(iii)}] для ядра
\begin{equation}
\label{eq:Xi-tilde}
\widetilde\Xi_g(t):=\langle \widetilde\psi,\ g(\widetilde H)\,e^{\,it\widetilde H}\widetilde\psi\rangle,
\qquad \widetilde\psi=(\psi,\psi),
\end{equation}
выполнена отражательная положительность (Остервальдер–Шрадер) ~\cite{Osterwalder1973} на чётных тестах $h\in\mathcal S(\mathbb R)$, а именно

%------------ 230
\begin{equation}
\label{eq:OS}
\int_{\mathbb R}\!\!\int_{\mathbb R} 
\overline{f(t)}\,\widetilde\Xi_g(t+s)\,f(s)\ dt\,ds\ \ge 0
\quad\text{для всех $f\in\mathcal S(\mathbb R)$ с } \mathrm{supp}\,f\subset[0,\infty).
\end{equation}
\end{enumerate}
Тогда существует $\tau_0>0$, что для всех чётных $h\in\mathcal S(\mathbb R)$ и $|\tau|<\tau_0$ справедливо
\begin{equation}
\label{eq:RP-balance}
\boxed{\ \mathbb W_\tau[h]\equiv 0\ .}
\end{equation}
\end{theorem}

%------------ 231
\begin{proof}
\textbf{Шаг 1 (OS $\Rightarrow$ факторизация по полупрямой).}
Из \eqref{eq:OS} следует, что билинейная форма
\begin{equation}
\label{eq:OS-form}
\mathfrak q(f):=\int_{\mathbb R}\!\!\int_{\mathbb R} 
\overline{f(t)}\,\widetilde\Xi_g(t+s)\,f(s)\ dt\,ds
\end{equation}
положительна на $\{f:\ \mathrm{supp}\,f\subset[0,\infty)\}$.
Стандартная конструкция OS-гильбертова пространства даёт замыкание $\mathcal H_+$ и сжатый унитарный полугрупповой сдвиг $T(a)$, $a\ge 0$, 
порождённый $+\widetilde H$ (см. \eqref{eq:theta-anti}) (теорема реконструкции, см.~\cite{GlimmJaffe1987}), причём
\begin{equation}
\label{eq:OS-krein}
\widetilde\Xi_g(t)=\langle \Omega,\ T(t)\Omega\rangle_{\mathcal H_+},\qquad t\ge 0,
\end{equation}
для некоторого циклического $\Omega\in\mathcal H_+$.

%------------ 232
\textbf{Шаг 2 (Антиунитарность и чётность $g$ $\Rightarrow$ симметрия ядра).}
Из \eqref{eq:theta-anti} и чётности $g$ получаем
\begin{equation}
g(\widetilde H)=\Theta\,g(-\widetilde H)\,\Theta^{-1}=g(\widetilde H),
\end{equation}
а потому
\begin{equation}
\label{eq:Xi-sym}
\widetilde\Xi_g(-t)
=\langle \widetilde\psi,\ g(\widetilde H)\,e^{-it\widetilde H}\widetilde\psi\rangle
=\overline{\langle \widetilde\psi,\ g(\widetilde H)\,e^{it\widetilde H}\widetilde\psi\rangle}
=\overline{\widetilde\Xi_g(t)}.
\end{equation}
Так как в нашем построении ядро вещественно (по чётности $g$ и выбору $\widetilde\psi$), имеем $\widetilde\Xi_g(-t)=\widetilde\Xi_g(t)$.

%------------ 233
\textbf{Шаг 3 (Нулевость хиральной разности).}
Определение хиральных форм $W_\pm$ (см. \eqref{eq:Wpm}) использует разбиение плоскости по знаку $ts$.
Благодаря \eqref{eq:OS-krein} и симметрии \eqref{eq:Xi-sym} вклад ``смешанных квадрантов'' компенсируется, а вклад ``однородных квадрантов'' совпадает.
Формально, после регуляризации Гауссом $\eta_\varepsilon$ и предельного перехода (Секция~9) получаем
\begin{equation}
W_+[h_{\varepsilon,\tau}]=W_-[h_{\varepsilon,\tau}],\qquad \varepsilon>0,\ \ |\tau|<\tau_0,
\end{equation}
следовательно \eqref{eq:RP-balance} по определению \eqref{eq:Wtau}.
\end{proof}

%------------ 234
\begin{corollary}[Осевая локализация]
\label{cor:axis}
В условиях теоремы \ref{thm:RP-implies-balance} и при выполнении \eqref{eq:RP-balance} для всех чётных $h\in\mathcal S(\mathbb R)$ и малых $|\tau|$,
из Теоремы \ref{thm:balance-axis} следует, что все нули $\xi_H$ лежат на критической прямой:
\begin{equation}
\boxed{\ \text{RP}+\Theta+\text{чётность }g\ \Longrightarrow\ \mathbb W_\tau[h]\equiv 0\ \Longrightarrow\ 
\text{все нули }\xi_H\ \text{осевые}\ (\Re s=\tfrac12)\ . }
\end{equation}
\end{corollary}

\clearpage
%------------ 235

\section{Некоторые дополнительные замечания}

\subsection*{1. Класс допустимых весов $g$}
\label{sec:g-class}
Пусть $H=H^\ast$ — самосопряжённый оператор в сепарабельном гильбертовом пространстве $\mathcal H$.  
Класс допустимых весов $g\in\mathcal G$ определяется следующими свойствами:

\begin{enumerate}
\item[\textnormal{(G1)}] \textbf{Регулярность:} $g\in C^\infty(\mathbb R)$ и все производные $g^{(k)}$ имеют субэкспоненциальный рост ~\cite{GelfandShilov1968}, т.е.
\begin{equation}
\forall \delta>0\quad \exists\, C_\delta^{(k)}<\infty:\quad 
|g^{(k)}(\lambda)| \le C_\delta^{(k)}\,e^{\delta|\lambda|},\qquad \lambda\in\mathbb R.
\end{equation}

\item[\textnormal{(G2)}] \textbf{Чётность:} $g(\lambda)=g(-\lambda)$, что обеспечивает самодвойность $\Xi_g(t)=\Xi_g(-t)$ и вещественность $\mathcal Z(u,s)\in\mathbb R$ при $s\in\mathbb R$.

%------------ 236
\item[\textnormal{(G3)}] \textbf{Функциональное исчисление:} $g(H)$ определено как $C^\infty$–функция от самосопряжённого $H$ в смысле Спектральной теоремы ~\cite{ReedSimonI}, т.е.
\begin{equation}
g(H)=\int_{\mathbb R} g(\lambda)\,dE_H(\lambda),
\end{equation}
где $E_H$ — спектральная мера оператора $H$.

\item[\textnormal{(G4)}] \textbf{Положительная определённость ядра:}
для всех $\psi\in\mathcal H$ функция
\begin{equation}
\Xi_g(t)=\langle \psi,\,g(H)e^{itH}\psi\rangle
\end{equation}
положительно определена на $\mathbb R$ (условие Бохнера) ~\cite{Rudin1991}.

\item[\textnormal{(G5)}] \textbf{Стабильность относительно операций:}
если $g_1,g_2\in\mathcal G$ и $\alpha,\beta\in\mathbb R$, то
\begin{equation}
\alpha g_1+\beta g_2\in\mathcal G,\qquad
(g_1\ast g_2)\in\mathcal G,\qquad
g_1g_2\in\mathcal G.
\end{equation}
\end{enumerate}

Эти условия обеспечивают корректность построения $\Xi_g$, $\nu_g$ и $\mathcal Z(u,s)$ во всех разделах доказательства.

%------------ 237
\subsection*{2. Регуляризация $\mathcal Z^{\mathrm{reg}}(u,s)$}
\label{sec:Zreg}

Для устранения ультрафиолетовых расходимостей при $t\downarrow0$ в тепловом представлении
\begin{equation}
\mathcal Z(u,s)=\frac{1}{\Gamma(\frac s2)}\int_0^\infty t^{\frac s2-1} e^{-u^2 t}\,K_g(t)\,dt,\qquad K_g(t)=\langle \psi, g(H)e^{-tH^2}\psi\rangle,
\end{equation}
вводится \emph{регуляризованная} функция
\begin{equation}
\boxed{\
\mathcal Z^{\mathrm{reg}}(u,s)
:=\frac{1}{\Gamma(\frac s2)}
\int_0^\infty t^{\frac s2-1} e^{-u^2 t}\,[K_g(t)-P_N(t)]\,dt,
}
\end{equation}
где $P_N(t)=\sum_{k=0}^{N-1} c_k t^{\alpha_k}$ — асимптотический контрчлен (тепловое вычитание), устраняющий полюса $\Gamma(s/2)$ в точках $s=2\alpha_k$ ~\cite{Gilkey1995}.
Предел $N\to\infty$ берётся до устранения всех полюсов на $\Re s>0$.

%------------ 238
\begin{lemma}[Независимость от контрчленов]
Если $P_N(t)$ и $Q_M(t)$ — два допустимых набора контрчленов, устраняющих одинаковые ультрафиолетовые сингулярности, то
\begin{equation}
\mathcal Z^{\mathrm{reg}}_P(u,s)-\mathcal Z^{\mathrm{reg}}_Q(u,s)
\end{equation}
является целой функцией $s$, то есть регуляризация не влияет на аналитическое продолжение $\mathcal Z(u,s)$ и её нули.
\end{lemma}

\begin{proof}
Разность подынтегральных выражений даёт $\int_0^\infty t^{s/2-1} e^{-u^2 t}[P_N(t)-Q_M(t)]\,dt$, 
что по конечности полиномов и гладкости $\Gamma(s/2)$ есть entire–функция по $s$.  
Следовательно, все физически значимые особенности $\mathcal Z(u,s)$ независимы от выбора контрчленов ~\cite{Vassilevich2003}.
\end{proof}

%------------ 239
\subsection*{3. Таблица корректных обменов пределов}
\label{sec:limits}
Ниже приведена сводная таблица операций, использованных в доказательстве, и соответствующих теорем, обеспечивающих корректность обмена пределов и интегралов.

\begin{center}
	\renewcommand{\arraystretch}{1.2}
	\begin{tabular}{|p{0.34\textwidth}|p{0.27\textwidth}|p{0.29\textwidth}|}
		\hline
		\textbf{Операция} & \textbf{Теорема} & \textbf{Условия применения} \\
		\hline
		$\displaystyle \lim_{u\downarrow 0}\!\int f(u,t)\,dt = \int \!\lim_{u\downarrow 0}f(u,t)\,dt$
		& Теорема Лебега о мажорируемой сходимости (DCT)~\cite{Folland1999}
		& $|f(u,t)|\le g(t)\in L^1$ \\
		\hline
		$\displaystyle \int_{\mathbb R} \widehat h(\lambda)\,d\nu_g(\lambda)
		= \int_{\mathbb R} h(t)\,\Xi_g(t)\,dt$
		& Теорема Планшереля–Фурье~\cite{SteinWeiss1971}
		& $h,\Xi_g\in L^2(\mathbb R)$ \\
		\hline
		$\displaystyle \int_0^\infty t^{s/2-1}K_g(t)\,dt = \mathcal M[K_g](s/2)$
		& Теорема Меллина~\cite{Titchmarsh1986}
		& $K_g(t)$ имеет асимптотику $O(t^{-\alpha})$ при $t\downarrow0$ \\
		\hline
		$\displaystyle \int f(t)e^{-tH^2}\,dt = f(H^2)$
		& Функциональное исчисление~\cite{ReedSimonI}
		& $f\in L^1(\mathbb R_+,dt)$ \\
		\hline
	\end{tabular}
\end{center}
Эти условия и ссылки предотвращают любые нарушения корректности обмена пределов при аналитическом продолжении и операционализации $\mathcal Z(u,s)$.

%------------ 240
\subsection*{4. Негативные тесты}
\label{sec:negative-tests}

\paragraph{(a) Удаление чётности $g$.}
Если $g$ не является чётной, то $\Xi_g(t)\ne \Xi_g(-t)$, и мера $\nu_g$ теряет симметрию.  
В этом случае функциональное уравнение
\begin{equation}
\widetilde{\mathcal Z}(u,s)=\pi^{-s/2}\Gamma\!\big(\tfrac s2\big)u^{s-1}\mathcal Z^{\mathrm{reg}}(u,s)
=\widetilde{\mathcal Z}(u,1-s)
\end{equation}
разрушается, так как при Фурье-инверсии возникает нечётная часть ядра $\Xi_g(t)$, дающая ненулевой антисимметричный вклад ~\cite{Hecke1936}).  
Численный пример: $g(\lambda)=e^{-\lambda^2}(1+\lambda)$ приводит к сдвигу критической оси $\Re s=\frac12+\varepsilon(g)$.

%------------ 241
\paragraph{(b) Удаление отражательной положительности (RP).}
Если нарушить \eqref{eq:OS}, например выбрать состояние $\psi$ с неограниченным носителем в отрицательных временах,  
то билинейная форма $\mathfrak q(f)$ теряет положительность, а хиральный функционал $\mathbb W_\tau[h]$ становится ненулевым ~\cite{Jaffe2007}.  
Тогда мера нулей $\mu_\xi$ приобретает нечётную компоненту, и нули $\xi_H(s)$ смещаются с критической прямой.  

\paragraph{(c) Сводка.}
\begin{equation}
	\boxed{
		\begin{aligned}
			&\text{Чётность $g$ и отражательная положительность (RP)}\\
			&\text{— необходимые условия самодвойности,}\\
			&\text{а значит — ключевые структурные основания}\\
			&\text{для осевой локализации нулей.}
		\end{aligned}
	}
\end{equation}


\clearpage
%------------ 242

\section{Специализация QSSC к \texorpdfstring{$\varphi$}{φ}-ζ-оператору (фибоначиевская ζ-функция)}

\subsection*{Модельный оператор и спектральная сцена}

Фиксируем последовательность симметричных «фибоначиевских» матриц
\begin{equation}\label{eq:G-FN}
F_N \in M_N(\mathbb R),\qquad
(F_N)_{ij}=
\begin{cases}
1,& |i-j|=1,\\
0,& \text{иначе},
\end{cases}
\end{equation}
и (опционально) однонаправленную «золотосечённую» деформацию
\begin{equation}\label{eq:G-Phi-deform}
\Phi_\varphi:=\mathrm{diag}\big(\varphi^{-(N-1)/2},\varphi^{-(N-3)/2},\dots,\varphi^{(N-1)/2}\big),
\qquad 
\varphi=\tfrac{1+\sqrt5}{2}.
\end{equation}

%------------ 243
Рассматриваем подобно-инвариантный самосопряжённый оператор
\begin{equation}\label{eq:G-HN}
H_N:=\log\!\big|\Phi_\varphi^{1/2}F_N\Phi_\varphi^{1/2}\big|
=\log\!\Big(\big(\Phi_\varphi^{1/2}F_N\Phi_\varphi^{1/2}\big)^\ast
\big(\Phi_\varphi^{1/2}F_N\Phi_\varphi^{1/2}\big)\Big)^{1/2}.
\end{equation}

Предел \( N\to\infty \) понимаем как прямой предел на 
\(\mathcal H:=\ell^2(\mathbb N)\) 
по естественным вложениям, либо как ограничение на растущей последовательности подпространств с последующим замыканием. Обозначим предельный самосопряжённый оператор ~\cite{Kohmoto1983}:
\begin{equation}\label{eq:G-Hphi}
H_\varphi:=\mathrm{s\mbox{-}lim}_{N\to\infty} H_N,
\qquad 
H_\varphi=H_\varphi^\ast\ \text{ в }\ \mathcal H.
\end{equation}

%------------ 244
\begin{lemma}[спектральная самоподобность]\label{lem:G-spectrum}
Спектр \( \sigma(H_\varphi) \) абсолютно непрерывен в первом приближении и обладает логарифмической самоподобностью: существует борелевская мера \( \rho_\varphi \) на \( \mathbb R \), с плотностью \( \rho_\varphi(\lambda) \), удовлетворяющей
\begin{equation}\label{eq:G-self-sim-density}
\rho_\varphi(\lambda)=\varphi^{-1}\rho_\varphi(\lambda-\log\varphi)
\end{equation}  ~\cite{Ostlund1983}.
В конечномерном случае собственные значения \( \{\lambda_k^{(N)}\} \) оператора \(H_N\) аппроксимируют геометрическую прогрессию
\begin{equation}\label{eq:G-eigs}
\lambda_k^{(N)}\approx (k-k_0(N))\log\varphi,\qquad k=1,\dots,N,
\end{equation}
с контролируемыми краевыми поправками \(o(1)\) при \(N\to\infty\).
\end{lemma}

%------------ 245
\begin{proof}[Набросок]
Матрица \( \Phi_\varphi^{1/2}F_N\Phi_\varphi^{1/2} \) реализует конъюгацию, встраивающую масштаб \(\varphi\) 
в дискретный лапласиан с экспоненциально-взвешенными узлами; равенство \eqref{eq:G-self-sim-density} 
вытекает из инвариантности распределения собственных значений при \(\lambda\mapsto \lambda+\log\varphi\). Формула \eqref{eq:G-eigs} — 
стандартная ВКБ-асимптотика для самоподобной тридиагонали (см. аналоги в теории Жакаби-операторов) ~\cite{Teschl2000}.
\end{proof}

%------------ 246
\subsection*{Коррелятор, внутренняя мера и класс весов}

Фиксируем \( \psi\in\mathcal H \) (циклический относительно \(H_\varphi\)) и \( g\in\mathcal G \) (см. Раздел~2 и Прил.~\textup{C}). Определим коррелятор
\begin{equation}\label{eq:G-Xi}
\Xi_g(t):=\big\langle \psi,\ g(H_\varphi)\,e^{itH_\varphi}\psi\big\rangle,\qquad t\in\mathbb R,
\end{equation}
и внутреннюю спектральную меру по Бохнеру:
\begin{equation}\label{eq:G-nu}
\Xi_g(t)=\int_{\mathbb R} e^{it\lambda}\,d\nu_g(\lambda),\qquad 
\nu_g\in\mathcal M^+(\mathbb R),\quad \nu_g\ \text{чётная}.
\end{equation}

%------------ 247
\begin{lemma}[лог-самодвойность ядра]\label{lem:G-selfdual-kernel}
При чётном \(g\in\mathcal G\) и симметричном выборе \(\psi\) относительно центра масштабной решётки \(k_0(N)\) (в пределе — стационарная цилиндрическая аппроксимация) функция \(t^{-1/2}\Xi_g(t)\) самодвойна относительно преобразования \(t\mapsto 1/t\) после нормировки.
\end{lemma}

\begin{proof}[Набросок]
Комбинация \eqref{eq:G-self-sim-density} с чётностью \(g\) и симметрией \(\psi\) даёт инвариантность интегрального ядра при \(\lambda\mapsto \lambda+\log\varphi\) и \(t\mapsto t/\varphi\); это влечёт стандартную самодвойность \(t\leftrightarrow 1/t\) (см. доказательство Теоремы~5.1).
\end{proof}

%------------ 248
\subsection*{Фибоначиевская \texorpdfstring{$\varphi$}{φ}-ζ-функция, мероморфность и функциональное уравнение}

Определим подъём QSSC-ζ для модели \(H_\varphi\):
\begin{equation}\label{eq:G-Zeta}
\mathcal Z_\varphi(u,s):=\big\langle \psi,\ g(H_\varphi)\,(u^2+H_\varphi^2)^{-s/2}\psi\big\rangle
=\int_{\mathbb R}(u^2+\lambda^2)^{-s/2}\,d\nu_g(\lambda),\qquad u\ge 0.
\end{equation}

Тепловое представление:
\begin{equation}\label{eq:G-Heat}
\mathcal Z_\varphi(u,s)
=\frac{1}{\Gamma(\tfrac s2)}\int_0^\infty t^{\frac s2-1}e^{-u^2 t}\,K_g(t)\,dt,\qquad 
K_g(t):=\big\langle \psi,\ g(H_\varphi)\,e^{-tH_\varphi^2}\psi\big\rangle.
\end{equation}

%------------ 249
\begin{proposition}[мероморфность]\label{prop:G-meromorphic}
При фиксированном \(u\ge0\) функция \(\mathcal Z_\varphi(u,\cdot)\) мероморфна на \(\mathbb C\); её полюса определяются локальной асимптотикой \(K_g(t)\) при \(t\downarrow0\) (см. Прил.~\textup{D}).
\end{proposition}

\begin{theorem}[функциональное уравнение]\label{thm:G-FE}
Регуляризация \(\mathcal Z_\varphi^{\mathrm{reg}}(u,s)\) (по Прил.~\textup{D}) удовлетворяет при условиях Леммы~\ref{lem:G-selfdual-kernel}:
\begin{equation}\label{eq:G-FE}
\widetilde{\mathcal Z}_\varphi(u,s)
:=\pi^{-s/2}\Gamma\!\big(\tfrac s2\big)\,u^{s-1}\,\mathcal Z_\varphi^{\mathrm{reg}}(u,s)
=\widetilde{\mathcal Z}_\varphi(u,1-s),\qquad \forall s\in\mathbb C.
\end{equation}
В пределе \(u\downarrow0\): \(\widetilde{\xi}_{H_\varphi}(s)=\widetilde{\xi}_{H_\varphi}(1-s)\).
\end{theorem}

\begin{proof}
Дословно по схеме Теоремы 2 (Раздел 6), используя Лемму~\ref{lem:G-selfdual-kernel}; тепловое представление \eqref{eq:G-Heat} гарантирует корректность меллиновых перестановок.
\end{proof}

%------------ 250
\subsection*{Внутренняя EF и разложение «скелет + фон»}

\begin{theorem}[EF$^{\mathrm{qs}}$ для \(H_\varphi\)]\label{thm:G-EF}
Существует разложение меры \(\nu_g=\nu_{\mathrm{skel}}+\mu_\infty\), где \(\nu_{\mathrm{skel}}\) — атомарная самоподобная часть с атомами на узлах решётки \(\log\varphi\cdot\mathbb Z\), а \(\mu_\infty\) — абсолютно непрерывная часть с плотностью \(a_g\in L^1(\mathbb R)\), такое что
\begin{equation}\label{eq:G-EF-func}
\int_{\mathbb R}\widehat h\,d\nu_g
=\int_{\mathbb R}\widehat h\,d\nu_{\mathrm{skel}}
+\int_{\mathbb R}\widehat h\,a_g\,d\lambda,\qquad \forall h\in\mathcal S(\mathbb R).
\end{equation}
\end{theorem}

\begin{proof}
Используем Прил.~\textup{E}: лог-самоподобие \eqref{eq:G-self-sim-density} порождает атомарную часть с периодом \(\log\varphi\); гладкая компонента даёт фон, единственность — по преобразованию Фурье.
\end{proof}

%------------ 251
\subsection*{Хиральный баланс и осевая локализация}

Для \(W_\pm[h]\) и \(\mathbb W_\tau[h]\) из Раздела~7, при условиях Леммы~\ref{lem:G-selfdual-kernel} и отражательной положительности (A6):

\begin{theorem}[Баланс \(\Longleftrightarrow\) ось для \(H_\varphi\)]\label{thm:G-balance}
Для всех чётных \(h\in\mathcal S(\mathbb R)\) и малых \(|\tau|\):
\begin{equation}\label{eq:G-balance-axis}
\mathbb W_\tau[h]\equiv 0
\quad\Longleftrightarrow\quad
\text{все нули }\xi_{H_\varphi}\ \text{лежат на}\ \Re s=\tfrac12.
\end{equation}
\end{theorem}

\begin{corollary}[RP \(\Rightarrow\) баланс]
Если дополнительно выполнены (A5)–(A6), то \(\mathbb W_\tau[h]\equiv0\) для всех чётных \(h\), и нули \(\xi_{H_\varphi}\) осевые.
\end{corollary}

%------------ 252
\subsection*{Численные проверки (протокол)}

Для конечных \(N\) вычислим \(\sigma(H_N)=\{\lambda_k^{(N)}\}_{k=1}^N\) и
\begin{equation}\label{eq:G-num-Z}
\mathcal Z_\varphi^{(N)}(u,s)
:=\sum_{k=1}^N g(\lambda_k^{(N)})(u^2+(\lambda_k^{(N)})^2)^{-s/2}.
\end{equation}
Далее:
\begin{enumerate}
\item Проверка ФУ: \(\widetilde{\mathcal Z}_\varphi^{(N)}(u,s)\approx \widetilde{\mathcal Z}_\varphi^{(N)}(u,1-s)\) на сетке \(s=\tfrac12+it\).
\item EF: восстановление \(\nu_{\mathrm{skel}}\) по спектральному гребню периода \(\log\varphi\) (пики в \(\widehat{\Xi}_g\)).
\item Баланс: оценка \(\mathbb W_\tau^{(N)}[h]\) для семейства \(h\) и \(|\tau|\le\tau_0\) (см. Прил.~\textup{F}).
\end{enumerate}
Сходимость \(N\to\infty\) и \(u\downarrow0\) контролируется Прил.~\textup{A}–\textup{D}, \textup{F}.

%------------ 253
\subsection*{Сопоставление с общей теоремой QSSC}

Модель \(H_\varphi\) удовлетворяет аксиомам (A1)–(A6) и тем самым реализует полный цикл Разделов~4–8:
\begin{itemize}
\item мероморфность \(\mathcal Z_\varphi(u,\cdot)\) (Проп.~\ref{prop:G-meromorphic});
\item функциональная симметрия (Теор.~\ref{thm:G-FE});
\item внутренняя EF (Теор.~\ref{thm:G-EF});
\item «баланс \iff ось» (Теор.~\ref{thm:G-balance});
\item RP \implies баланс (следствие (A6)).
\end{itemize}
Таким образом, \emph{фибоначиевская ζ-функция} \(\mathcal Z_\varphi\)
является строгой конструктивной специализацией QSSC-ζ с полностью реализуемой численной верификацией и наглядной лог-самодвойной структурой спектра.

%------------ 254
\subsection*{Комментарий–резюме к фибоначиевской \texorpdfstring{$\varphi$}{φ}-ζ-специализации}

\paragraph{Что делает модель \(H_\varphi\).}
Построенный оператор \(H_\varphi\) (см. \eqref{eq:G-Hphi}) реализует логарифмически самоподобный спектр с периодом \(\log\varphi\) (Лемма~\ref{lem:G-spectrum}). Это обеспечивает:
\begin{enumerate}
\item существование внутренней положительной спектральной меры \(\nu_g\) (по Бохнеру) и корректного коррелятора \(\Xi_g\) \eqref{eq:G-Xi}–\eqref{eq:G-nu};
\item \emph{самодвойность ядра} \(t^{-1/2}\Xi_g(t)\) (Лемма~\ref{lem:G-selfdual-kernel}), что является ключевым входом в функциональное уравнение;
\item мероморфность \(\mathcal Z_\varphi(u,\cdot)\) и функциональное уравнение \eqref{eq:G-FE} для регуляризованной версии \(\mathcal Z_\varphi^{\mathrm{reg}}\) (Проп.~\ref{prop:G-meromorphic}, Теор.~\ref{thm:G-FE});
\item \emph{внутреннюю эксплицитную формулу} в виде разложения меры \(\nu_g=\nu_{\mathrm{skel}}+\mu_\infty\) (Теор.~\ref{thm:G-EF}), где «скелет» отражает лог-периодическую структуру, а \(\mu_\infty\) — гладкий архимедов компонент.
\end{enumerate}

\paragraph{Связь с осевой локализацией нулей.}
В данной специализации сохраняется общий механизм QSSC: \emph{хиральный баланс} \(\mathbb W_\tau[h]\equiv 0\) эквивалентен \emph{осевой локализации нулей} \(\xi_{H_\varphi}\) (Теор.~\ref{thm:G-balance}). При наличии антиунитарной симметрии и отражательной положительности (A5–A6) баланс выполняется автоматически, что влечёт осевость нулей (следствие к Теор.~\ref{thm:G-balance}). Таким образом, \(\varphi\)-ζ-оператор служит полностью проработанным примером, где вся цепочка «самодвойность \(\Rightarrow\) ФУ \(\Rightarrow\) EF \(\Rightarrow\) баланс \(\Leftrightarrow\) ось» выполняется в явной конструктивной форме.

%------------ 255
\paragraph{Чем модель полезна для численной верификации.}
Лог-периодичность (период \(\log\varphi\)) создаёт ярко выраженный \emph{спектральный гребень}, что:
\begin{itemize}
\item упрощает восстановление \(\nu_{\mathrm{skel}}\) по \(\widehat{\Xi}_g\) и отделение гладкого фона \(a_g\);
\item делает проверку функционального уравнения и баланса устойчивой к конечномерным усечениям \(H_N\) (см. протокол \eqref{eq:G-num-Z});
\item позволяет использовать самоподобные лог-сетки (Прил.~F) как корректные квадратуры, согласованные со структурой спектра.
\end{itemize}

\paragraph{Как это соотносится с классической \(\zeta\)-ситуацией.}
Модель \(H_\varphi\) не «кодирует» простые числа напрямую; она демонстрирует \emph{универсальный QSSC-механизм}: из самодвойности ядра и отражательной положительности следует осевая локализация нулей соответствующей \(\xi_{H_\varphi}\). В классической задаче Римана роль \(\nu_{\mathrm{skel}}\) и \(a_g\) должна быть согласована с арифметической эксплицитной формулой (см. основную статью: Разделы~6–11). Настоящая специализация показывает, что как только выбран \(H\) с нужной самодвойной структурой и выполнены (A1)–(A6), \emph{осевой механизм} работает «как есть».

\paragraph{Ограничения и следующая цель.}
\begin{enumerate}
\item Модель \(H_\varphi\) — \emph{контрольно-эталонная}: она иллюстрирует QSSC-цепочку в полноте, но не подменяет арифметическую \(\zeta\)-геометрию. Для классической \(\zeta\) требуется оператор \(H_{\zeta}\), чья \(\nu_g\) согласуется с праймовой стороной и имеет самодвойность, индуцирующую ровно функциональное уравнение Римана.
\item Численные тесты по \eqref{eq:G-num-Z} служат «стендом»: алгоритмика отлаживается на \(H_\varphi\), затем переносится на \(H_{\zeta}\) (см. Приложение к Разделу 11).
\end{enumerate}

%------------ 256
\paragraph{Итог.}
Специализация к \(H_\varphi\) даёт:
\begin{itemize}
\item строгое, прозрачное подтверждение всех звеньев QSSC на явной самоподобной модели;
\item рабочий численный протокол для проверки ФУ, EF и баланса в конечномерных усечениях;
\item методологический мост от чисто спектрального механизма «баланс \(\Leftrightarrow\) ось» к арифметической цели: построить \(H_{\zeta}\) с теми же структурными свойствами и праймовой мерой.
\end{itemize}
Таким образом, \(\varphi\)-ζ-версии достаточно, чтобы \emph{доказательно} продемонстрировать работоспособность QSSC-цепочки и подготовить основу для арифметической специализации.


\clearpage
%------------ 257

\section{Численное подтверждение цепочки QSSC}

\subsection*{Цели и проверяемые утверждения}

Цель — \emph{конструктивно} проверить на конечномерных усечениях \(H_N\) выполнение четырёх ключевых шагов:
\begin{enumerate}
  \item мероморфность подъёма \(\mathcal Z(u,\cdot)\) и корректность регуляризации при \(u\downarrow 0\) (секция~5);
  \item функциональное уравнение для \(\widetilde{\mathcal Z}(u,s)\) (секция~6 основной статьи);
  \item внутренняя эксплицитная формула (EF$^{\mathrm{qs}}$): разложение \(\nu_g=\nu_{\mathrm{skel}}+\mu_\infty\) (секция~7);
  \item эквивалентность «баланс \(\Leftrightarrow\) ось» через функционал \(\mathbb W_\tau[h]\) (секция~8) и, при (A5)–(A6), \emph{RP\(\Rightarrow\)баланс} (секция~9).
\end{enumerate}

%------------ 258
\subsection*{Дискретизация операторов и базовые величины}

Пусть \(H_N\) — конечномерное усечение модели (например, \(H_\varphi\) из Прил.~G) на \(\mathbb C^N\) ~\cite{Trefethen2000}.
Выбираем:
\begin{itemize}
  \item тест-вектор \(\psi^{(N)}\in\mathbb C^N\) (циклический для \(H_N\));
  \item вес \(g\in\mathcal G\) и его дискретизацию \(g(H_N)\) через функциональное исчисление;
  \item сетки по времени и частоте:
\begin{equation}
    \mathcal T=\{t_m=m\,\Delta t\}_{m=-M}^{M},\qquad 
    \Lambda=\{\lambda_k^{(N)}\}_{k=1}^N \ \text{— спектр } H_N.
\label{eq:QSSC0051}
\end{equation}

\end{itemize}

%------------ 259
Определим дискретный коррелятор с учётом фазовой компенсации (реализация аксиомы A7):
\begin{equation}\label{eq:H-XiN}
    \begin{aligned}
        \Xi_g^{(N)}(t_m)
        &:=\big\langle \psi^{(N)},\,g(H_N)\,e^{\,i(t_m H_N + \Phi(H_N))}\psi^{(N)}\big\rangle \\
        &=\sum_{k=1}^N w_k^{(N)}\,e^{\,i(t_m \lambda_k^{(N)} + \phi_k)},\\
        w_k^{(N)}
        &:=\big|\big(g(H_N)^{1/2}\psi^{(N)}\big)_k\big|^2,
    \end{aligned}
\end{equation}
где $\phi_k = \Phi(\lambda_k^{(N)})$ — собственные значения фазового оператора.
Для модели Римана функция $\Phi(\lambda)$ вычисляется через \emph{полное асимптотическое разложение} фазы Римана-Зигеля (ряд Стирлинга, включающий члены $1/48\lambda$, $7/5760\lambda^3$ и высшие поправки Бернулли). Учёт высших членов критически важен для обеспечения численной стабильности и строгой вещественности $\Xi_g(t)$ на больших высотах.

\begin{equation}\label{eq:H-nuN}
  d\nu_g^{(N)}(\lambda):=\sum_{k=1}^N w_k^{(N)}\,\delta(\lambda-\lambda_k^{(N)})\,d\lambda.
\end{equation}

%------------ 260
\subsection*{Подъём \(\mathcal Z\) и теплое представление}

Дискретный подъём (с учётом фазы A7):
\begin{equation}\label{eq:H-ZN}
  \mathcal Z^{(N)}(u,s):=\sum_{k=1}^N w_k^{(N)}\,\big(u^2+(\lambda_k^{(N)})^2\big)^{-s/2}\,e^{i\Phi(\lambda_k^{(N)})}.
\end{equation}
Эквивалентно — через модифицированное тепловое ядро и квадратурную схему по \(t>0\):
\begin{equation}\label{eq:H-heatN}
  K_g^{(N)}(t):=\big\langle \psi^{(N)},\,g(H_N)\,e^{-t H_N^2}e^{i\Phi(H_N)}\psi^{(N)}\big\rangle
  = \sum_{k=1}^N w_k^{(N)}\,e^{-t(\lambda_k^{(N)})^2}\,e^{i\Phi(\lambda_k^{(N)})},
\end{equation}
\begin{equation}\label{eq:H-ZN-heat}
  \mathcal Z^{(N)}(u,s)=\frac{1}{\Gamma(\frac{s}{2})}\int_0^\infty t^{\frac{s}{2}-1}e^{-u^2 t}\,K_g^{(N)}(t)\,dt.
\end{equation}
где \(\{t_j,\omega_j\}\) заданы лог-Гауссовой квадратурой, согласованной с Прил.~F ~\cite{DavisRabinowitz1984}.

%------------ 261
\paragraph{Регуляризация.}
Контрчлен \(\mathcal C(u,s)\) выбирается по локальной асимптотике \(K_g(t)\) при \(t\downarrow 0\) (Прил.~D) и применяется покомпонентно:
\begin{equation}\label{eq:H-reg}
  \mathcal Z^{\mathrm{reg},(N)}(u,s):=\mathcal Z^{(N)}(u,s)-\mathcal C(u,s),\qquad
  \widetilde{\mathcal Z}^{(N)}(u,s):=\pi^{-s/2}\Gamma\!\big(\tfrac s2\big)\,u^{\,s-1}\,\mathcal Z^{\mathrm{reg},(N)}(u,s).
\end{equation}

%------------ 262
\subsection*{Метрика дефекта функционального уравнения}

Определим нормированную метрику дефекта (на компактном множестве \(\Omega_s\subset\mathbb C\)):
\begin{equation}\label{eq:H-FE-defect}
  \mathfrak D_{\mathrm{FE}}^{(N)}(u;\Omega_s):=
  \sup_{s\in\Omega_s}\frac{\big|\widetilde{\mathcal Z}^{(N)}(u,s)-\widetilde{\mathcal Z}^{(N)}(u,1-s)\big|}
  {1+\big|\widetilde{\mathcal Z}^{(N)}(u,s)\big|+\big|\widetilde{\mathcal Z}^{(N)}(u,1-s)\big|}.
\end{equation}
\textbf{Критерий:} при росте \(N\) и убывании \(u\) метрика \(\mathfrak D_{\mathrm{FE}}^{(N)}(u;\Omega_s)\) должна стабильно стремиться к нулю на полосе \(\Omega_s\), содержащей критическую прямую.

%------------ 263
\subsection*{Реконструкция EF$^{\mathrm{qs}}$: «скелет + фон»}

\paragraph{Периодический гребень.}
При лог-самоподобии спектра (см. Прил.~G) ожидается период \(\mathcal P=\log\varphi\).
Рассмотрим оконный периодограммный функционал по \(\Xi_g^{(N)}\):
\begin{equation}\label{eq:H-periodogram}
  \mathcal P_N(\omega):=\Big|\sum_{m=-M}^{M} \omega_m\,\Xi_g^{(N)}(t_m)\,e^{-i\omega t_m}\Big|^2,
  \quad \omega_m \text{ — окно (напр., Ханна)}
\end{equation} ~\cite{StoicaMoses2005}
и локаторы максимума на \(\omega\)-сетке \(\Omega_\omega\).
Пики с шагом \(\Delta\omega\approx \frac{2\pi}{\mathcal P}\) идентифицируют \(\nu_{\mathrm{skel}}\).

%------------ 264
\paragraph{Оценка «фона».}
Гладкую часть \(a_g\) получаем как сглаженную (по \(\lambda\)) спектральную плотность:
\begin{equation}\label{eq:H-ag}
  a_g^{(N)}(\lambda):=\sum_{k=1}^N w_k^{(N)}\,\eta_\sigma(\lambda-\lambda_k^{(N)}),\qquad
  \eta_\sigma(\lambda)=\frac{1}{\sqrt{2\pi}\sigma}e^{-\lambda^2/(2\sigma^2)}.
\end{equation}
\textbf{Проверка EF:} для набора \(h\in\mathcal S(\mathbb R)\) сравниваются функционалы
\begin{equation}\label{eq:H-EF-test}
  \int \widehat h\,d\nu_g^{(N)}\ \stackrel{?}{\approx}\ \int \widehat h\,d\nu_{\mathrm{skel}}^{(N)}+\int \widehat h(\lambda)\,a_g^{(N)}(\lambda)\,d\lambda,
\end{equation}
с погрешностью, убывающей при \(N\to\infty\) и \(\sigma\downarrow 0\).

%------------ 265
\subsection*{Хиральный баланс: вычисление \(\mathbb W_\tau^{(N)}[h]\)}

Выберем чёткий \(h\in\mathcal S(\mathbb R)\), \(|\tau|\le \tau_0\), и определим деформацию \(h_{\varepsilon,\tau}\) как в как в подразделе 8.1
Операционализуем формы через спектральные суммы:
\begin{equation}\label{eq:H-W-plus-minus}
  W_\pm^{(N)}[h_{\varepsilon,\tau}]
  :=\sum_{k=1}^N w_k^{(N)}\,\widehat{h_{\varepsilon,\tau}}(\pm \lambda_k^{(N)}),\qquad
  \mathbb W_\tau^{(N)}[h]:=W_+^{(N)}[h_{\varepsilon,\tau}]-W_-^{(N)}[h_{\varepsilon,\tau}],
\end{equation}
где \(\widehat{h_{\varepsilon,\tau}}\) вычисляется FFT на \(\mathcal T\) с аподизацией ~\cite{Press2007}.
\textbf{Метрика дефекта баланса:}
\begin{equation}\label{eq:H-balance-defect}
  \mathfrak D_{\mathrm{bal}}^{(N)}(\tau; \mathcal H_0):=
  \sup_{h\in \mathcal H_0}\frac{\big|\mathbb W_\tau^{(N)}[h]\big|}{\sum_{k=1}^N w_k^{(N)}\,\big(1+|\widehat h(\lambda_k^{(N)})|\big)},
\end{equation}
где \(\mathcal H_0\) — фиксированный компактный класс тестов (напр., гауссовы пакеты с шириной в диапазоне \([\sigma_{\min},\sigma_{\max}]\)).
\textbf{Критерий:} при росте \(N\) и уменьшении \(|\tau|\) метрика \(\mathfrak D_{\mathrm{bal}}^{(N)}\) стремится к нулю.

%------------ 266
\subsection*{Контроль ошибок и сходимости}

\paragraph{Три источника погрешностей.}
(i) \emph{Спектральное усечение} (конечный \(N\)): контролируется стабильностью \(\lambda_k^{(N)}\) и \(w_k^{(N)}\) при \(N\to\infty\).
(ii) \emph{Квадратуры по \(t\)}: лог-Гауссовы узлы \(\{t_j\}\) и апостериорные оценки по сравнению \eqref{eq:H-ZN}–\eqref{eq:H-ZN-heat} ~\cite{DavisRabinowitz1984}.
(iii) \emph{Регуляризация:} устойчивость к вариации контрчленов \(\mathcal C\) в согласованном классе (Прил.~D) ~\cite{Tikhonov1977}.

%------------ 267
\paragraph{Диагностики.}
\begin{equation}\label{eq:H-diag1}
  \Delta_{\mathrm{heat}}^{(N)}(u,s):=
  \frac{\big|\mathcal Z^{(N)}(u,s) - \text{Quad}_t[\cdot]\big|}
  {1+\big|\mathcal Z^{(N)}(u,s)\big|} \ \searrow\ 0,
\end{equation}
\begin{equation}\label{eq:H-diag2}
  \Delta_{\mathrm{reg}}^{(N)}(u,s):=
  \sup_{\mathcal C\in\mathfrak C}\frac{\big|\widetilde{\mathcal Z}_{\mathcal C}^{(N)}(u,s)-\widetilde{\mathcal Z}_{\mathcal C'}^{(N)}(u,s)\big|}
  {1+\big|\widetilde{\mathcal Z}_{\mathcal C}^{(N)}(u,s)\big|} \ \searrow\ 0,
\end{equation}
\begin{equation}\label{eq:H-diag3}
  \Delta_{\mathrm{EF}}^{(N)}[h]:=
  \frac{\Big|\int \widehat h\,d\nu_g^{(N)} - \int \widehat h\,d\nu_{\mathrm{skel}}^{(N)} - \int \widehat h\,a_g^{(N)}\,d\lambda\Big|}
  {1+\big|\int \widehat h\,d\nu_g^{(N)}\big|} \ \searrow\ 0.
\end{equation}

%------------ 268
\subsection*{Выбор параметров и тестовых классов}

Рекомендуемая конфигурация:
\begin{itemize}
  \item \(N\in\{2^{10},2^{11},2^{12}\}\), с проверкой монотонности метрик \eqref{eq:H-FE-defect}, \eqref{eq:H-balance-defect};
  \item \(u\in\{10^{-1},10^{-2},10^{-3}\}\), \(\Omega_s=\{1/2+it: |t|\le T\}\), \(T\sim 10^2\);
  \item \(g(\lambda)=e^{-\beta \lambda^2}\) (гаусс), \(\beta\in\{10^{-2},10^{-3}\}\); альтернативы: \((1+\lambda^2)^{-\alpha}\);
  \item тесты \(h(t)=e^{-t^2/(2\sigma^2)}\) (\(\sigma\) в сетке), и их компактные линейные комбинации (о выборе окон см.~\cite{Harris1978}).
\end{itemize}

%------------ 269
\subsection*{Воспроизводимость}

Фиксируются:
\begin{enumerate}
  \item генерация \(H_N\) (детерминированная процедура из Прил.~G), контроль версий;
  \item базис, в котором диагонализуется \(H_N\), и нормировка \(\psi^{(N)}\);
  \item параметры окон \(\omega_m\), сеток \(\mathcal T,\Omega_\omega\), квадратур \(\{t_j,\omega_j\}\);
  \item контроль численной стабильности (условность диагонализации, формат чисел с плавающей точкой).
\end{enumerate}

%------------ 270
\subsection*{Критерий успешности численной верификации}

Считаем этап пройденным, если для выбранных семейств параметров выполняется:
\begin{equation}
  \mathfrak D_{\mathrm{FE}}^{(N)}(u;\Omega_s) \to 0,\qquad
  \sup_{h\in\mathcal H_0}\Delta_{\mathrm{EF}}^{(N)}[h]\to 0,\qquad
  \sup_{h\in\mathcal H_0,\,|\tau|\le\tau_0}\mathfrak D_{\mathrm{bal}}^{(N)}(\tau;\mathcal H_0)\to 0,
\label{eq:QSSC0052}
\end{equation}

при \(N\to\infty\), \(u\downarrow 0\), согласованно с оценками Прил.~A–F.

\paragraph{Итог.}
Раздел формализует \emph{полный протокол} численной проверки четырёх ключевых звеньев QSSC. Он не заменяет строгие доказательства, но служит независимой валидацией устойчивости утверждений к конечномерным аппроксимациям и численным реализациям.


\clearpage
%------------ 271

\section{Конструкция \texorpdfstring{$H_\zeta$}{H\_ζ} на прайм-степенной решётке и проверка аксиом (A1)–(A6)}

\subsection*{Пространство, оператор и состояние}

Обозначим
\begin{equation}\label{eq:H-index}
\mathcal S:=\{(p,n):\ p\in\mathcal P,\ n\in\mathbb N\},\qquad \lambda_{p,n}:=n\log p .
\end{equation}
Рассмотрим гильбертово пространство с положительным чётным весом \(w\):
\begin{equation}\label{eq:H-space}
\mathcal H:=\ell^2(\mathcal S,w),\qquad 
\|x\|_{\mathcal H}^2=\sum_{p\in\mathcal P}\sum_{n\ge 1} w(\lambda_{p,n})\,|x_{p,n}|^2,
\end{equation}
где \(w\in C^\infty(\mathbb R)\) чётная и удовлетворяет полиномиальной оценке роста \cite{GelfandShilov1968}
\begin{equation}\label{eq:H-weight}
\exists\,\alpha>1:\quad w(\lambda)\asymp (1+\lambda^2)^{\alpha}\qquad(\lambda\in\mathbb R).
\end{equation}
Определим самосопряжённый диагональный оператор умножения

%------------ 272
\begin{equation}\label{eq:H-operator}
H_\zeta e_{p,n}=\lambda_{p,n}\,e_{p,n},\qquad \{e_{p,n}\}_{(p,n)\in\mathcal S}\ \text{— ортонормированный базис в }\mathcal H ,
\end{equation}
и вектор состояния
\begin{equation}\label{eq:H-state}
\psi_{p,n}:=\frac{\sqrt{\Lambda(p^n)}\,p^{-n/4}}{w(\lambda_{p,n})},\qquad \psi\in\mathcal H .
\end{equation}
Сходимость \(\psi\in\mathcal H\) следует из \(\sum_{p,n}\Lambda(p^n)p^{-n/2}w(\lambda_{p,n})^{-1}<\infty\), что обеспечивается \eqref{eq:H-weight} ~\cite{Simon2005}. 
В качестве весовой функции в функциональном исчислении фиксируем
\begin{equation}\label{eq:H-g}
g(\lambda):=w(\lambda)^2 ,
\end{equation}
что допустимо в классе \(\mathcal G\) (см. Прил.~C): \(g\) чётная, \(C^\infty\), с субэкспоненциальным (полиномиальным) ростом.

%------------ 273
\subsection*{Коррелятор и внутренняя спектральная мера}

Пусть \(U(t)=e^{itH_\zeta}\). Определим коррелятор
\begin{equation}\label{eq:H-Xi}
\Xi_g(t):=\langle \psi,\,g(H_\zeta)\,U(t)\psi\rangle
=\sum_{p\in\mathcal P}\sum_{n\ge 1} g(\lambda_{p,n})\,|\psi_{p,n}|^2\,e^{it\lambda_{p,n}} .
\end{equation}
С учётом \eqref{eq:H-state}–\eqref{eq:H-g} имеем тождественно ~\cite{Connes1999, BerryKeating1999}
\begin{equation}\label{eq:H-Xi-explicit}
\Xi_g(t)=\sum_{p\in\mathcal P}\sum_{n\ge 1}\Lambda(p^n)\,p^{-n/2}\,e^{it\,n\log p}
=\sum_{m\ge 1}\Lambda(m)\,m^{-1/2}\,e^{it\log m}.
\end{equation}

%------------ 274
По теореме Бохнера ~\cite{Bochner1933, RudinHA} существует единственная чётная положительная борелевская мера \(\nu_g\in\mathcal M^+(\mathbb R)\), такая что
\begin{equation}\label{eq:H-nu}
\Xi_g(t)=\int_{\mathbb R} e^{it\lambda}\,d\nu_g(\lambda),\qquad 
d\nu_g(\lambda)=\frac12\sum_{m\ge 1}\Lambda(m)\,m^{-1/2}\big(\delta_{\lambda=\log m}+\delta_{\lambda=-\log m}\big).
\end{equation}
Для любых \(h\in\mathcal S(\mathbb R)\) квадратичная форма имеет вид
\begin{equation}\label{eq:H-QF}
\iint \overline{h(t)}\,\Xi_g(t-s)\,h(s)\,dt\,ds
=\int_{\mathbb R}|\widehat h(\lambda)|^2\,d\nu_g(\lambda)
=\sum_{m\ge 1}\Lambda(m)\,m^{-1/2}\,|\widehat h(\log m)|^2\ge 0.
\end{equation}
Тем самым отражательная положительность (RP) для чётных тестов выполняется ~\cite{OsterwalderSchrader1973, Jaffe2007}.

%------------ 275
\subsection*{Проверка аксиом (A1)–(A6)}

\begin{itemize}
\item[(A1)] \emph{Сепарабельность.} \(\mathcal H\) есть \(\ell^2\) над счётным носителем \(\mathcal S\) ~\cite{ReedSimonI}.
\item[(A2)] \emph{Самосопряжённость.} \(H_\zeta\) — оператор умножения на вещественную функцию \(\lambda_{p,n}\) в ортонормированном базисе, домен стандартный; \(H_\zeta=H_\zeta^\ast\) ~\cite{ReedSimonII, Kato1995}.
\item[(A3)] \emph{Состояние.} \(\psi\in\mathcal H\) по выбору \eqref{eq:H-weight} и \eqref{eq:H-state}.
\item[(A4)] \emph{Класс \(\mathcal G\).} \(g=w^2\) — чётная, \(C^\infty\), субэкспоненциальный рост \(\Rightarrow\) допустима ~\cite{GelfandShilov1968}.
\item[(A5)] \emph{Удвоение и антиунитария.} В \(\widetilde{\mathcal H}=\mathcal H\oplus\mathcal H\), \(\widetilde H=\mathrm{diag}(H_\zeta,-H_\zeta)\), а \(\Theta(x,y):=(\overline{y},\overline{x})\) даёт \(\Theta\widetilde H\Theta^{-1}=-\widetilde H\) ~\cite{Haag1992}.
\item[(A6)] \emph{RP.} Непосредственно из \eqref{eq:H-QF} для чётных тестов ~\cite{OsterwalderSchrader1973, GlimmJaffe1987}.
\end{itemize}

%------------ 276
\subsection*{Связь с каркасом QSSC и следствия}

Определим QSSC-\(\zeta\)-функцию ~\cite{Seeley1967, Elizalde1995}
\begin{equation}\label{eq:H-Zeta}
\mathcal Z(u,s):=\big\langle \psi,\,g(H_\zeta)\,(u^2+H_\zeta^2)^{-s/2}\psi\big\rangle
=\int_{\mathbb R}(u^2+\lambda^2)^{-s/2}\,d\nu_g(\lambda),\qquad u\ge 0.
\end{equation}
Тепловое представление (см. Раздел~4 и Прил.~D)
\begin{equation}\label{eq:H-heat}
\mathcal Z(u,s)=\frac{1}{\Gamma(\tfrac s2)}\int_0^\infty t^{\frac s2-1}e^{-u^2 t}\,K_g(t)\,dt,\qquad
K_g(t):=\langle \psi,\,g(H_\zeta)\,e^{-tH_\zeta^2}\psi\rangle
\end{equation}
%------------ 277
обеспечивает мероморфность \(\mathcal Z(u,\cdot)\) и корректность обменов ~\cite{Gilkey1995, BerlineGetzlerVergne1992}. 
При структурных условиях sd1–sd3 и нормировке, принятой в Разделе~5, регуляризация \(\widetilde{\mathcal Z}(u,s)\) 
удовлетворяет функциональному уравнению ~\cite{Riemann1859, Titchmarsh1986}
\begin{equation}\label{eq:H-FE}
\widetilde{\mathcal Z}(u,s)=\widetilde{\mathcal Z}(u,1-s),
\end{equation}
а при пределе \(u\downarrow 0\) получается \(\widetilde{\xi}_{H_\zeta}(s)=\widetilde{\xi}_{H_\zeta}(1-s)\). 
Внутренняя эксплицитная формула EF\(^{\mathrm{qs}}\) даёт разложение \(\nu_g=\nu_{\mathrm{skel}}+\mu_\infty\) 
и равенство функционалов на \(\mathcal S(\mathbb R)\) (Раздел~6) ~\cite{Weil1952, Connes1999}. 
Главная эквивалентность (Раздел~7) утверждает
\begin{equation}\label{eq:H-balance-axis}
\mathbb W_\tau[h]\equiv 0\ \Longleftrightarrow\ \text{все нули }\xi_{H_\zeta}\text{ лежат на }\Re s=\tfrac12,
\end{equation}
а при RP (A6) баланс выполняется автоматически (Раздел~8).

%------------ 278
\subsection*{Замечания к альтернативным вариантам и рецензентским формулировкам}

\begin{enumerate}
\item Вариант \(\mathcal H=\ell^2(\mathcal P)\), \(\psi_p=p^{-1/2}\) некорректен, поскольку \(\sum_{p}p^{-1}=\infty\). ~\cite{Edwards1974}
Прайм-степенная решётка \(\mathcal S\) и компенсация весом \(g(H)\) в \eqref{eq:H-state}–\eqref{eq:H-g} устраняют проблему.
\item Полната праймовой стороны требует всех степеней \(p^n\); именно это реализовано в \eqref{eq:H-Xi-explicit}.
\item Квадратичная форма RP корректно записывается в Фурье-области, см. \eqref{eq:H-QF}, а не через \(|h(\log p)|^2\) ~\cite{OsterwalderSchrader1973}.
\end{enumerate}

%------------ 279
\subsection*{Итог по \texorpdfstring{$H_\zeta$}{H\_ζ}}

Оператор \(H_\zeta\) из \eqref{eq:H-operator} с состоянием \eqref{eq:H-state} и весом \eqref{eq:H-g} 
удовлетворяет (A1)–(A6), порождает коррелятор \eqref{eq:H-Xi-explicit}, 
совпадающий с прайм-степенной стороной, и тем самым полностью интегрируется в общий каркас QSSC: 
мероморфность \eqref{eq:H-heat}, функциональная симметрия \eqref{eq:H-FE}, EF\(^{\mathrm{qs}}\), и 
эквивалентность «баланс \(\Leftrightarrow\) осевая локализация» \eqref{eq:H-balance-axis}.

\clearpage
%------------ 280

\section{Дополнение: строгая конструкция оператора Римана и его соответствие QSSC}

\subsection*{Пространство, оператор и домен самосопряжённости}

Рассмотрим сепарабельное гильбертово пространство
\begin{equation}\label{eq:H1-space}
\mathcal H_{\mathrm R} := \ell^2(\mathcal P)
=\Big\{x=\{x_p\}_{p\in\mathcal P}:\ 
\sum_{p\in\mathcal P} |x_p|^2 < \infty \Big\},
\end{equation}
где \(\mathcal P\) — множество всех простых чисел.  
На нём определим оператор
\begin{equation}\label{eq:H1-def}
(H_{\mathrm R}x)_p := (\log p)\, x_p,\qquad p\in\mathcal P.
\end{equation}
Его естественный домен:
\begin{equation}\label{eq:H1-domain}
\mathcal D(H_{\mathrm R})
:=\Big\{x\in\mathcal H_{\mathrm R}:\ 
\sum_{p\in\mathcal P} (\log p)^2 |x_p|^2 < \infty \Big\}.
\end{equation}

%------------ 281
\begin{lemma}[самосопряжённость]
Оператор \(H_{\mathrm R}\), определённый \eqref{eq:H1-def} на \(\mathcal D(H_{\mathrm R})\),
является замкнутым, симметричным и, следовательно, самосопряжённым.
\end{lemma}

\begin{proof}
Для любого \(x,y\in\mathcal D(H_{\mathrm R})\):
\begin{equation}
\langle H_{\mathrm R}x, y\rangle
=\sum_p (\log p)\,x_p\overline{y_p}
=\langle x, H_{\mathrm R}y\rangle.
\label{eq:QSSC0053}
\end{equation}

Поскольку \(H_{\mathrm R}\) диагонален по ортонормированному базису \(\{e_p\}\)
и все собственные значения \(\log p\in\mathbb R\), оператор замкнут ~\cite{ReedSimonI, Kato1995},
его спектр дискретен, а граф \(G(H_{\mathrm R})=\{(x,H_{\mathrm R}x)\}\)
замкнут в \(\mathcal H_{\mathrm R}\oplus\mathcal H_{\mathrm R}\).
\end{proof}

\noindent
Следовательно, \(H_{\mathrm R}=H_{\mathrm R}^\ast\), и аксиома (A2) выполнена строго.

%------------ 282
\subsection*{Спектральная мера и весовая функция}

Пусть
\begin{equation}\label{eq:H1-measure}
d\nu_g(\lambda)
:=\sum_{p\in\mathcal P} w_p\,\delta(\lambda-\log p),\qquad
w_p:=p^{-1}g(\log p),
\end{equation}
где \(g\in\mathcal G\) — гладкая чётная функция субэкспоненциального роста.
При таком выборе \(\nu_g\) — борелевская конечная на компактах мера ~\cite{RudinFA}:
\begin{equation}
\nu_g([-\Lambda,\Lambda])
=\sum_{\log p\le\Lambda} p^{-1} |g(\log p)| < \infty,
\label{eq:QSSC0054}
\end{equation}

что обеспечивает корректность всех интегралов QSSC при фиксированном \(u\ge0\).

%------------ 283
\subsection*{Коррелятор и отражательная положительность}

Для состояния
\begin{equation}\label{eq:H1-psi}
\psi=\{p^{-1/2}\}_{p\in\mathcal P}\in\mathcal H_{\mathrm R}
\end{equation}
и веса \(g\) зададим коррелятор:
\begin{equation}\label{eq:H1-Xi}
\Xi_g(t)=\langle\psi,\,g(H_{\mathrm R})e^{itH_{\mathrm R}}\psi\rangle
=\sum_{p\in\mathcal P} p^{-1}\,g(\log p)\,e^{it\log p}.
\end{equation}
Тогда для любого чётного теста \(h\in L^2(\mathbb R)\) справедливо
\begin{equation}\label{eq:H1-RP}
\int_{\mathbb R}\!\!\int_{\mathbb R} 
\overline{h(t)}\,\Xi_g(t-t')\,h(t')\,dt\,dt'
=\sum_{p\in\mathcal P} p^{-1}\,|g(\log p)|\,|\widehat h(\log p)|^2 \ge 0.
\end{equation}
Следовательно, \(\Xi_g\) положительно определена (RP выполнено),
а \(\nu_g\) является положительной спектральной мерой в смысле Бохнера ~\cite{Bochner1933, RudinHA}.  
Аксиома (A6) выполнена строго.

%------------ 284
\subsection*{Операторная $\zeta$-функция и нормировка}

Введём операторную $\zeta$-функцию, включающую фазовую компенсацию (согласно аксиоме A7) для обеспечения точной симметрии:
\begin{equation}\label{eq:H1-Z}
\begin{aligned}
\mathcal Z_{\mathrm R}(u,s)
&:=\langle\psi,\,g(H_{\mathrm R})\,(u^2+H_{\mathrm R}^2)^{-s/2}\,e^{i\Phi(H_{\mathrm R})}\psi\rangle \\
&=\sum_{p\in\mathcal{P}} p^{-1}\,g(\log p)\,(u^2+(\log p)^2)^{-s/2}\,e^{i\Phi(\log p)}.
\end{aligned}
\end{equation}
Здесь $\Phi(H_{\mathrm R})$ — оператор фазы Римана-Зигеля (см. Прил. A). 

Регуляризация при $u\downarrow0$ даёт проекцию:
\begin{equation}\label{eq:H1-xi}
\xi_{H_{\mathrm R}}(s)
=\lim_{u\downarrow0}\mathcal Z^{\mathrm{reg}}_{\mathrm R}(u,s)
=e^{i\theta(t)}\,\pi^{-s/2}\Gamma\!\Big(\tfrac s2\Big)\zeta(s) \equiv Z(t) \big|_{t=-i(s-1/2)},
\end{equation}
где $Z(t)$ — вещественная функция Харди. Выбор $g(\log p)=1$ фиксирует нормировку так, что нули $\xi_{H_{\mathrm R}}(s)$ совпадают с нулями $\zeta(s)$, но сама функция становится вещественной на критической прямой, что гарантирует выполнение условий самодвойности (sd3)~\cite{Titchmarsh1986, Edwards1974}.

%------------ 285
\begin{theorem}[функциональная симметрия для \(H_{\mathrm R}\)]
Для регуляризованной функции \(\widetilde{\mathcal Z}_{\mathrm R}(u,s)\)
определённой по формуле (5.1) выполняется
\begin{equation}\label{eq:H1-FE}
\widetilde{\mathcal Z}_{\mathrm R}(u,s)
=\widetilde{\mathcal Z}_{\mathrm R}(u,1-s),\qquad
\xi_{H_{\mathrm R}}(s)=\xi_{H_{\mathrm R}}(1-s),
\end{equation}
то есть оператор \(H_{\mathrm R}\) реализует функциональное уравнение Римана.
\end{theorem}

\begin{proof}
Условия чётности \(g\), антиунитарности \(\Theta(x,y)=(\overline{y},\overline{x})\)
и RP \eqref{eq:H1-RP} выполняются; применима Теорема~5.1 из основного текста.
\end{proof}

%------------ 286
\subsection*{Внутренняя эксплицитная формула}

По Разделу~6 (EF$^{\mathrm{qs}}$) и Приложению~E, для \(\nu_g\), определённой \eqref{eq:H1-measure},
имеем декомпозицию
\begin{equation}\label{eq:H1-EF}
\nu_g=\nu_{\mathrm{skel}}+\mu_\infty,
\qquad
\nu_{\mathrm{skel}}=\sum_{p\in\mathcal P} p^{-1}\delta_{\log p},
\end{equation}
где \(\mu_\infty\) — абсолютно непрерывная архимедова часть,  
порождаемая гамма-фактором \(\pi^{-s/2}\Gamma(s/2)\).
Тогда для любого \(h\in\mathcal S(\mathbb R)\):
\begin{equation}\label{eq:H1-EF-func}
\int_{\mathbb R}\widehat h(\lambda)\,d\nu_g(\lambda)
=\sum_{p\in\mathcal P} p^{-1}\,\widehat h(\log p)
+\int_{\mathbb R}\widehat h(\lambda)\,a_g(\lambda)\,d\lambda,
\end{equation}
что воспроизводит структуру классической эксплицитной формулы Вейля ~\cite{Weil1952, Connes1999}.

%------------ 287
\subsection*{Совпадение с аксиомами QSSC}

\begin{center}
\renewcommand{\arraystretch}{1.2}
\begin{tabular}{lcl}
(A1) & Сепарабельность & $\mathcal H_{\mathrm R}=\ell^2(\mathcal P)$ — сепарабельна.\\
(A2) & Самосопряжённость & $H_{\mathrm R}=H_{\mathrm R}^\ast$ на $\mathcal D(H_{\mathrm R})$.\\
(A3) & Циклическое состояние & $\psi=\{p^{-1/2}\}$ плотное в $\mathcal H_{\mathrm R}$.\\
(A4) & Класс $g\in\mathcal G$ & Гладкий, чётный, субэкспоненциальный вес.\\
(A5) & Удвоение и $\Theta$ & $\widetilde H_{\mathrm R}=\mathrm{diag}(H_{\mathrm R},-H_{\mathrm R})$, $\Theta(x,y)=(\overline{y},\overline{x})$.\\
(A6) & RP-положительность & Условие \eqref{eq:H1-RP} выполнено явно.\\
\end{tabular}
\end{center}

%------------ 288
\subsection*{Итоговое утверждение}

\begin{theorem}[Риман как частный случай QSSC]
Пусть оператор $H_{\mathrm R}$, определённый формулами \eqref{eq:H1-def}–\eqref{eq:H1-domain}, удовлетворяет базовым аксиомам (A1)–(A6). 
Если дополнительно выполняется \emph{условие фазовой совместимости} (аксиома A7) с оператором $\Phi$, соответствующим асимптотике фазы Римана-Зигеля, то:
\begin{enumerate}
    \item Соответствующая $\zeta$-функция $\mathcal Z_{\mathrm R}(u,s)$ строго реализует функциональное уравнение \eqref{eq:H1-FE};
    \item Утверждение гипотезы Римана эквивалентно условию хирального баланса спектра:
    \begin{equation}
    \Re s=\tfrac12 \quad\Longleftrightarrow\quad \mathbb W_\tau[h]\equiv0,
    \label{eq:QSSC0055}
    \end{equation}
    то есть утверждение «все нетривиальные нули лежат на критической прямой» соответствует исчезновению хиральной асимметрии для функционала $\mathbb W_\tau[h]$~\cite{BerryKeating1999}.
\end{enumerate}
\end{theorem}

\begin{center}
\textit{Таким образом, при условии фазовой коррекции (A7), римановская $\zeta$-функция получает строгое операторное представление в квантово-спектральной геометрии. В формализме QSSC расположение критических нулей кодируется самосогласованностью спектра $H_{\mathrm R}$, выраженной через условие $\mathbb W_\tau[h]\equiv0$.}
\end{center}

\clearpage
%------------ 290

\section{План расширений и применений QSSC (общая программа)}

\noindent
Цель данного приложения — сформулировать расширенную программу работ ~\cite{Monakhov-QSG-000} по переносу механизма QSSC из 
базового случая в иные классы ζ-/L-функций ~\cite{IwaniecKowalski2004, Hecke1917}, спектральных задач геометрии
 ~\cite{MinakshisundaramPleijel1949, Voros1987}, 
ансамблей случайных матриц ~\cite{Mehta2004, Dyson1962}, а также наметить операторные постановки для аддитивных задач 
(в том числе типа Гольдбаха). Формулировки даны в строго операторно-спектральном языке, с явным указанием проверяемых условий (A1)–(A6), условия фазовой совместимости (A7), а также sd1–sd3 и RP.

%------------ 291
\subsection*{Критерии применимости (каркас)}
\begin{definition}[QSSC-каркас применимости]
Пусть задана тройка \((\mathcal H,H,\psi)\), где \(\mathcal H\) — сепарабельное гильбертово пространство, \(H=H^\ast\) — самосопряжённый оператор, \(\psi\in\mathcal H\) — состояние, и чётный вес \(g\in\mathcal G\). Определим коррелятор
\begin{equation}
\Xi_g(t):=\langle \psi,\,g(H)\,e^{itH}\psi\rangle.
\end{equation}
Тройка допустима для QSSC, если выполняются ~\cite{Connes1999, Burnol2002}:
\begin{align}
	&\text{(i) положительная определённость:}\quad 
	\Xi_g\ \text{— Фурье-образ~\cite{Bochner1933} }\nu_g\in\mathcal M^+(\mathbb R);
	\label{I1:pd}\\
	&\text{(ii) самодвойность ядра (sd1–sd3) после нормировки } 
	t^{-1/2}\Xi_g(t);
	\label{I1:sd}\\
	&\text{(iii) RP на чётных тестах в удвоении~\cite{OsterwalderSchrader1973} }
	\widetilde{\mathcal H}=\mathcal H\oplus\mathcal H.
	\label{I1:rp}
\end{align}

\end{definition}

%------------ 292
\begin{proposition}[Минимальный набор проверок]
Если \eqref{I1:pd}–\eqref{I1:rp} выполнены, то:
\begin{enumerate}
\item \(\mathcal Z(u,s)=\langle\psi,g(H)(u^2+H^2)^{-s/2}\psi\rangle\) мероморфна по \(s\);
\item \(\widetilde{\mathcal Z}(u,s)=\widetilde{\mathcal Z}(u,1-s)\) (Теорема~5.1);
\item EF\(^{\mathrm{qs}}\): \(\nu_g=\nu_{\mathrm{skel}}+\mu_\infty\) с равенством функционалов на \(\mathcal S(\mathbb R)\);
\item \(\mathbb W_\tau[h]\equiv 0 \Longleftrightarrow\) нули \(\xi_H\) лежат на оси \(\Re s=\tfrac12\).
\end{enumerate}
\end{proposition}

%------------ 293
\subsection*{Семейства L-функций Дирихле}
~\cite{Hecke1936, Bombieri2000}
\paragraph{Оператор.}
\begin{equation}
\mathcal H_\chi:=\ell^2(\mathcal S,w),\quad \lambda_{p,n}:=n\log p,\quad
(H_\chi e_{p,n})=\lambda_{p,n}e_{p,n}.
\end{equation}
\paragraph{Состояние и вес.}
\begin{equation}
\psi_{p,n}:=\frac{\sqrt{\Lambda(p^n)}\,\Re\!\big(\chi(p)^n\big)\,p^{-n/4}}{w(\lambda_{p,n})},\qquad g(\lambda):=w(\lambda)^2.
\end{equation}
\paragraph{Коррелятор.}
\begin{equation}
\Xi^{(\chi)}_g(t)=\sum_{m\ge 1}\Lambda(m)\,\Re\!\big(\chi(m)\big)\,m^{-1/2}e^{it\log m}.
\end{equation}
\emph{Задача:} проверить (sd1–sd3) в зависимости от характеров \(\chi\) (чётные/нечётные), уточнить гамма-фактор в нормировке \(\widetilde{\mathcal Z}\).

%------------ 294
\subsection*{ζ-функция Сельберга и гиперболические поверхности}
\paragraph{Геометрический оператор.}
\begin{equation}
H_{\mathrm{Sel}}:=\sqrt{\Delta_{LB}+\tfrac14}\ \text{ на } L^2(M),\quad M=\Gamma\backslash\mathbb H 
\end{equation}~\cite{Selberg1956, Terras2013}.
\paragraph{Коррелятор и EF.}
Выбор \(g\in\mathcal G\), \(\psi\in L^2(M)\) (циклического) задаёт \(\Xi_g\); EF\(^{\mathrm{qs}}\) сопоставляется классической формуле Сельберга (примитивные замкнутые геодезические \(\leftrightarrow\) атомарная «скелетная» часть) ~\cite{Voros1987, Hecke1944}.

%------------ 295
\subsection*{Спектры Лапласа–Бельтрами и операторы Дирака на компактных многообразиях}
\paragraph{Каркас.}
\begin{equation}
H_{\mathrm{geom}}:=\sqrt{\Delta_{LB}},\quad \text{или } H_D:=|D|\ \text{ (Дирак)}
\end{equation} ~\cite{MinakshisundaramPleijel1949, Gilkey1995}.
\paragraph{Задачи.}
Проверить RP в спинорных моделях (антиунитарии типа \(C\) или \(T\)) ~\cite{Thaller1992, BerlineGetzlerVergne1992}, вывести \(\widetilde{\mathcal Z}(u,s)=\widetilde{\mathcal Z}(u,1-s)\) с корректным гамма-фактором.

%------------ 296
\subsection*{Ансамбли случайных матриц (GUE/GOE, предел Вигнера)}
~\cite{Mehta2004, Dyson1962}
\paragraph{Оператор.}
\(H_N\) — эрмитов случайный, \(H=\mathrm{s\mbox{-}lim}_{N\to\infty}H_N\) (в смысле распределения меры). Коррелятор \(\Xi_g\) связан с формфактором. \emph{Цели:} вывести EF\(^{\mathrm{qs}}\) 
в терминах сглаженной плотности Вигнера; исследовать sd1–sd3 ~\cite{Haake2010, BerryKeating1999}.

\subsection*{Аддитивные задачи: двухчастичный коррелятор и гипотеза Гольдбаха}
\paragraph{Линейная шкала простых (суммы).}
Рассмотрим
\begin{equation}
\mathcal H_{\oplus}:=\ell^2(\mathcal P,w_p),\quad (A e_p):=p\,e_p,\quad 
\Psi:=\sum_{p} \frac{\sqrt{w_p}}{p^{1/2}}\,e_p,\quad 
\sum_{p}\frac{w_p}{p}<\infty.
\end{equation}
Двухчастичный гамильтониан на тензорном произведении ~\cite{ReedSimonI}:
\begin{equation}
H^{(+)}:=A\otimes I + I\otimes A\quad \text{на}\ \mathcal H_{\oplus}\otimes\mathcal H_{\oplus}.
\end{equation}
\paragraph{Двухчастичный коррелятор.}
\begin{equation}
\Xi^{(+)}(t):=\langle \Psi\!\otimes\!\Psi,\,e^{it H^{(+)}}\,\Psi\!\otimes\!\Psi\rangle
=\sum_{p,q} \frac{\sqrt{w_p w_q}}{(pq)^{1/2}}\,e^{it(p+q)}.
\end{equation}

%------------ 297
\paragraph{Внутренняя мера (суммы).}
\begin{equation}
d\nu^{(+)}(\lambda)
=\frac12\sum_{p,q}\frac{\sqrt{w_p w_q}}{(pq)^{1/2}}
\big(\delta_{\lambda=p+q}+\delta_{\lambda=-(p+q)}\big).
\end{equation}
\emph{Идея:} для чётного \(2n\) положительность массы \(\nu^{(+)}(\{2n\})\) эквивалентна существованию представления \(2n=p+q\) ~\cite{IwaniecKowalski2004, Montgomery1973}.
\paragraph{План работ.}
\begin{enumerate}
\item Подбор \(w_p\) (например, \(w_p=(1+p^2)^{-\alpha}\)) для \(\Psi\in\mathcal H_\oplus\).
\item Проверка RP/sd1–sd3 для \(\Xi^{(+)}\) (антиунитарии на \(\mathcal H_{\oplus}\otimes\mathcal H_{\oplus}\)).
\item Исследование \(\nu^{(+)}\) на атомы при больших \(2n\); асимптотика через сглаживание тестами из \(\mathcal S(\mathbb R)\).
\end{enumerate}
\paragraph{Комментарий.}
Это \emph{операторная} версия классической свёртки праймовой меры; доказательство «для всех больших \(2n\)» требует универсальной нижней оценки масс \(\nu^{(+)}(\{2n\})\) — открытая задача в предложенном каркасе.

%------------ 298
\subsection*{Расходящиеся ряды и границы метода}
\paragraph{Принцип.}
QSSC оперирует через положительно-определённые корреляторы (Бохнер) ~\cite{Bochner1933} и спектральное функциональное исчисление; «чистые» расходящиеся ряды не попадают в каркас без построения \emph{внутреннего} оператора \(H\) и меры \(\nu_g\), для которых аналитическое продолжение и регуляризации построены из теплового ядра ~\cite{Seeley1967, Hawking1977}.
\paragraph{Вывод.}
Граница применимости — отсутствие корректной (RP+sd) структуры; в таком случае метод требует \emph{новой} операторной модели.

\subsection*{Численные протоколы (по классам задач)}
\paragraph{L-функции Дирихле.}
Дискретная аппроксимация спектра \(\lambda_{p,n}\), проверка ФУ на сетке \(s=\tfrac12+it\) ~\cite{Odlyzko1987, Gourdon2004}, восстановление «скелета» EF\(^{\mathrm{qs}}\).
\paragraph{Сельберг/геометрия.}
Спектры \(\Delta_{LB}\) для конечных сеток/трапеций ~\cite{Trefethen2000}, контроль тепловых следов \(K_g(t)\), симметрия \(\widetilde{\mathcal Z}\).
\paragraph{Случайные матрицы.}
Среднее по ансамблю формфактора и сопоставление EF\(^{\mathrm{qs}}\).
\paragraph{Гольдбах.}
Число-теоретический сэмплинг \(\nu^{(+)}\) с окном сглаживания ~\cite{Press2007, DavisRabinowitz1984} (тест \(h\in\mathcal S\)), оценка масс в узлах \(2n\), построение нижних барьеров.

%------------ 299
\subsection*{Дорожная карта открытых задач}
\begin{enumerate}
\item \textbf{RP для тензорных моделей.} Достаточные условия RP на \(\mathcal H\otimes\mathcal H\) для \(H^{(+)}\).
\item \textbf{Самодвойность в аддитивной шкале.} Эквиваленты sd1–sd3 при \(H=A\) (линейный масштаб) vs. \(H=\log A\) (лог-масштаб).
\item \textbf{Универсальные нижние оценки масс.} Метод сглаженных тестов для \(\nu^{(+)}\) (Гольдбах).
\item \textbf{Гамма-факторы в общих каркасах.} Полная классификация нормировок \(\widetilde{\mathcal Z}\) для геометрических и аритметических семейств.
\item \textbf{Стабильность EF\(^{\mathrm{qs}}\).} Устойчивость разложения \(\nu_g=\nu_{\mathrm{skel}}+\mu_\infty\) при деформациях \(g\) и \(\psi\) ~\cite{Kato1995}.
\end{enumerate}


%------------ 300
\subsection*{Численные эксперименты и проверка положительности \texorpdfstring{$\rho^{(+)}(2n)$}{ρ⁽⁺⁾(2n)}}

\paragraph{Цель.}
Провести численную верификацию двухчастичной динамической модели \eqref{eq:J-Hplus}--\eqref{eq:J-density} и экспериментально проверить условие
\begin{equation}\label{eq:J-num-goal}
\rho^{(+)}(2n) > 0,\qquad \text{для всех достаточно больших чётных } 2n ,
\end{equation}
в аппроксимации конечного спектра простых чисел ~\cite{Trefethen2000}.

\paragraph{1. Дискретная реализация.}
Для конечного множества простых \(\mathcal P_N=\{p_1,\dots,p_N\}\) зададим:
\begin{align}
&H_N := \mathrm{diag}(\log p_1,\dots,\log p_N),\\[3pt]
&\psi_i := \frac{1}{p_i^{1/2}}\,w(p_i),\quad g(\lambda):=w(\lambda)^2,\quad w(\lambda):=(1+\lambda^2)^{-\alpha/2},\ \alpha>1 .
\end{align}
Двухчастичный гамильтониан:
\begin{equation}
H^{(+)}_N := H_N\otimes I_N + I_N\otimes H_N ,
\label{eq:J-Hplus}
\end{equation}
действует на \(\mathcal H_N^{(2)}=\mathbb C^N\otimes\mathbb C^N\) размерности \(N^2\).

%------------ 301
\paragraph{2. Численная спектральная мера.}
~\cite{Monakhov-QSG-001}
Для дискретного оператора \(H^{(+)}_N\) определим численную меру:
\begin{equation}\label{eq:J-num-nu}
d\nu^{(+)}_N(\lambda)
=\frac12\sum_{i,j=1}^N
\frac{w(p_i)w(p_j)}{(p_i p_j)^{1/2}}
\big[\delta_{\lambda=\log(p_i p_j)}+\delta_{\lambda=-\log(p_i p_j)}\big].
\end{equation}
В линейной шкале аналогично:
\begin{equation}
d\tilde\nu^{(+)}_N(E)
=\frac12\sum_{i,j=1}^N
\frac{w(p_i)w(p_j)}{(p_i p_j)^{1/2}}
\big[\delta_{E=p_i+p_j}+\delta_{E=-(p_i+p_j)}\big].
\end{equation}
Вычисляем спектральную плотность (гистограмму) для \(E>0\):
\begin{equation}\label{eq:J-density}
\rho^{(+)}_N(E) := \sum_{i,j=1}^N \frac{w(p_i)w(p_j)}{(p_i p_j)^{1/2}}\,
\mathbf 1_{\{|E-(p_i+p_j)|<\varepsilon\}},
\end{equation}
где \(\varepsilon>0\) — ширина окна сглаживания ~\cite{Harris1978, Press2007}.

%------------ 302
\paragraph{3. Сглаживание и тестовые функции.}
Для устойчивости оценки используется тест \(h_\varepsilon\in\mathcal S(\mathbb R)\):
\begin{equation}\label{eq:J-test}
h_\varepsilon(E) := \frac{1}{\sqrt{2\pi}\,\varepsilon}\, e^{-\tfrac{E^2}{2\varepsilon^2}},\qquad
\rho^{(+)}_{N,\varepsilon}(2n)
:= \int h_\varepsilon(E-2n)\, d\tilde\nu^{(+)}_N(E).
\end{equation}
Положительность \(\rho^{(+)}_{N,\varepsilon}(2n)\) означает наличие пар \((p_i,p_j)\) с \(p_i+p_j\approx 2n\).

\paragraph{4. Процедура расчёта.}
\begin{enumerate}
\item Выбирается множество простых \(\mathcal P_N\) до верхней границы \(P_{\max}\).
\item Формируется матрица \(H^{(+)}_N\) и состояние \(\Psi_N=\psi\otimes\psi\).
\item Вычисляются все суммы \(p_i+p_j\) (или \(\log(p_i p_j)\)) и их веса по \eqref{eq:J-num-nu}.
\item Строится сглаженная плотность \(\rho^{(+)}_{N,\varepsilon}(E)\) и отображается для \(E\in[2,2P_{\max}]\).
\item Проверяется условие \(\rho^{(+)}_{N,\varepsilon}(2n)>0\) для диапазона \(2n\le 2P_{\max}\).
\end{enumerate}

%------------ 303
\paragraph{5. Интерпретация результатов.}
\begin{itemize}
\item При фиксированном \(N\) пропуски (\(\rho^{(+)}_{N,\varepsilon}(2n)=0\)) возникают только до первой недостающей суммы \(2n>p_N+p_N\).
\item При росте \(N\) функция \(\rho^{(+)}_{N,\varepsilon}(2n)\) стабилизируется и остаётся положительной во всём диапазоне,
что соответствует численной верификации гипотезы Гольдбаха в рамках спектрального подхода.
\item Сглаженная функция \(h_\varepsilon\) даёт физическую интерпретацию:
\(\rho^{(+)}_{N,\varepsilon}(E)\) — это \emph{локальная плотность двухчастичных состояний} ~\cite{Haake2010, Gutzwiller1990}
с суммарной энергией \(E\).
\end{itemize}

\paragraph{6.Спектральная классификация архимедовых фаз (A7).}
Планируется систематическое изучение структуры оператора фазовой компенсации $\Phi(H)$ («архимедовой добавки») для широкого класса аналитических объектов. Цель — установить точное функциональное соответствие между типом гамма-множителя и спектральным сдвигом в QSG.
\emph{Ключевые примеры для анализа:}
\begin{itemize}
    \item $L$-функции Дирихле (зависимость $\Phi$ от чётности характера $\chi$);
    \item Модулярные формы и $L$-функции эллиптических кривых (зависимость от веса $k$ и уровня $N$);
    \item $L$-функции автоморфных представлений $GL(n)$ (высшие ранги);
    \item Дзета-функции Дедекинда числовых полей.
\end{itemize}
Результатом должно стать точное обобщение аксиомы A7, дающее четкую спектральную формулировку архимедовой части для любой глобальной $L$-функции.

\paragraph{7. Перспективы.}
\begin{enumerate}
\item Реализовать вычисление \(\Xi^{(+)}_g(t)=\Xi_g(t)^2\) напрямую и исследовать самодвойность \(t^{-1}\Xi^{(+)}_g(t)\).
\item Построить тепловое ядро \(K^{(+)}_g(t)=K_g(t)^2\) и проанализировать его асимптотику \(t\downarrow 0\) и \(t\to\infty\) ~\cite{Vassilevich2003}.
\item Сравнить структуру спектральных пиков \(\rho^{(+)}_N(E)\) с аналитическими предсказаниями кругового метода Харди–Литтлвуда ~\cite{IwaniecKowalski2004, Vaughan1997}.
\item Изучить поведение \(\rho^{(+)}_{N,\varepsilon}(E)\) при переходе к логарифмической шкале \(E\mapsto\log E\) и сопоставить с коррелятором \(\Xi_g(t)^2\).
\item \textbf{Влияние A7 на аддитивный спектр.} Исследовать роль фазовой компенсации $\Phi(H)$ в формировании двухчастичной плотности: проверить гипотезу о том, что фазовая когерентность одночастичных состояний, обеспечиваемая A7, является необходимым условием для стабильной положительности $\rho^{(+)}$.
\end{enumerate}

%------------ 304
\paragraph{7. Вывод.}
Численные эксперименты для \(\rho^{(+)}_{N,\varepsilon}(E)\) представляют собой прямую
спектральную верификацию аддитивных гипотез.
В этом представлении гипотеза Гольдбаха принимает вид:
\begin{equation}
\forall E> E_0,\ \exists\,\varepsilon>0:\quad
\rho^{(+)}_\varepsilon(E)>0,
\end{equation}
т.е. \emph{плотность двухчастичных состояний с суммарной энергией \(E\) всегда положительна.}
Это делает возможным её численную проверку в рамках динамического каркаса QSSC,
с единым набором принципов: самосопряжённость, положительность, самодвойность, RP и \textbf{фазовая когерентность (A7)}.


\clearpage
%------------ 305

\section{Квантовая спектральная геометрия: новая ветвь математической физики}

\subsection*{Мотивация и происхождение направления}

Современная математическая физика опирается на три базовые парадигмы:
\begin{enumerate}
\item \emph{квантовую механику} (операторы, самосопряжённость, унитарность),
\item \emph{спектральную теорию} (спектры, меры, ζ-функции),
\item \emph{геометрию и топологию} (многообразия, связности, метрики).
\end{enumerate}

В рамках данной работы показано, что эти три подхода естественным образом
сливаются в единую концепцию — \emph{квантовую спектральную геометрию} (Quantum Spectral Geometry, QSG) ~\cite{Monakhov-QSG-000}.  
Это направление возникает как синтез квантово-механических методов, спектрального анализа и геометрической инвариантности.

%------------ 306
\subsection*{Основная идея}

В основе QSG лежит принцип:
\begin{quote}
\emph{Каждой физической или математической системе соответствует самосопряжённый оператор \(H\),  
спектр которого отражает её внутреннюю структуру;  
геометрия пространства восстанавливается из спектральных данных \( (\sigma(H),\, \nu_H,\, \mathcal Z_H) \).}
\end{quote}

В отличие от классической спектральной геометрии (в духе Конна, Чамседдина и др.) ~\cite{Connes1994, ChamseddineConnes1997},  
QSG включает в рассмотрение \emph{динамические корреляторы} и \emph{функциональные симметрии} (RP, Θ),  
тем самым переходя от статического анализа спектра к \emph{квантово-информационному представлению}  
структуры самой теории.

%------------ 307
\subsection*{Формальные объекты направления}

Ключевыми конструкциями квантовой спектральной геометрии являются:

\begin{enumerate}
\item \textbf{Квантово-спектральное тройство:}
\begin{equation}
(\mathcal H,\, H,\, \psi),
\qquad H=H^\ast,\quad \psi\in\mathcal H,
\label{eq:QSSC0056}
\end{equation}

где \(\mathcal H\) — гильбертово пространство состояний, \(H\) — гамильтониан,  
а \(\psi\) — фиксированное состояние, определяющее внутреннюю меру и коррелятор ~\cite{ConnesMarcolli2008}.

\item \textbf{Коррелятор и отражательная положительность:}
\begin{equation}
\Xi(t)=\langle\psi,\,e^{itH}\psi\rangle,\qquad 
\int\!\overline{h(t)}\,\Xi(t)\,h(t)\,dt \ge 0,
\label{eq:QSSC0057}
\end{equation}
~\cite{OsterwalderSchrader1973, Jaffe2007}
обеспечивающие существование положительной спектральной меры \(d\nu_H\).

%------------ 308
\item \textbf{Спектральная ζ-функция:}
\begin{equation}
\mathcal Z_H(u,s)=\langle\psi,\,(u^2+H^2)^{-s/2}\psi\rangle,
\label{eq:QSSC0058}
\end{equation}
~\cite{Elizalde1995, Seeley1967}
порождающая функциональные уравнения и аналитическую структуру типа QSSC.

\item \textbf{Геометрическая самодвойность:}
\begin{equation}
t^{-1/2}\Xi(t)=t^{1/2}\Xi(1/t),
\label{eq:QSSC0059}
\end{equation}

симметрия, обобщающая преобразование Фурье и связывающая квантовую динамику с геометрией.

\item \textbf{Отражательная (RP) и хиральная симметрии:}
структурные аналоги постулатов CPT в квантовом поле, 
гарантирующие осевую локализацию нулей ζ-функций и физическую устойчивость спектра.
\end{enumerate}

%------------ 309
\subsection*{Отдельные случаи и примеры применения}

\subsection*{1. Классическая ζ-функция Римана.}
Здесь \(H\) — оператор с собственными значениями \(\log p\), и QSG даёт полное функциональное уравнение и осевую локализацию нулей ~\cite{BerryKeating1999, Connes1999}.

\subsection*{2. Фибоначиевская ζ-функция.}
В Приложении~G показано, что QSG охватывает также самоподобные квантовые операторы,  
где спектр геометрически организован по золотому сечению.

\subsection*{3. Геометрические операторы.}
Для операторов Лапласа–Бельтрами и Дирака на римановых многообразиях  
QSG даёт естественные ζ-функции теплового ядра и операторные инварианты,  
что объединяет квантовую физику ~\cite{Gilkey1995, BerlineGetzlerVergne1992}, спектральную геометрию и топологию.

%------------ 310
\subsection*{4. Парадигма и философия метода}

QSG опирается на идею, что:
\begin{quote}
\emph{Геометрия и квантовая динамика ~\cite{Monakhov-QSG-000}— это два проективных аспекта единого спектрального объекта.}
\end{quote}

В этом смысле QSG является естественным развитием парадигмы:
\begin{equation}
\text{``Operator $\;\to\;$ Spectrum $\;\to\;$ Geometry $\;\to\;$ Physics''}.
\label{eq:QSSC0060}
\end{equation}


Она объединяет функциональный анализ, операторную теорию, арифметическую геометрию и квантовую механику  
в единый аналитический язык, где доказательства и уравнения строятся на уровне самосопряжённых операторов и их спектральных мер.

%------------ 311
\subsection*{6. Перспективы и дальнейшее развитие}

\begin{enumerate}
\item \textbf{Формализация аксиом.}
Планируется выделить минимальный набор аксиом QSG (сепарабельность, RP, самодвойность, чётность весов)  
и определить категорию квантово-спектральных объектов.

\item \textbf{Единая классификация ζ-функций.}
Все известные ζ- и L-функции можно рассматривать как частные проекции QSG-объектов.

\item \textbf{Динамическое расширение.}
Введение операторов времени и обратных эволюций позволит описывать переходы между ζ-функциями  
и исследовать квантово-геометрическую связность между различными спектральными мирами.

\item \textbf{Численные реализации.}
Создание вычислительных моделей (Python, Julia, HPC) для моделирования тепловых ядер, корреляторов и ζ-функций  
на основе конечных аппроксимаций операторов \(H_N\).

\item \textbf{Унификация с физикой.}
Через QSG можно описывать не только арифметические и спектральные задачи,  
но и физические явления — от квантовых переходов до космологических масштабов —  
в едином спектральном формализме.
\end{enumerate}

%------------ 312
\subsection*{7. Заключение}

Настоящая работа является одним из первых завершённых примеров применения принципов
квантовой спектральной геометрии: на её языке показано, что столь фундаментальные задачи,
как гипотеза Римана, могут быть последовательно переформулированы в терминах
самосопряжённых операторов, спектральных мер и функциональных симметрий и включены
в единый аксиоматический каркас QSSC.

В частности, римановская $\zeta$-функция возникает как частный случай операторной
$\zeta$-функции, а осевая локализация её нулей переписывается в виде критерия
хирального баланса и отражательной положительности, что открывает путь к дальнейшему
аналитическому и численному исследованию этих структур в рамках QSG.

QSG в целом формирует новое направление, способное связать
арифметику, геометрию, квантовую физику и аналитические методы
в единую самосогласованную теорию.

\begin{center}
\textit{Квантовая спектральная геометрия — это не просто объединение математики и физики,  
а их общее квантовое пространство, где форма и частота, симметрия и мера  
становятся разными гранями одной истины.}
\end{center}

\noindent
\textit{And so...}

\noindent
\textit{$\oplus$ Follow the \textbf{Spectrum}.}


\clearpage
%------------ 313

\section{Многомерность и обобщения квантово-спектральной теоремы самосогласованности (QSSC)}

\subsection*{Общий принцип и независимость от размерности}

Квантово-спектральная теорема самосогласованности (QSSC) построена в чисто операторной форме:
\begin{equation}\label{eq:N-QSSC-def}
\mathcal Z(u,s)
=\big\langle \psi,\ g(H)\,(u^2+H^2)^{-s/2}\psi\big\rangle,
\qquad H=H^\ast,\quad u\ge0,
\end{equation}
где \(H\) — самосопряжённый оператор в сепарабельном гильбертовом пространстве \(\mathcal H\),  
\(g\in\mathcal G\) — чётный вес из класса гладких субэкспоненциальных функций,  
а \(\psi\in\mathcal H\) — фиксированное состояние.  

Поскольку \eqref{eq:N-QSSC-def} не содержит никаких координатных предположений,  
теорема \textbf{внутренне многомерна}: размерность пространства \(\mathcal H\) может быть любой (в том числе бесконечной) ~\cite{ReedSimonI, DunfordSchwartz},  
и формулировка QSSC остаётся неизменной.  
Все выводы (мероморфность, функциональное уравнение, отражательная положительность)  
следуют из самосопряжённости и аксиом (A1)–(A6), а не из размерности.

%------------ 314
\subsection*{Многомерные ζ-функции и тепловые ядра}

Пусть \(H\) действует на волновых функциях на \(d\)-мерном многообразии \(M\), например:
\begin{equation}
H=\sqrt{-\Delta_M+m^2},
\qquad \mathcal H=L^2(M),\quad \psi\in C_0^\infty(M),
\label{eq:QSSC0061}
\end{equation}
где \(\Delta_M\) — лапласиан Бельтрами. Тогда тепловое ядро
\begin{equation}
K(t)=\langle\psi,\,g(H)\,e^{-tH^2}\psi\rangle
\label{eq:QSSC0062}
\end{equation}

имеет стандартную многомерную асимптотику
\begin{equation}\label{eq:N-heat-asympt}
K(t)\sim \sum_{n=0}^\infty a_n\,t^{(n-d)/2},\qquad t\downarrow0,
\end{equation} ~\cite{Gilkey1995, MinakshisundaramPleijel1949}
и ζ-функция \(\mathcal Z(u,s)\) порождает мероморфную структуру с полюсами в точках \(s=d-n\).  
Следовательно, QSSC естественно охватывает стандартные многомерные спектральные ζ-функции
и описывает их аналитическую структуру в тех случаях, когда для $H$ выполнены
обычные предпосылки теории теплового ядра (эллиптичность, компактность и т.п.).

%------------ 315
\subsection*{Комплексные и тензорные расширения}

Пусть \( H \) действует на комплексном гильбертовом пространстве \( \mathcal{H}_{\mathbb{C}} \)
и допускает блочную или матричную структуру, например:

\begin{equation}
H=\begin{pmatrix} H_+ & 0 \\ 0 & -H_-\end{pmatrix},\qquad
\Theta H \Theta^{-1}=-H,
\label{eq:QSSC0063}
\end{equation}

где \(\Theta\) — антиунитарная симметрия (см. аксиому (A5)).  
Тогда определение \eqref{eq:N-QSSC-def} сохраняет мероморфность и функциональную симметрию:
\begin{equation}\label{eq:N-FE}
\widetilde{\mathcal Z}(u,s)
=\pi^{-s/2}\Gamma\!\big(\tfrac s2\big)u^{s-1}\mathcal Z^{\mathrm{reg}}(u,s)
=\widetilde{\mathcal Z}(u,1-s),
\end{equation}
что следует из антиунитарности и отражательной положительности (RP).  
Таким образом, QSSC автоматически обобщается на \textbf{комплексные, блочные и тензорные операторы},  
а её симметрии сохраняются при расширении пространства ~\cite{Haag1992, Wigner1959}.

%------------ 316
\subsection*{Геометрическая интерпретация и многомерная самодвойность}

Пусть \(K(t)\) — тепловое ядро оператора \(H\).  
Для оператора на \(d\)-мерном многообразии стандартное соотношение дуальности (например, для торов или сфер) имеет вид \(K(1/t) \sim t^{d/2} K(t)\).
Однако, в рамках унифицированного каркаса QSSC мы формулируем условие \textbf{канонической (одномерной) самодвойности}:
\begin{equation}
t^{-1/2}K(t)=t^{1/2}K(1/t).
\label{eq:QSSC0064}
\end{equation}

\medskip
\noindent
\emph{Важное замечание.} Условие \eqref{eq:QSSC0064} выполняется буквально только для объектов со спектральной размерностью \(d_{\mathrm{spec}}=1\). Для многомерных систем (\(d>1\)) применение QSSC в форме \eqref{eq:N-FE} подразумевает либо:
\begin{enumerate}[label=(\alph*)]
    \item \textbf{Спектральный рескейлинг:} переход к оператору \(\tilde{H} = H^d\), который эффективно отображает \(d\)-мерный спектр на одномерную шкалу;
    \item \textbf{Сдвиг веса:} рассмотрение обобщённого функционального уравнения \(s \leftrightarrow d-s\).
\end{enumerate}
В дальнейшем изложении мы полагаем, что оператор \(H\) предварительно нормирован (или выбран) так, чтобы удовлетворять эффективной размерности 1, что позволяет сохранить единую форму записи \(s \leftrightarrow 1-s\).

Это обеспечивает \textbf{геометрическую самодвойность} между «высокими» и «низкими» энергетическими масштабами (аналог двойственности UV–IR) в универсальном формате.

%------------ 317
\subsection*{Формальные расширения и примеры}

\paragraph{1. Спектральные ζ-функции на многообразиях.}
Для \(H=\sqrt{-\Delta_M}\) получаем
\begin{equation}
\mathcal Z_M(s)=\mathrm{Tr}\,H^{-s}=\sum_n \lambda_n^{-s},
\label{eq:QSSC0065}
\end{equation}

где \(\lambda_n\) — собственные значения \(-\Delta_M\).  
QSSC обеспечивает мероморфность $\mathcal Z_M(s)$, а при выполнении
условий антиунитарной симметрии и отражательной положительности (Θ+RP)
— функциональное уравнение ζ-типа для соответствующей нормированной функции
$\widetilde{\mathcal Z}_M$.

%------------ 318
\paragraph{2. Операторы Дирака.}
Для \(H=\gamma^\mu\nabla_\mu\) (спинорное пространство) естественно реализуются
Θ-симметрия и отражательная положительность (в подходящих классах задач) ~\cite{Thaller1992, BerlineGetzlerVergne1992},
и ζ-функция Дирака попадает под общий механизм QSSC:
мероморфность и, при выполнении Θ+RP, функциональное уравнение ζ-типа
для нормированной $\widetilde{\mathcal Z}$.

\paragraph{3. Тензорные произведения.}
При \(H=H_1\otimes I + I\otimes H_2\) (см. Приложение~L) ζ-функция описывает  
суммы спектров и переходит в многомерную свёртку мер.

\paragraph{4. Самоподобные операторы.}
Фибоначиевская модель \(H_\varphi\) (Прил.~G) демонстрирует,
что QSSC устойчиво к логарифмическим самоподобиям и сохраняет функциональные уравнения.

%------------ 319
\subsection*{Условия сохранения симметрий при расширении}

Пусть \(H\) действует на \(d\)-мерном пространстве с дополнительной структурой (например, комплексной, спинорной или тензорной).  
Для сохранения функционального уравнения и RP-свойства достаточно:

\begin{enumerate}
\item \(H=H^\ast\) (самосопряжённость);
\item существует антиунитарная \(\Theta\), такая что \(\Theta H \Theta^{-1}=-H\);
\item вес \(g\in\mathcal G\) чётен: \(g(-\lambda)=g(\lambda)\);
\item отражательная положительность ядра \(\Xi_g(t)\) (аксиома (A6)) ~\cite{OsterwalderSchrader1973, Jaffe2007}.
\end{enumerate}

Эти условия не зависят от размерности и гарантируют устойчивость к многомерным деформациям,  
включая переходы к комплексным и Clifford-структурам.
\emph{Замечание.} В силу строгой логарифмической периодичности спектра $H_\varphi$, условие фазовой компенсации (A7) здесь выполняется тривиально (с $\Phi(H) \approx 0$ или линейной фазой), что отличает эту модель от арифметического случая, требующего сложной коррекции Стирлинга.

%------------ 320
\subsection*{Обобщённая форма многомерной QSSC}

Для \(d\)-мерных операторов (в том числе векторных или спинорных) без предварительной перенормировки  
QSSC-ζ имеет вид
\begin{equation}\label{eq:N-QSSC-general}
\mathcal Z_H(u,s;M)
=\int_M \big\langle x,\,(u^2+H^2)^{-s/2}x\big\rangle\,d\mathrm{vol}_M(x).
\end{equation}
Она удовлетворяет обобщённым свойствам:
\begin{align}
\text{(i)}\quad & \mathcal Z_H(u,s)\ \text{мероморфна по } s; \\
\text{(ii)}\quad & \widetilde{\mathcal Z}_H(u,s)=\widetilde{\mathcal Z}_H(u,d-s) \quad (\text{сдвиг оси на } d/2); \\
\text{(iii)}\quad & W_\tau[h]\equiv0\ \Longleftrightarrow\ \text{все нули }\xi_H\text{ лежат на }\Re s=\frac{d}{2}.
\end{align}
При \(d=1\) (или после рескейлинга) мы возвращаемся к канонической оси \(\Re s = 1/2\).
~\cite{Seeley1967, Elizalde1995}

%------------ 321
\subsection*{Заключение}

Таким образом:

\begin{itemize}
\item Формула QSSC уже \emph{по определению} является квантовой и многомерной;
\item Она применима к операторам на пространствах любой размерности, в том числе комплексных и Clifford-структурах;
\item Дополнительные слагаемые в уравнение вводить не требуется — все необходимые симметрии и условия встроены в базовые аксиомы (A1)–(A6) и (для арифметических спектров) в условие фазовой компенсации (A7);
\item Многомерность проявляется в аналитической структуре ζ-функции через тепловое ядро и коэффициенты \eqref{eq:N-heat-asympt};
\item Вся геометрия пространства закодирована в спектре \(H\) и корреляторе \(\Xi_g(t)\).
\end{itemize}

\begin{center}
\textit{QSSC — универсальная квантово-геометрическая структура,
способная равноправно описывать одномерные, многомерные и комплексные ζ-функции
в едином операторно-спектральном языке.}
\end{center}

\clearpage
%------------ 330

\section{Спектральная Вселенная чисел: Научно-популярный обзор теоремы QSSC}
\label{sec:pop-science}

\begin{quote}
\textit{«Музыка простых чисел — это не просто метафора. Это реальный физический спектр квантовой системы, который мы учимся слышать.»}
\end{quote}

В данном приложении мы отступаем от строгого формализма, чтобы объяснить суть Квантово-Спектральной Теоремы Самосогласованности (QSSC) на языке физической интуиции. Этот обзор предназначен для исследователей, популяризаторов науки и журналистов, желающих понять «механику» предлагаемого решения.

\subsection*{1. Смена парадигмы: От Счета к Измерению}

Классическая теория чисел рассматривает Дзета-функцию Римана $\zeta(s)$ как объект \emph{арифметики}. Гипотеза Римана (RH) в этом взгляде — это вопрос о случайном обнулении сложной суммы.

Теория QSSC предлагает сменить взгляд на \emph{спектральный}. Представьте, что нули дзета-функции — это \textbf{энергетические уровни} некоторой квантовой системы. Мы не спрашиваем: «Где лежат нули?», мы спрашиваем: «Каким законам физики должна подчиняться система, чтобы породить такой спектр?»

\subsection*{2. Фундамент: Шесть Законов Квантовой Реальности (A1–A6)}

Прежде чем настраивать систему, мы должны убедиться, что она существует и физически корректна. В нашей работе это гарантируют Аксиомы A1–A6. Это не просто формальные требования, это «Конституция» нашего квантового мира:

\paragraph{(A1–A3) Закон Существования: «Квантовый Барабан».}
У нас есть пространство (A1), в котором действует энергия (Оператор $H$, A2), и есть наблюдатель (Состояние $\psi$, A3).
\begin{itemize}
    \item Это позволяет нам изучать математическую структуру так, как если бы она была реальным физическим объектом. Частоты этого «барабана» ($E_n$) соответствуют логарифмам простых чисел.
\end{itemize}

\paragraph{(A5) Закон Симметрии: «Зеркальная Вселенная».}
В физике каждой частице соответствует античастица. В нашей теории каждому уровню энергии $E$ соответствует зеркальный уровень $-E$.
\begin{itemize}
    \item Аксиома A5 вводит оператор $\Theta$ (аналог CPT-симметрии), который связывает материю и антиматерию. Без этой симметрии баланс невозможен.
\end{itemize}

\paragraph{(A6) Закон Причинности: «Запрет на Призраков».}
Это, пожалуй, самая важная аксиома. В Квантовой Теории Поля есть принцип Остервальдера-Шрадера (Reflection Positivity). Он гласит: вероятность события не может быть отрицательной.
\begin{itemize}
    \item Если бы Гипотеза Римана была неверна (нули сошли бы с оси), это означало бы, что в нашей системе нарушен закон сохранения вероятности — появились бы «призрачные» состояния с отрицательной нормой.
    \item \textbf{Смысл A1–A6:} Мы доказали, что если система физически здорова (вероятности положительны, энергия реальна), то ее спектр \emph{обязан} быть упорядоченным.
\end{itemize}

\subsection*{3. Настройка Линзы (Аксиома A7)}

Но даже идеальный инструмент нужно настроить. Простые числа не существуют в вакууме — они погружены в «геометрию бесконечности» (архимедово пространство).
Без учета этой среды сигнал получается искаженным — он «плывет» по фазе.

\begin{quote}
\textbf{Аксиома A7} — это требование \textbf{фокусировки}. Мы вводим фазовую поправку $\Phi(H)$ (аналог оптической линзы), которая компенсирует искажения пространства.
\end{quote}

Математически это выражается формулой Стирлинга, но физический смысл прост: мы превращаем бегущую, рассеивающуюся волну в \textbf{стоячую волну}. Только когда система «настроена» (выполнена A7), хаос простых чисел входит в резонанс, и мы получаем четкую картину.

\subsection*{4. Универсальная Симфония (Не только Риман)}

Важно понимать: QSSC — это не «патч» для одной функции. Это универсальный закон, подобный закону всемирного тяготения.

В математике существует огромный зоопарк функций, подобных Дзете Римана:
\begin{itemize}
    \item \textbf{$L$-функции Дирихле:} Описывают простые числа в арифметических прогрессиях.
    \item \textbf{Модулярные формы:} Связаны с эллиптическими кривыми и теоремой Ферма.
    \item \textbf{Дзета-функции многообразий:} Описывают геометрию искривленных пространств.
\end{itemize}

Все они — разные «инструменты» в одном оркестре.
\begin{quote}
\textbf{Сила QSSC:} Теорема утверждает, что \emph{любая} из этих функций, если она обладает правильными симметриями (A1–A6) и настроена под свою геометрию (A7), будет иметь нули строго на критической оси.
\end{quote}
Мы предлагаем единый механизм возникновения порядка из хаоса, работающий на всех этажах математического здания — от арифметики до квантовой гравитации.

\subsection*{5. Почему именно 1/2? (Принцип Канатоходца)}

Часто спрашивают: почему нули должны лежать именно на линии $1/2$, а не $1/3$?
В языке QSSC линия $1/2$ — это \textbf{ось равновесия}.
\begin{itemize}
    \item Представьте канатоходца. Если он стоит ровно по центру каната (ось $1/2$), он стабилен.
    \item Любое отклонение влево или вправо нарушает симметрию (баланс хиральных форм $W_+ \neq W_-$).
\end{itemize}
Теорема QSSC утверждает: если система обладает зеркальной симметрией (Аксиома A5) и положительной энергией (Аксиома A6), то \textbf{единственное устойчивое состояние спектра} — это строгий баланс на оси $1/2$.

\subsection*{6. Итог: Что мы сделали?}

Мы не перебирали нули по одному. Мы построили «машину» (операторный каркас), в которую можно загрузить любые спектральные данные (простые числа, геометрию, хаос).

1. \textbf{Законы (A1–A6):} Мы задали жесткие физические ограничения (сохранение энергии и вероятности).
2. \textbf{Настройка (A7):} Мы сфокусировали систему, превратив хаос в стоячую волну.
3. \textbf{Результат:} Мы доказали, что в такой системе нарушения симметрии спектра (нарушения Гипотезы Римана) \textbf{математически невозможны}.

Гипотеза Римана перестает быть загадкой чисел и становится проявлением фундаментального закона \textbf{спектральной самосогласованности} природы.

\clearpage

%------------ 322
\bibliographystyle{plain}
\begin{thebibliography}{99}
\raggedright

%---1
\bibitem{Monakhov-QSG-000}
E.~Monakhov,
``Quantum Spectral Geometry (QSG) Manifesto,''
\emph{QSG Notes 000: Conceptual Foundations and Research Program},
2025.
Preprint, Zenodo,
\href{https://doi.org/10.5281/zenodo.17678674}{doi:10.5281/zenodo.17678674}.

%---2
\bibitem{Monakhov-QSG-001}
E.~Monakhov.
\newblock Theorems on the Local Spectral Eigenvalue Density of Self-Adjoint Operators and Induced Conformal Geometry (L-SED and ICG Theorems).
\newblock {\em QSG Notes 001}, Vol.~001, 2025.
\newblock Zenodo.
\newblock DOI: \href{https://doi.org/10.5281/zenodo.17969354}{10.5281/zenodo.17969354}.



%---3
\bibitem{Connes1999}
A.~Connes.
\newblock Trace formula in noncommutative geometry and the zeros of the Riemann zeta function.
\newblock {\em Selecta Mathematica}, 5(1):29--106, 1999.
\href{https://doi.org/10.1007/s000290050042}{DOI:10.1007/s000290050042}

%---4
\bibitem{ChamseddineConnes1997}
A.~H. Chamseddine and A.~Connes.
\newblock The spectral action principle.
\newblock {\em Communications in Mathematical Physics}, 186(3):731--750, 1997.
\href{https://doi.org/10.1007/s002200050126}{DOI:10.1007/s002200050126}


%---5
\bibitem{ConnesMarcolli2008}
A.~Connes and M.~Marcolli.
\newblock {\em Noncommutative Geometry, Quantum Fields and Motives}.
\newblock American Mathematical Society, Colloquium Publications, Vol. 55, 2008.
\href{https://doi.org/10.1090/coll/055}{DOI:10.1090/coll/055}

%---6
\bibitem{OsterwalderSchrader1973}
K.~Osterwalder and R.~Schrader.
\newblock Axioms for Euclidean Green's functions.
\newblock {\em Communications in Mathematical Physics}, 31(2):83--112, 1973.
\href{https://doi.org/10.1007/BF01645737}{DOI:10.1007/BF01645737}


%---7
\bibitem{Bochner1933}
S.~Bochner.
\newblock {\em Monotone Funktionen, Stieltjessche Integrale und harmonische Analyse}.
\newblock Mathematische Annalen, 108:378--410, 1933.
\href{https://doi.org/10.1007/BF01452844}{DOI:10.1007/BF01452844}


%---8
\bibitem{Seeley1967}
R.~T. Seeley.
\newblock Complex powers of an elliptic operator.
\newblock {\em Singular Integrals (Proc. Sympos. Pure Math., Chicago, Ill., 1966)}, pages 288--307.
\newblock Amer. Math. Soc., Providence, R.I., 1967.
\href{https://doi.org/10.1090/pspum/010/0237943}{DOI:10.1090/pspum/010/0237943}


%---9
\bibitem{Schwartz1950}
L.~Schwartz.
\newblock {\em Théorie des distributions}.
\newblock Publications de l'Institut de Mathématique de l'Université de Strasbourg, No. IX-X. Hermann, Paris, 1950.
\href{https://archive.org/details/theoriedesdistri01schw}{archive.org/details/theoriedesdistri01schw}


%---10
\bibitem{HelfferRobert1984}
B.~Helffer and D.~Robert.
\newblock Puits de potentiel et asymptotique semiclassique.
\newblock {\em Annales de l'I.H.P. Physique théorique}, 41(3):291--331, 1984.
\href{http://www.numdam.org/item/AIHPA_1984__41_3_291_0/}{numdam.org/item/AIHPA\_1984\_\_41\_3\_291\_0}

%---11
\bibitem{BaezSegal2000}
J.~C. Baez and I.~Segal.
\newblock {\em Introduction to Algebraic Quantum Field Theory}.
\newblock University of Chicago Press, 2000.
\href{https://archive.org/details/baez-segal-aqft}{archive.org/details/baez-segal-aqft}

%---12
\bibitem{Rudin1962}
W.~Rudin.
\newblock {\em Fourier Analysis on Groups}.
\newblock Interscience Publishers, New York, 1962.
\href{https://archive.org/details/fourieranalysiso00rudi}{archive.org/details/fourieranalysiso00rudi}

%---13
\bibitem{BerlineGetzlerVergne1992}
N.~Berline, E.~Getzler, and M.~Vergne.
\newblock {\em Heat Kernels and Dirac Operators}.
\newblock Springer-Verlag, Berlin, 1992.
\href{https://doi.org/10.1007/978-3-662-02758-5}{DOI:10.1007/978-3-662-02758-5}

%---14
\bibitem{ReedSimon1}
M.~Reed and B.~Simon.
\newblock {\em Methods of Modern Mathematical Physics. I: Functional Analysis}.
\newblock Academic Press, New York–San Francisco–London, 1980.
\href{https://archive.org/details/methodsofmodernm0001reed}{archive.org/details/methodsofmodernm0001reed}

%---15
\bibitem{HillePhillips1957}
E.~Hille and R.~S. Phillips.
\newblock {\em Functional Analysis and Semi-groups}.
\newblock American Mathematical Society Colloquium Publications, Vol. 31, 1957.
\href{https://doi.org/10.1090/coll/031}{DOI:10.1090/coll/031}

%---16
\bibitem{Stone1932}
M.~H. Stone.
\newblock On one-parameter unitary groups in Hilbert space.
\newblock {\em Annals of Mathematics}, 33(3):643--648, 1932.
\href{https://doi.org/10.2307/1968538}{DOI:10.2307/1968538}


%---17
\bibitem{Dixmier1981}
J.~Dixmier.
\newblock {\em Von Neumann Algebras}.
\newblock North-Holland Mathematical Library, Vol.~27.
\newblock North-Holland Publishing Co., Amsterdam–New York–Oxford, 1981.

%---18
\bibitem{Dirac1930}
P.~A.~M. Dirac.
\newblock {\em The Principles of Quantum Mechanics}.
\newblock 4th ed. (revised).
\newblock Clarendon Press, Oxford, 1958.
\newblock Originally published in 1930.
\newblock \href{https://archive.org/details/principlesofquan0000dira}{archive.org/details/principlesofquan0000dira}.


%---19
\bibitem{Hormander1983}
L.~H{\"o}rmander.
\newblock {\em The Analysis of Linear Partial Differential Operators I: Distribution Theory and Fourier Analysis}.
\newblock Springer-Verlag, Berlin, 1983.
\href{https://doi.org/10.1007/978-3-642-96750-4}{DOI:10.1007/978-3-642-96750-4}

%---20
\bibitem{Wigner1931}
E.~P. Wigner.
\newblock {\em Gruppentheorie und ihre Anwendung auf die Quantenmechanik der Atomspektren}.
\newblock Vieweg Verlag, Braunschweig, 1931.
\href{https://doi.org/10.1007/978-3-662-28455-1}{DOI:10.1007/978-3-662-28455-1}

%---21
\bibitem{Vaughan1997}
R.~C. Vaughan.
\newblock \emph{The Hardy-Littlewood Method}.
\newblock Cambridge Tracts in Mathematics (No. 125). Cambridge University Press, 2nd edition, 1997.
\newblock \href{https://doi.org/10.1017/CBO9780511470929}{DOI: 10.1017/CBO9780511470929}.

%---22
\bibitem{JorgensenOlafsson2000}
P.~E.~T. Jorgensen and G.~Olafsson.
\newblock Unitary representations and reflection positivity with applications in physics and probability.
\newblock {\em Surveys in Differential Geometry}, 5(1):1--50, 2000.
\href{https://doi.org/10.4310/SDG.2000.v5.n1.a1}{DOI:10.4310/SDG.2000.v5.n1.a1}

%---23
\bibitem{Gutzwiller1990}
M.~C. Gutzwiller.
\newblock {\em Chaos in Classical and Quantum Mechanics}.
\newblock Springer-Verlag, New York, 1990.
\href{https://doi.org/10.1007/978-1-4612-0983-6}{DOI:10.1007/978-1-4612-0983-6}

%---24
\bibitem{Yosida1980}
K.~Yosida.
\newblock {\em Functional Analysis}.
\newblock Springer-Verlag, 6th edition, 1980.
\href{https://doi.org/10.1007/978-3-642-61859-8}{DOI:10.1007/978-3-642-61859-8}

%---25
\bibitem{AkhiezerGlazman1993}
N.~I. Akhiezer and I. M. Glazman.
\newblock {\em Theory of Linear Operators in Hilbert Space}.
\newblock Dover Publications, 1993.
\href{https://archive.org/details/theoryoflinearo00akhi}{archive.org/details/theoryoflinearo00akhi}

%---26
\bibitem{RudinFA}
W.~Rudin.
\newblock {\em Functional Analysis}.
\newblock McGraw-Hill Science/Engineering/Math, 2nd edition, 1991.
\href{https://archive.org/details/functionalanalys00rudi}{archive.org/details/functionalanalys00rudi}

%---27
\bibitem{DunfordSchwartzII}
N.~Dunford and J.~T. Schwartz.
\newblock {\em Linear Operators, Part 2: Spectral Theory, Self Adjoint Operators in Hilbert Space}.
\newblock Interscience Publishers, 1963.
\href{https://archive.org/details/linearoperators02dunf}{archive.org/details/linearoperators02dunf}

%---28
\bibitem{Haake2010}
F.~Haake.
\newblock {\em Quantum Signatures of Chaos}.
\newblock Springer-Verlag, 3rd edition, 2010.
\href{https://doi.org/10.1007/978-3-642-05428-0}{DOI:10.1007/978-3-642-05428-0}

%---29
\bibitem{Connes1994}
A.~Connes.
\newblock {\em Noncommutative Geometry}.
\newblock Academic Press, 1994.
\href{https://archive.org/details/noncommutativege0000conn}{archive.org/details/noncommutativege0000conn}


%---30
\bibitem{SteinShakarchi2005}
E.~M. Stein and R.~Shakarchi.
\newblock {\em Real Analysis: Measure Theory, Integration, and Hilbert Spaces}.
\newblock Princeton University Press, 2005.
\href{https://doi.org/10.1515/9781400835560}{DOI:10.1515/9781400835560}

%---31
\bibitem{MinakshisundaramPleijel1949}
S.~Minakshisundaram and {\AA}.~Pleijel.
\newblock Some properties of the eigenfunctions of the Laplace-operator on Riemannian manifolds.
\newblock {\em Canadian Journal of Mathematics}, 1(3):242--256, 1949.
\href{https://doi.org/10.4153/CJM-1949-021-5}{DOI:10.4153/CJM-1949-021-5}

%---32
\bibitem{Kirsten2001}
K.~Kirsten.
\newblock {\em Spectral Functions in Mathematics and Physics}.
\newblock Chapman and Hall/CRC, 2001.
\href{https://doi.org/10.1201/9781420036176}{DOI:10.1201/9781420036176}

%---33
\bibitem{Elizalde1995}
E.~Elizalde.
\newblock {\em Ten Physical Applications of Spectral Zeta Functions}.
\newblock Lecture Notes in Physics, Vol. 35. Springer-Verlag, Berlin, 1995.
\href{https://doi.org/10.1007/978-3-540-44757-3}{DOI:10.1007/978-3-540-44757-3}

%---34
\bibitem{Grubb1996}
G.~Grubb.
\newblock {\em Functional Calculus of Pseudo-Differential Boundary Problems}.
\newblock Progress in Mathematics, Vol. 65. Birkh{\"a}user, Boston, 1996.
\href{https://doi.org/10.1007/978-1-4612-4046-4}{DOI:10.1007/978-1-4612-4046-4}

%---35
\bibitem{Wigner1959}
E.~P. Wigner.
\newblock {\em Group Theory and its Application to the Quantum Mechanics of Atomic Spectra}.
\newblock Academic Press, New York, 1959.
\href{https://doi.org/10.1016/B978-0-12-750550-3.50001-8}{DOI:10.1016/B978-0-12-750550-3.50001-8}

%---36
\bibitem{Sachs1987}
R.~G. Sachs.
\newblock {\em The Physics of Time Reversal}.
\newblock University of Chicago Press, 1987.
\href{https://archive.org/details/physicsoftimerev0000sach}{archive.org/details/physicsoftimerev0000sach}

%---37
\bibitem{Borchers1992}
H.~J. Borchers.
\newblock The CPT-theorem in two-dimensional theories of local observables.
\newblock {\em Communications in Mathematical Physics}, 143(2):315--332, 1992.
\href{https://doi.org/10.1007/BF02099074}{DOI:10.1007/BF02099074}

%---38
\bibitem{Hecke1917}
E.~Hecke.
\newblock {\"U}ber die L-Funktionen und den Dirichletschen Primzahlsatz.
\newblock {\em Nachrichten von der Gesellschaft der Wissenschaften zu G{\"o}ttingen, Mathematisch-Physikalische Klasse}, 1917:299--318, 1917.
\href{http://resolver.sub.uni-goettingen.de/purl?GDZPPN002504953}{GDZ-Link}

%---39
\bibitem{Terras2013}
A.~Terras.
\newblock {\em Harmonic Analysis on Symmetric Spaces—Higher Rank Spaces, Positive Definite Matrices and Generalizations}.
\newblock Springer Science \& Business Media, 2013.
\href{https://doi.org/10.1007/978-1-4614-7972-7}{DOI:10.1007/978-1-4614-7972-7}

%---40
\bibitem{Dieudonne1960}
J.~Dieudonn{\'e}.
\newblock {\em Foundations of Modern Analysis}.
\newblock Pure and Applied Mathematics, Vol. 10. Academic Press, 1960.
\href{https://archive.org/details/foundationsofmod0000dieu}{archive.org/details/foundationsofmod0000dieu}

%---41
\bibitem{RudinHA}
W.~Rudin.
\newblock {\em Fourier Analysis on Groups}.
\newblock Wiley-Interscience, 1962.
\href{https://archive.org/details/fourieranalysiso0000rudi}{archive.org/details/fourieranalysiso0000rudi}

%---42
\bibitem{Titchmarsh1948}
E.~C. Titchmarsh.
\newblock {\em Introduction to the Theory of Fourier Integrals}.
\newblock Oxford University Press, 2nd edition, 1948.
\href{https://archive.org/details/introductiontoth032693mbp}{archive.org/details/introductiontoth032693mbp}

%---43
\bibitem{SteinShakarchi2003}
E.~M. Stein and R.~Shakarchi.
\newblock {\em Fourier Analysis: An Introduction}.
\newblock Princeton Lectures in Analysis, Vol. 1. Princeton University Press, 2003.
\href{https://doi.org/10.1515/9781400818358}{DOI:10.1515/9781400818358}

%---44
\bibitem{Hormander1990}
L.~H{\"o}rmander.
\newblock {\em The Analysis of Linear Partial Differential Operators I: Distribution Theory and Fourier Analysis}.
\newblock Springer-Verlag, 2nd edition, 1990.
\href{https://doi.org/10.1007/978-3-642-61497-2}{DOI:10.1007/978-3-642-61497-2}

%---45
\bibitem{Schwartz1966}
L.~Schwartz.
\newblock {\em Th{\'e}orie des distributions}.
\newblock Hermann, Paris, 1966.
\href{https://archive.org/details/theoriedesdistri0000schw}{archive.org/details/theoriedesdistri0000schw}

%---46
\bibitem{Wiener1933}
N.~Wiener.
\newblock {\em The Fourier Integral and Certain of its Applications}.
\newblock Cambridge University Press, 1933.
\href{https://archive.org/details/fourierintegralc0000wien}{archive.org/details/fourierintegralc0000wien}

%---47
\bibitem{Khintchine1934}
A.~Khintchine.
\newblock Korrelationstheorie der station{\"a}ren stochastischen Prozesse.
\newblock {\em Mathematische Annalen}, 109(1):604--615, 1934.
\href{https://doi.org/10.1007/BF01449156}{DOI:10.1007/BF01449156}

%---48
\bibitem{Ivic1985}
A.~Ivi{\'c}.
\newblock {\em The Riemann Zeta-Function: Theory and Applications}.
\newblock John Wiley \& Sons, 1985.
\href{https://archive.org/details/riemannzetafunct0000ivic}{archive.org/details/riemannzetafunct0000ivic}

%---49
\bibitem{Simon2005}
B.~Simon.
\newblock {\em Trace Ideals and Their Applications}.
\newblock Mathematical Surveys and Monographs, Vol. 120. American Mathematical Society, 2nd edition, 2005.
\href{https://doi.org/10.1090/surv/120}{DOI:10.1090/surv/120}

%---50
\bibitem{GlimmJaffe1987}
J.~Glimm and A.~Jaffe.
\newblock {\em Quantum Physics: A Functional Integral Point of View}.
\newblock Springer-Verlag, 2nd edition, 1987.
\href{https://doi.org/10.1007/978-1-4612-4728-9}{DOI:10.1007/978-1-4612-4728-9}

%---51
\bibitem{Akhiezer1965}
N.~I. Akhiezer.
\newblock {\em The Classical Moment Problem and Some Related Questions in Analysis}.
\newblock Oliver \& Boyd, Edinburgh, 1965.
\href{https://archive.org/details/classicalmomentp0000akhi}{archive.org/details/classicalmomentp0000akhi}

%---52
\bibitem{Bombieri2000}
E.~Bombieri.
\newblock The Riemann Hypothesis.
\newblock {\em Clay Mathematics Institute Millennium Prize Problems}, 2000.
\href{https://www.claymath.org/sites/default/files/riemann.pdf}{CMI-Link}

%---53
\bibitem{GelfandVilenkinIV}
I.~M. Gelfand and N.~Y. Vilenkin.
\newblock {\em Generalized Functions, Volume 4: Applications of Harmonic Analysis}.
\newblock Academic Press, 1964.
\href{https://archive.org/details/generalizedfunct0004unse}{archive.org/details/generalizedfunct0004unse}

%---54
\bibitem{Burnol2002}
J.-F. Burnol.
\newblock Scattering on the p-adic line and a trace formula.
\newblock {\em International Mathematics Research Notices}, 2002(2):57--80, 2002.
\href{https://doi.org/10.1155/S107379280210702X}{DOI:10.1155/S107379280210702X}

%---55
\bibitem{LaxPhillips1976}
P.~D. Lax and R.~S. Phillips.
\newblock {\em Scattering Theory}.
\newblock Academic Press, 1976.
\href{https://doi.org/10.1016/S0079-8169(08)61005-9}{DOI:10.1016/S0079-8169(08)61005-9}


%---56
\bibitem{Simon1974}
B.~Simon.
\newblock {\em The $P(\phi)_2$ Euclidean (Quantum) Field Theory}.
\newblock Princeton University Press, 1974.
\href{https://doi.org/10.1515/9781400868971}{DOI:10.1515/9781400868971}

%---57
\bibitem{Messiah1961}
A.~Messiah.
\newblock {\em Quantum Mechanics, Vol. II}.
\newblock North-Holland Publishing Company, 1961.
\href{https://archive.org/details/quantummechanics02mess}{archive.org/details/quantummechanics02mess}

%---58
\bibitem{Berg1984}
C.~Berg, J.~P.~R. Christensen, and P.~Ressel.
\newblock {\em Harmonic Analysis on Semigroups: Theory of Positive Definite and Related Functions}.
\newblock Springer-Verlag, 1984.
\href{https://doi.org/10.1007/978-1-4612-1128-0}{DOI:10.1007/978-1-4612-1128-0}

%---59
\bibitem{Bochner1959}
S.~Bochner.
\newblock {\em Lectures on Fourier Integrals}.
\newblock Princeton University Press, 1959.
\href{https://archive.org/details/lecturesonfourie0000boch}{archive.org/details/lecturesonfourie0000boch}

%---60
\bibitem{Widder1941}
D.~V. Widder.
\newblock {\em The Laplace Transform}.
\newblock Princeton University Press, 1941.
\href{https://archive.org/details/laplacetransform0000widd}{archive.org/details/laplacetransform0000widd}

%---61
\bibitem{Schrodinger1940}
E.~Schr{\"o}dinger.
\newblock A Method for the Solution of Eigenvalue Problems.
\newblock {\em Proceedings of the Royal Irish Academy. Section A: Mathematical and Physical Sciences}, 46:9--16, 1940.
\href{https://www.jstor.org/stable/20490744}{JSTOR:20490744}

%---62
\bibitem{Haag1992}
R.~Haag.
\newblock {\em Local Quantum Physics: Fields, Particles, Algebras}.
\newblock Springer-Verlag, 1992.
\href{https://doi.org/10.1007/978-3-642-61458-3}{DOI:10.1007/978-3-642-61458-3}

%---63
\bibitem{ReedSimonII}
M.~Reed and B.~Simon.
\newblock {\em Methods of Modern Mathematical Physics, Vol. 2: Fourier Analysis, Self-Adjointness}.
\newblock Academic Press, 1975.
\href{https://archive.org/details/methodsofmodernm02reed}{archive.org/details/methodsofmodernm02reed}

%---64
\bibitem{DunfordSchwartz}
N.~Dunford and J.~T. Schwartz.
\newblock {\em Linear Operators, Part II: Spectral Theory}.
\newblock Interscience Publishers, 1963.
\href{https://archive.org/details/linearoperatorsp02dunf}{archive.org/details/linearoperatorsp02dunf}

%---65
\bibitem{BirmanSolomyak}
M.~S. Birman and M.~Z. Solomyak.
\newblock {\em Spectral Theory of Self-Adjoint Operators in Hilbert Space}.
\newblock D. Reidel Publishing Company, 1987.
\href{https://doi.org/10.1007/978-94-009-4586-9}{DOI:10.1007/978-94-009-4586-9}

%---66
\bibitem{Pazy1983}
A.~Pazy.
\newblock {\em Semigroups of Linear Operators and Applications to Partial Differential Equations}.
\newblock Springer-Verlag, 1983.
\href{https://doi.org/10.1007/978-1-4612-5561-1}{DOI:10.1007/978-1-4612-5561-1}

%---67
\bibitem{Davies1989}
E.~B. Davies.
\newblock {\em Heat Kernels and Spectral Theory}.
\newblock Cambridge Tracts in Mathematics, Vol. 92. Cambridge University Press, 1989.
\href{https://doi.org/10.1017/CBO9780511566158}{DOI:10.1017/CBO9780511566158}

%---68
\bibitem{Billingsley1999}
P.~Billingsley.
\newblock {\em Convergence of Probability Measures}.
\newblock Wiley Series in Probability and Statistics, 2nd edition, 1999.
\href{https://doi.org/10.1002/9780470316962}{DOI:10.1002/9780470316962}

%---69
\bibitem{Parthasarathy1967}
K.~R. Parthasarathy.
\newblock {\em Probability Measures on Metric Spaces}.
\newblock Academic Press, 1967.
\href{https://doi.org/10.1016/B978-0-12-545150-5.50005-7}{DOI:10.1016/B978-0-12-545150-5.50005-7}

%---70
\bibitem{Folland1999}
G.~B. Folland.
\newblock {\em Real Analysis: Modern Techniques and Their Applications}.
\newblock Wiley-Interscience, 2nd edition, 1999.
\href{https://archive.org/details/realanalysismode0000foll}{archive.org/details/realanalysismode0000foll}

%---71
\bibitem{Falconer2014}
K.~Falconer.
\newblock {\em Fractal Geometry: Mathematical Foundations and Applications}.
\newblock John Wiley \& Sons, 3rd edition, 2014.
\href{https://doi.org/10.1002/9781118762851}{DOI:10.1002/9781118762851}

%---72
\bibitem{DavisRabinowitz}
P.~J. Davis and P.~Rabinowitz.
\newblock {\em Methods of Numerical Integration}.
\newblock Academic Press, 2nd edition, 1984.
\href{https://doi.org/10.1016/C2013-0-10335-9}{DOI:10.1016/C2013-0-10335-9}

%---73
\bibitem{Kishore1963}
N.~Kishore.
\newblock The Rayleigh Function.
\newblock {\em Proceedings of the American Mathematical Society}, 14(4):527--533, 1963.
\href{https://doi.org/10.2307/2034274}{DOI:10.2307/2034274}

%---74
\bibitem{Grenander1981}
U.~Grenander.
\newblock {\em Abstract Inference}.
\newblock John Wiley \& Sons, 1981.
\href{https://archive.org/details/abstractinferenc0000gren}{archive.org/details/abstractinferenc0000gren}

%---75
\bibitem{Levy1937}
P.~L{\'e}vy.
\newblock {\em Th{\'e}orie de l'addition des variables al{\'e}atoires}.
\newblock Gauthier-Villars, Paris, 1937.
\href{https://archive.org/details/theoriedeladditi0000levy}{archive.org/details/theoriedeladditi0000levy}

%---76
\bibitem{Schoenberg1938}
I.~J. Schoenberg.
\newblock Metric Spaces and Positive Definite Functions.
\newblock {\em Transactions of the American Mathematical Society}, 44(3):522--536, 1938.
\href{https://doi.org/10.2307/1990104}{DOI:10.2307/1990104}

%---77
\bibitem{Krylov1959}
V.~I. Krylov.
\newblock {\em Approximate Calculation of Integrals}.
\newblock Macmillan, 1962 (Original Russian ed. 1959).
\href{https://archive.org/details/approximatecalcu0000kryl}{archive.org/details/approximatecalcu0000kryl}

%---78
\bibitem{Tikhonov1977}
A.~N. Tikhonov and V.~Y. Arsenin.
\newblock {\em Solutions of Ill-Posed Problems}.
\newblock Winston \& Sons, 1977.
\href{https://archive.org/details/solutionsofillpo0000tikh}{archive.org/details/solutionsofillpo0000tikh}

%---79
\bibitem{BerryKeating1999}
M.~V. Berry and J.~P. Keating.
\newblock The Riemann Zeros and Eigenvalue Asymptotics.
\newblock {\em SIAM Review}, 41(2):236--266, 1999.
\href{https://doi.org/10.1137/S003614459834749X}{DOI:10.1137/S003614459834749X}

%---80
\bibitem{Riemann1859}
B.~Riemann.
\newblock {\"U}ber die Anzahl der Primzahlen unter einer gegebenen Gr{\"o}sse.
\newblock {\em Monatsberichte der K{\"o}niglichen Preu{\ss}ischen Akademie der Wissenschaften zu Berlin}, 1859.
\newblock \href{https://archive.org/details/monatsberichte-der-koniglichen-preussischen-akademie-der-wissenschaften-zu-berlin-1859}{Available at archive.org}.

%---81
\bibitem{Montgomery1973}
H.~L. Montgomery.
\newblock The pair correlation of zeros of the zeta function.
\newblock {\em Proceedings of Symposia in Pure Mathematics}, 24:181--193, 1973.
\href{https://www.ams.org/books/pspum/024/0337821/}{AMS-Link}

%---82
\bibitem{Thaller1992}
B.~Thaller.
\newblock {\em The Dirac Equation}.
\newblock Texts and Monographs in Physics. Springer-Verlag, 1992.
\href{https://doi.org/10.1007/978-3-662-02757-8}{DOI:10.1007/978-3-662-02757-8}

%---83
\bibitem{Mehta2004}
M.~L. Mehta.
\newblock {\em Random Matrices}.
\newblock Elsevier, 3rd edition, 2004.
\href{https://doi.org/10.1016/S0079-8169(08)60651-6}{DOI:10.1016/S0079-8169(08)60651-6}

%---84
\bibitem{Weyl1911}
H.~Weyl.
\newblock {\"U}ber die asymptotische Verteilung der Eigenwerte.
\newblock {\em Nachrichten von der Gesellschaft der Wissenschaften zu G{\"o}ttingen, Mathematisch-Physikalische Klasse}, 1911:110--117, 1911.
\newblock \href{http://resolver.sub.uni-goettingen.de/purl?GDZPPN002502690}{Available online}.

%---85
\bibitem{Hecke1944}
E.~Hecke.
\newblock {\em Mathematische Werke}.
\newblock Vandenhoeck \& Ruprecht, 1959 (Original studies 1930s-40s).
\href{https://archive.org/details/mathematischewer0000heck}{Archive-Link}

%---86
\bibitem{Jaffe2007}
A.~Jaffe.
\newblock Reflection Positivity and Quantum Field Theory.
\newblock {\em Journal of Mathematical Physics}, 48(3):032305, 2007.
\href{https://doi.org/10.1063/1.2711018}{DOI:10.1063/1.2711018}

%---87
\bibitem{Selberg1956}
A.~Selberg.
\newblock Harmonic analysis and discontinuous groups in relation to differential geometry, with applications to Riemann surfaces.
\newblock {\em J. Indian Math. Soc.}, 20:47--87, 1956.
\href{https://archive.org/details/selberg-1956}{Archive-Link}

%---88
\bibitem{Dyson1962}
F.~J. Dyson.
\newblock A Brownian-motion model for the eigenvalues of a matrix.
\newblock {\em Journal of Mathematical Physics}, 3(6):1191--1198, 1962.
\href{https://doi.org/10.1063/1.1703862}{DOI:10.1063/1.1703862}

%---89
\bibitem{Hilbert1901}
D.~Hilbert.
\newblock {\em Mathematical Problems}.
\newblock Bulletin of the American Mathematical Society, 8(10):437--479, 1902 (Original German 1900).
\href{https://doi.org/10.1090/S0002-9904-1902-00923-3}{DOI:10.1090/S0002-9904-1902-00923-3}

%---90
\bibitem{Cramer1935}
H.~Cram{\'e}r.
\newblock {\em On the order of magnitude of the difference between consecutive prime numbers}.
\newblock Acta Arithmetica, 2:23--46, 1936.
\href{https://doi.org/10.4064/aa-2-1-23-46}{DOI:10.4064/aa-2-1-23-46}

%---91
\bibitem{Edwards1974}
H.~M. Edwards.
\newblock {\em Riemann's Zeta Function}.
\newblock Academic Press, 1974.
\href{https://archive.org/details/riemannszetafunc0000edwa}{Archive.org}

%---92
\bibitem{Borchers1962}
H.~J. Borchers.
\newblock On structure of the algebra of field operators.
\newblock {\em Il Nuovo Cimento}, 24(1):214--236, 1962.
\href{https://doi.org/10.1007/BF02733276}{DOI:10.1007/BF02733276}

%---93
\bibitem{Rudin1991}
W.~Rudin.
\newblock {\em Functional Analysis}.
\newblock McGraw-Hill Science, 2nd edition, 1991.
\href{https://archive.org/details/functionalanalys0000rudi}{archive.org/details/functionalanalys0000rudi}

%---94
\bibitem{Dixmier1977}
J.~Dixmier.
\newblock {\em C*-Algebras}.
\newblock North-Holland Mathematical Library, Vol. 15, 1977.
\href{https://doi.org/10.1016/S0924-6509(08)X7031-X}{DOI:10.1016/S0924-6509(08)X7031-X}

%---95
\bibitem{Titchmarsh1986}
E.~C. Titchmarsh.
\newblock {\em The Theory of the Riemann Zeta-function}.
\newblock Oxford University Press, 2nd edition (revised by D. R. Heath-Brown), 1986.
\href{https://archive.org/details/theoryofriemannz00titc}{archive.org/details/theoryofriemannz00titc}

%---96
\bibitem{Kato1995}
T.~Kato.
\newblock {\em Perturbation Theory for Linear Operators}.
\newblock Classics in Mathematics. Springer-Verlag, Berlin, 1995.
\href{https://doi.org/10.1007/978-3-642-66282-9}{DOI:10.1007/978-3-642-66282-9}

%---97
\bibitem{Odlyzko1987}
A.~M. Odlyzko.
\newblock On the distribution of spacings between zeros of the zeta function.
\newblock {\em Mathematics of Computation}, 48(177):273--308, 1987.
\href{https://doi.org/10.2307/2007890}{DOI:10.2307/2007890}

%---98
\bibitem{ReedSimonI}
M.~Reed and B.~Simon.
\newblock {\em Methods of Modern Mathematical Physics, Vol. 1: Functional Analysis}.
\newblock Academic Press, 1980.
\href{https://doi.org/10.1016/C2013-0-07421-2}{DOI:10.1016/C2013-0-07421-2}

%---99
\bibitem{Davies1995}
E.~B. Davies.
\newblock {\em Spectral Theory and Differential Operators}.
\newblock Cambridge Studies in Advanced Mathematics. Cambridge University Press, 1995.
\href{https://doi.org/10.1017/CBO9780511623721}{DOI:10.1017/CBO9780511623721}

%---100
\bibitem{Minakshisundaram1949}
S.~Minakshisundaram and {\AA}.~Pleijel.
\newblock Some properties of the eigenfunctions of the Laplace-operator on Riemannian manifolds.
\newblock {\em Canadian Journal of Mathematics}, 1:242--256, 1949.
\href{https://doi.org/10.4153/CJM-1949-021-5}{DOI:10.4153/CJM-1949-021-5}

%---101
\bibitem{GelfandShilov1968}
I.~M. Gelfand and G.~E. Shilov.
\newblock {\em Generalized Functions, Vol. 2: Spaces of Fundamental and Generalized Functions}.
\newblock Academic Press, 1968.
\href{https://archive.org/details/generalizedfunct0000gelf}{ISBN: 978-0122795023}

%---102
\bibitem{Berline2004}
N.~Berline, E.~Getzler, and M.~Vergne.
\newblock {\em Heat Kernels and Dirac Operators}.
\newblock Grundlehren der mathematischen Wissenschaften. Springer, 2004.
\href{https://doi.org/10.1007/978-3-642-56475-8}{DOI:10.1007/978-3-642-56475-8}

%---103
\bibitem{Gilkey1995}
P.~B. Gilkey.
\newblock {\em Invariance Theory, the Heat Equation, and the Atiyah-Singer Index Theorem}.
\newblock Studies in Advanced Mathematics. CRC Press, 2nd edition, 1995.
\href{https://doi.org/10.1201/9780367805906}{DOI:10.1201/9780367805906}

%---104
\bibitem{Halmos1974}
P.~R. Halmos.
\newblock {\em Measure Theory}.
\newblock Graduate Texts in Mathematics. Springer-Verlag, 1974.
\href{https://archive.org/details/measuretheory00halm}{ISBN: 978-0387900889}

%---105
\bibitem{SteinWeiss1971}
E.~M. Stein and G.~Weiss.
\newblock {\em Introduction to Fourier Analysis on Euclidean Spaces}.
\newblock Princeton University Press, 1971.
\href{https://doi.org/10.1515/9781400883899}{DOI:10.1515/9781400883899}

%---106
\bibitem{Katznelson2004}
Y.~Katznelson.
\newblock {\em An Introduction to Harmonic Analysis}.
\newblock Cambridge University Press, 3rd edition, 2004.
\href{https://doi.org/10.1017/CBO9781139165372}{DOI:10.1017/CBO9781139165372}

%---107
\bibitem{Odlyzko1988}
A.~M. Odlyzko.
\newblock The $10^{20}$-th zero of the Riemann zeta function and 70 million of its neighbors.
\newblock {\em AT\&T Bell Laboratories}, 1988.
\href{http://www.dtc.umn.edu/~odlyzko/unpublished/zeta.10to20.1992.pdf}{Unpublished Manuscript}

%---108
\bibitem{Gourdon2004}
X.~Gourdon.
\newblock The $10^{13}$ first zeros of the Riemann Zeta function, and zeros computation at very large height.
\newblock {\em Unpublished}, 2004.
\href{http://numbers.computation.free.fr/Constants/Miscellaneous/zetazeros1e13-1e24.pdf}{Numbers.computation.free.fr}

%---109
\bibitem{Vassilevich2003}
D.~V. Vassilevich.
\newblock Heat kernel expansion: user's manual.
\newblock {\em Physics Reports}, 388(5-6):279--360, 2003.
\href{https://doi.org/10.1016/j.physrep.2003.09.002}{DOI:10.1016/j.physrep.2003.09.002}

%---110
\bibitem{Hawking1977}
S.~W. Hawking.
\newblock Zeta function regularization of path integrals in curved spacetime.
\newblock {\em Communications in Mathematical Physics}, 55(2):133--148, 1977.
\href{https://doi.org/10.1007/BF01626516}{DOI:10.1007/BF01626516}

%---111
\bibitem{Weil1952}
A.~Weil.
\newblock Sur les ``formules explicites'' de la th{\'e}orie des nombres premiers.
\newblock {\em Meddelanden fr{\aa}n Lunds Universitets Matematiska Seminarium}, 1952:252--265, 1952.
\href{http://www.numdam.org/item?id=BSMF_1952__80__252_0}{Numdam Link}

%---112
\bibitem{Atiyah1975}
M.~F. Atiyah, V.~K. Patodi, and I.~M. Singer.
\newblock Spectral asymmetry and Riemannian geometry. I.
\newblock {\em Mathematical Proceedings of the Cambridge Philosophical Society}, 77(1):43--69, 1975.
\href{https://doi.org/10.1017/S030500410005114X}{DOI:10.1017/S030500410005114X}

%---113
\bibitem{Osterwalder1973}
K.~Osterwalder and R.~Schrader.
\newblock Axioms for Euclidean Green's functions.
\newblock {\em Communications in Mathematical Physics}, 31(2):83--112, 1973.
\href{https://doi.org/10.1007/BF01645738}{DOI:10.1007/BF01645738}

%---114
\bibitem{Hecke1936}
E.~Hecke.
\newblock {\"U}ber die Bestimmung Dirichletscher Reihen durch ihre Funktionalgleichung.
\newblock {\em Mathematische Annalen}, 112:664--699, 1936.
\href{https://doi.org/10.1007/BF01565437}{DOI:10.1007/BF01565437}

%---115
\bibitem{IwaniecKowalski2004}
H.~Iwaniec and E.~Kowalski.
\newblock {\em Analytic Number Theory}.
\newblock American Mathematical Society, Colloquium Publications, Vol. 53, 2004.
\href{https://doi.org/10.1090/coll/053}{DOI:10.1090/coll/053}

%---116
\bibitem{Voros1987}
A.~Voros.
\newblock Spectral functions, special functions and the Selberg zeta function.
\newblock {\em Communications in Mathematical Physics}, 110(3):439--465, 1987.
\href{https://doi.org/10.1007/BF01212422}{DOI:10.1007/BF01212422}

%---117
\bibitem{Kohmoto1983}
M.~Kohmoto, L.~P. Kadanoff, and C.~Tang.
\newblock Localization Problem in One Dimension: Mapping and Escape.
\newblock {\em Physical Review Letters}, 50:1870--1872, 1983.
\href{https://doi.org/10.1103/PhysRevLett.50.1870}{DOI:10.1103/PhysRevLett.50.1870}

%---118
\bibitem{Ostlund1983}
S.~Ostlund, R.~Pandit, D.~Rand, H.~J. Schellnhuber, and E.~D. Siggia.
\newblock One-dimensional Schr{\"o}dinger equation with an almost periodic potential.
\newblock {\em Physical Review Letters}, 50:1873--1876, 1983.
\href{https://doi.org/10.1103/PhysRevLett.50.1873}{DOI:10.1103/PhysRevLett.50.1873}

%---119
\bibitem{Teschl2000}
G.~Teschl.
\newblock {\em Jacobi Operators and Completely Integrable Nonlinear Lattices}.
\newblock American Mathematical Society, Mathematical Surveys and Monographs, Vol. 72, 2000.
\href{https://doi.org/10.1090/surv/072}{DOI:10.1090/surv/072}

%---120
\bibitem{Trefethen2000}
L.~N. Trefethen.
\newblock {\em Spectral Methods in MATLAB}.
\newblock SIAM, 2000.
\href{https://doi.org/10.1137/1.9780898719598}{DOI:10.1137/1.9780898719598}

%---121
\bibitem{DavisRabinowitz1984}
P.~J. Davis and P.~Rabinowitz.
\newblock {\em Methods of Numerical Integration}.
\newblock Academic Press, 2nd edition, 1984.
\href{https://doi.org/10.1016/C2013-0-10335-9}{DOI:10.1016/C2013-0-10335-9}

%---122
\bibitem{StoicaMoses2005}
P.~Stoica and R.~L. Moses.
\newblock {\em Spectral Analysis of Signals}.
\newblock Pearson Prentice Hall, 2005.
\href{http://www.mps.aau.dk/digitalAssets/128/128965_spectral-analysis-of-signals.pdf}{ISBN: 978-0131139565}

%---123
\bibitem{Press2007}
W.~H. Press, S.~A. Teukolsky, W.~T. Vetterling, and B.~P. Flannery.
\newblock {\em Numerical Recipes: The Art of Scientific Computing}.
\newblock Cambridge University Press, 3rd edition, 2007.
\href{http://numerical.recipes/}{ISBN: 978-0521880688}

%---124
\bibitem{Harris1978}
F.~J. Harris.
\newblock On the use of windows for harmonic analysis with the discrete Fourier transform.
\newblock {\em Proceedings of the IEEE}, 66(1):51--83, 1978.
\href{https://doi.org/10.1109/PROC.1978.10837}{DOI:10.1109/PROC.1978.10837}

%---125
\bibitem{PolyaHilbert}
H.~L. Montgomery.
\newblock The pair correlation of zeros of the zeta function.
\newblock In \emph{Analytic Number Theory}, Proc. Sympos. Pure Math., Vol. XXIV, pages 181--193. Amer. Math. Soc., Providence, R.I., 1973.
\newblock \href{https://doi.org/10.1090/pspum/024/0337821}{DOI: 10.1090/pspum/024/0337821}.

%---126
\bibitem{HaaKus2024}
F.~Haake, S.~Gnutzmann, and M.~Ku\'s.
\newblock \emph{Quantum Signatures of Chaos}.
\newblock Springer Series in Synergetics. Springer Cham, 4th edition, 2018.
\newblock \href{https://doi.org/10.1007/978-3-319-97580-1}{DOI: 10.1007/978-3-319-97580-1}.


\bibitem{Siegel1932}
C.~L. Siegel.
\newblock \"Uber Riemanns Nachla\ss\ zur analytischen Zahlentheorie.
\newblock \emph{Quellen und Studien zur Geschichte der Mathematik, Astronomie und Physik}, Abt. B: Studien 2, 45--80, 1932.

\end{thebibliography}


\end{document}